\chapter{模型压缩与加速}

\paragraph*{本章阅读锚点}
原始 3DGS 的速度优势来自显式高斯和 rasterization，但代价是场景里常有数百万个基元，每个基元还带位置、协方差、不透明度和球谐颜色。压缩与加速方法首先要判断“冗余在哪里”：有些冗余在高斯数量上，可以剪枝、合并或重新控制 densification；有些冗余在属性上，可以量化、码本化、熵编码或用神经场预测；有些冗余在视角和距离上，可以用 LoD、锚点或集群只渲染当前需要的细节；还有一类冗余在训练过程中，早期高分辨率渲染和后期无效更新都可以省掉。本章按技术路线读：先看率失真和编码压缩，再看剪枝稀疏化，随后看 Scaffold/LoD 这类结构化表示，最后讨论训练调度、视频流式传输以及少数跨表示压缩方法。

\section{编码与压缩}

编码压缩路线的锚点不是简单“把文件 zip 一下”，而是把 3DGS 改写成适合码率控制的表示。原始 3DGS 的高斯点云无序、属性维度高，位置和外观分布又高度非均匀；因此后续方法分别从率失真优化、锚点上下文、码本量化、神经解码器、点云编码器和视频编码器切入。可以把它们看成三层：第一层减少要存的基元或属性，第二层让属性分布更可预测，第三层把预测残差交给熵编码或标准 codec。



% METHOD: 2311.13681
\paragraph{Compact 3D Gaussian Representation（早期码本压缩）（2023.11）}
Compact 3D Gaussian Representation~\cite{2311_13681} 是 3DGS 压缩的早期基础方法之一，也是后续许多码本、神经属性场和剪枝方法的参照点。它直接面对原始 3DGS 的三类冗余：densification 会生成大量贡献很小的高斯；每个高斯都保存高阶球谐颜色系数，即使该区域几乎没有视角相关外观；尺度和旋转等几何属性在局部区域内高度相似，却仍以独立浮点向量保存。

方法包含三块。第一块是 volume-based learnable mask：为高斯学习一个 mask 参数，用 straight-through estimator 把二值 mask 接入训练；mask 同时作用在 opacity 和 scale 上，使体积很小、透明度低、对渲染贡献弱的高斯在训练中被逐步删除，而不是只靠原始 3DGS 的 opacity pruning。第二块是 grid-based neural field for color：不再为每个高斯单独保存全部 SH 系数，而是在高斯位置处查询 hash/grid 特征，再由小 MLP 解码视角相关颜色；这样颜色从“每个高斯一大组系数”变成“共享网格特征 + 按位置查询”。第三块是 residual vector quantization for geometry：尺度和旋转被映射到共享码本，最终只需要保存码本索引和少量必要参数。训练结束后，方法还可对 opacity、hash-grid 参数和 R-VQ 索引做 8-bit min-max quantization、参数剪枝和 Huffman coding，进一步减小落盘体积。相对后续 ContextGS/HAC++ 的复杂上下文模型，它更像“剪枝、神经属性场、VQ”三件套的起点，为后续压缩路线提供了模块化拆解方式；它的一个工程限制是原始 SIBR viewer 不能直接渲染 neural-field color，需要绕过或改造颜色分支。
\textbf{方法信息：} 代码：\url{https://github.com/maincold2/Compact-3DGS}；GitHub 星标很高，约 500 stars；引用很多，Semantic Scholar 约 461 次。
\textbf{设计动机与边界：} 相对原始 3DGS 以独立高斯和图像重建误差为主的设定，Compact 3D Gaussian Representation for Radiance Field 把问题落在以下可实现的改动上：C3DGS的压缩框架包含三个核心策略。第一，高斯体积掩码：不同于原始3D-GS仅基于不透明度剪枝，C3DGS引入基于体积的可学习掩码。掩码同时作用于不透明度和缩放属性，通过可微分训练自动识别体积小、对渲染贡献微弱的高斯并掩蔽掉。训练后使用硬掩码永久移除这些高斯。实验显示仅42%的高斯即可维持相当性能。第二，紧凑颜色表示：用基于网格的神经场（类似TensoRF的分解网格结构）替代每个高斯的球谐系数。颜色通过查询高斯位置的网格特征并经过小型MLP解码得到，将每个高斯的颜色参数从48维（4阶SH）减少到3维位置查询。第三，几何属性向量量化：对高斯的缩放和旋转属性使用码本量化，通过学习共享码本将连续值映射到有限集合，结合熵编码进一步压缩存储。其设计目的不是单独增加一个模块，而是用上述表示、监督或渲染约束消除原流程在目标任务中的具体失配。提出Compact 3DGS框架，从三个维度压缩3D-GS：1）可学习的体积掩码策略识别并移除非必要高斯，显著减少高斯数量；2）用基于网格的神经场替代球谐系数表示视角相关颜色，大幅降低每个高斯的特征维度；3）通过码本学习的向量量化紧凑表示高斯的几何属性（缩放、旋转）。实现25倍以上存储压缩。适用时还应保留论文给出的边界：极端压缩下质量损失明显；神经颜色场增加推理时的MLP计算；码本量化对属性分布有假设。
% METHOD: 2404.09458
\paragraph{CompGS（锚点基元与耦合基元）（2024.04）}
CompGS~\cite{2404_09458} 是压缩表示中的重要早期方法。公开代码包含训练、测试、模块和配置目录，README 明确围绕压缩 3DGS 的率失真曲线组织实验。它的问题起点是：原始高斯全部保存完整属性会造成大量重复，而简单剪枝又容易破坏渲染质量。

CompGS 把高斯分为信息完整的 anchor primitives 和更轻量的 coupled primitives。anchor 保存位置、协方差和参考嵌入 $f$；每个 anchor 关联若干 coupled primitives，后者只保存紧凑残差嵌入 $g_k$。解码时，把 $f$ 与 $g_k$ 拼接后送入小 MLP：几何分支预测平移、缩放和旋转，用仿射变换从 anchor 推出 coupled primitive 的位置与协方差；外观分支再结合视角 embedding 预测颜色和不透明度。训练目标同时约束重建误差和码率，码率项通过熵估计建模 anchor 与 coupled primitive 的编码成本，并用 Lagrange 权重控制 $R$ 和 $D$ 的取舍。相对 Compact 3D Gaussian 的属性码本化，CompGS 更进一步利用基元间相关性：不是每个高斯都独立压缩，而是让一部分高斯成为局部预测上下文。
\textbf{实现细节：} CompGS方法包含三个核心设计：（1）混合原语结构：将3D高斯分为两类。锚点原语（anchor primitives）数量少，包含完整属性（位置、协方差、参考嵌入 $f$）。每个锚点关联K个耦合原语（coupled primitives），耦合原语仅包含紧凑残差嵌入 $g_k$。（2）原语间预测：通过通道拼接融合锚点的 $f$和耦合原语的 $g_k$，通过MLP预测耦合原语的几何属性（位置、协方差，通过仿射变换）和外观属性（不透明度、颜色）。由于多数原语以残差形式存储，大幅减少了数据量。（3）码率约束优化：引入率失真优化框架，L = min(R + D)，其中R通过熵估计建模锚点和耦合原语的码率，D为渲染失真。联合优化所有原语实现紧凑表示。相比独立压缩每个高斯的方法，CompGS利用了高斯间的局部相似性（空间相邻的高斯具有相似的属性），通过预测消除冗余。 提出CompGS，一种面向3DGS的率失真优化压缩方法。核心创新是混合原语结构：少量锚点原语携带完整参考信息，多数耦合原语仅包含紧凑残差嵌入，通过锚点预测耦合原语的属性。这种结构利用了3D高斯间的局部相似性，消除了独立压缩中的冗余。进一步引入码率约束优化，联合优化渲染失真和码率消耗，实现最优的率失真权衡。在保持渲染质量的同时显著压缩3DGS数据。 混合结构增加了编解码复杂度；锚点选择策略对压缩效果影响较大。
\textbf{设计动机与边界：} 相对原始 3DGS 以独立高斯和图像重建误差为主的设定，CompGS 把问题落在以下可实现的改动上：CompGS方法包含三个核心设计：（1）混合原语结构：将3D高斯分为两类。锚点原语（anchor primitives）数量少，包含完整属性（位置、协方差、参考嵌入f\_）。每个锚点关联K个耦合原语（coupled primitives），耦合原语仅包含紧凑残差嵌入g\_k。（2）原语间预测：通过通道拼接融合锚点的f\_和耦合原语的g\_k，通过MLP预测耦合原语的几何属性（位置、协方差，通过仿射变换）和外观属性（不透明度、颜色）。由于多数原语以残差形式存储，大幅减少了数据量。（3）码率约束优化：引入率失真优化框架，L = min(R + D)，其中R通过熵估计建模锚点和耦合原语的码率，D为渲染失真。联合优化所有原语实现紧凑表示。相比独立压缩每个高斯的方法，CompGS利用了高斯间的局部相似性（空间相邻的高斯具有相似的属性），通过预测消除冗余。其设计目的不是单独增加一个模块，而是用上述表示、监督或渲染约束消除原流程在目标任务中的具体失配。提出CompGS，一种面向3DGS的率失真优化压缩方法。核心创新是混合原语结构：少量锚点原语携带完整参考信息，多数耦合原语仅包含紧凑残差嵌入，通过锚点预测耦合原语的属性。这种结构利用了3D高斯间的局部相似性，消除了独立压缩中的冗余。进一步引入码率约束优化，联合优化渲染失真和码率消耗，实现最优的率失真权衡。在保持渲染质量的同时显著压缩3DGS数据。适用时还应保留论文给出的边界：混合结构增加了编解码复杂度；锚点选择策略对压缩效果影响较大。
\textbf{方法信息：} 代码：\url{https://github.com/liuxiangrui/compgs}；引用很多，Semantic Scholar 约 113 次。



% METHOD: 2405.17083
\paragraph{F-3DGS（坐标与属性分解）（2024.05）}
F-3DGS~\cite{2405_17083} 的锚点是低秩分解。原始 3DGS 把每个高斯的位置、尺度、旋转和颜色独立存储，但密集区域里的高斯常有相似空间结构和相似属性；F-3DGS 的改动是把这些局部高斯簇看作可分解的张量或矩阵，而不是一串彼此独立的参数。

它先把密集区域的位置改成 factorized coordinates：不再显式保存所有 3D 点，而是在坐标轴或坐标平面上保存少量 1D/2D 因子，通过张量积展开出一组 3D 高斯位置；多组不同分辨率的因子可以覆盖不同尺度的结构。随后，颜色、尺度、旋转、不透明度等属性也用 factorized representation 表示，再由小 MLP 或线性组合解码为标准高斯属性。论文和代码都保留原始 3DGS rasterizer：实践中先训练一个原始 3DGS 并把其 \texttt{.ply} 作为 \texttt{--hist\_path}，用它的高斯分布初始化或约束分解表示；训练前段联合优化坐标和属性，后段固定坐标、继续优化 factorized attributes，以避免联合优化不稳定。渲染时把分解表示展开为普通高斯后再 splat，因此它压缩的是存储参数化，而不是替换渲染器。相对码本量化，F-3DGS 更像利用连续低秩结构来减少冗余；它适合几何和外观有局部重复模式的场景。
\textbf{方法信息：} 代码：\url{https://github.com/Xiangyu1Sun/Factorize-3DGS}。
\textbf{设计动机与边界：} 相对原始 3DGS 以独立高斯和图像重建误差为主的设定，F-3DGS 把问题落在以下可实现的改动上：F-3DGS观察到3DGS中的高斯存在大量结构模式和冗余，密集区域的高斯聚类可以通过分解技术用更少的参数近似。核心方法：1）坐标分解：将密集高斯聚类的位置表示为沿各坐标轴的分解。例如，一个3D区域内的N个高斯位置可通过低秩矩阵分解近似为三个因子矩阵的乘积（类似CP分解），将存储从O(N*3)降至O(R*(3+N))，R为秩。2）属性分解：对高斯的颜色（球谐系数）、缩放、旋转等属性同样应用分解策略。属性通过共享基向量和每高斯系数的线性组合表示，系数进一步通过分解压缩。3）近似策略：通过SVD或低秩近似选择最优的分解秩，平衡压缩率和渲染质量。训练时在原始参数化空间中优化，定期进行分解更新。4）渲染流程保持不变：分解后的高斯仍通过标准3DGS光栅化管线渲染，无需修改渲染代码。其设计目的不是单独增加一个模块，而是用上述表示、监督或渲染约束消除原流程在目标任务中的具体失配。提出F-3DGS，受矩阵和张量分解启发，通过高效分解表示密集高斯聚类，用显著更少的高斯近似原始高斯集合。将大量高斯及其颜色、缩放、旋转等属性通过沿各轴及其组合的有限信息编码，大幅减少存储开销同时保持渲染质量。适用时还应保留论文给出的边界：压缩率与质量之间存在权衡；分解秩选择需要调节；对极其稀疏分布的高斯压缩效果有限。
% METHOD: 2405.20721
\paragraph{ContextGS（锚点级自回归上下文）（2024.05）}
ContextGS~\cite{2405_20721} 继承 Scaffold-GS 的锚点思想，但把锚点之间的空间关系显式用于熵编码。公开 README 说明其核心是 anchor-level autoregressive model：粗层锚点预测细层锚点，从而减少 anchor 属性中的空间冗余。

具体做法是先对 Scaffold-GS anchor 做自底向上的层级划分：从细层 anchor 出发，用逐渐变大的体素尺寸合并或抽取更粗层 anchor，并通过目标层间比例和二分搜索控制每一层的规模。编码时按粗到细顺序处理。最粗层没有可用上下文，因此用低维量化特征作为 hyperprior；从第二层开始，已经解码的粗层 anchor 特征会输入上下文网络，预测当前层 anchor 属性的均值、尺度或概率分布，真实属性只需要以残差或符号形式进入熵编码。解码端按同样顺序恢复 anchor，再由 Scaffold-GS 的 anchor-to-Gaussian 解码器生成可渲染高斯；论文还使用 anchor forward，让粗层 anchor 可以直接参与后续层生成，避免上下文只服务编码而不服务渲染。相对 Scaffold-GS 只用 anchor 生成视角自适应高斯，ContextGS 把 anchor 进一步变成压缩上下文；相对普通 VQ，它不是只学全局码本，而是利用局部和层级结构预测当前属性。
\textbf{实现细节：} ContextGS（2024.05） 的核心思想是将信息流从"单向超先验"升级为双向自回归：已编码锚点的具体属性值被直接用作未编码锚点的上下文。直觉上，如果 HAC++ 是"先看地图，再猜房子"，那么 ContextGS 是"先看邻居的房子，再猜自己的房子"——后者显然提供了更精确的预测信息，因为邻居的实际特征值比一张抽象的空间地图包含了更丰富的场景细节。
\textbf{设计动机与边界：} 相对原始 3DGS 以独立高斯和图像重建误差为主的设定，ContextGS 把问题落在以下可实现的改动上：方法基于Scaffold-GS的锚点结构，提出创新的锚点划分策略和自回归编码框架。锚点划分采用自底向上方法，根据不同体素大小将锚点分为K个不相交层级，每层是更细层的相对均匀下采样。通过体素量化建立层级间可追溯的映射关系。自回归编码按从粗到细的顺序处理，粗层已解码的锚点特征可作为细层的先验条件，提升编码效率。对于最粗层无已编码锚点的情况，引入低维量化特征作为超先验。提出自适应参数搜索策略，根据场景特性自动确定体素缩放参数，避免逐场景调参。整体框架包括锚点划分、自回归熵编码和率失真优化训练。为了消除原始 3DGS 在该目标任务中的具体失配，方法把上述表示、监督或渲染约束作为可学习和可执行的改动，而不是单独堆叠模块。提出首个锚点级别的自回归上下文模型用于3DGS压缩。将锚点划分为不同层级，未编码锚点可基于已编码的粗层级锚点预测，实现更准确的建模和更高的编码效率。引入低维量化特征作为超先验。相比原始3DGS实现超过100倍压缩，相比Scaffold-GS实现15倍压缩，同时保持相当或更高的渲染质量。适用时还应保留论文给出的边界：自回归编码增加了解码时间，不完全适合实时渲染的快速加载场景；锚点层级数量需预设；对非均匀分布场景的划分策略可能不够优化。
\textbf{方法信息：} 代码：\url{https://github.com/wyf0912/contextgs}；GitHub 约 92 stars；引用很多，Semantic Scholar 约 113 次。



% METHOD: 2406.01597
\paragraph{RDO-Gaussian（端到端率失真优化）（2024.06）}
End-to-End Rate-Distortion Optimized 3D Gaussian Representation~\cite{2406_01597} 是本章的代表方法之一。它继承原始 3DGS 的显式高斯和可微 rasterization，但把“质量好不好”和“码率高不高”放进同一个优化目标，而不是先训练完整模型再事后压缩。原始 3DGS 的问题是存储粒度过细：一个高斯通常要保存位置、协方差、不透明度和大量球谐颜色系数，其中颜色球谐占了很大比例；但训练目标只关心渲染误差，不会主动判断某个高斯、某一阶球谐或某组属性到底值不值得为它付码率。

它的改变是用 $R+\lambda D$ 的方式联合控制失真和码率。训练流程通常先按原始 3DGS 稳定到中期，再接入三个可学习压缩模块：第一，动态高斯剪枝为每个高斯学习一个 mask，前向渲染时使用硬选择，反向传播时用 straight-through estimator 让保留/删除决策仍能被优化；第二，自适应球谐剪枝不再让所有区域使用同样数量的颜色系数，而是按区域和材质复杂度学习需要保留的球谐阶数；第三，熵约束向量量化把剩余属性映射到离散码本，并用码字概率估计 rate。这样，$\lambda$ 控制的是完整率失真曲线：增大码率约束时，模型会同时减少高斯数量、降低颜色阶数并改变属性量化精度。推理或存储时，只保留 mask 后的高斯、球谐子集和量化索引，再交给熵编码。相对 LightGaussian 这类“训练后剪枝加量化”的路线，RDO-Gaussian 的锚点意义在于把压缩变成训练目标本身；论文报告在多个真实和合成场景中可超过 $40\times$ 压缩，在更低码率点还能进一步提高压缩倍数。
\textbf{设计动机与边界：} 相对原始 3DGS 以独立高斯和图像重建误差为主的设定，End-to-End Rate-Distortion Optimized 3D Gaussian Representation 把问题落在以下可实现的改动上：RDO-Gaussian将3D高斯表示学习建模为R+lambda*D的联合优化问题。三个核心组件：（1）Gaussian Pruning——通过可学习的掩码动态剪枝冗余高斯，优化过程中自适应控制高斯数量；（2）Adaptive SH Pruning——观察到一个高斯59个参数中45个（约76%）用于SH系数但贡献有限，因此为每个高斯学习独立的SH阶数掩码，使不同区域和材质的高斯拥有不同数量的参数（漫反射物体用低阶SH，视角相关外观用高阶SH）；（3）Entropy-Constrained Vector Quantization (ECVQ)——对协方差和颜色参数进行熵约束向量量化，在优化过程中同时考虑量化误差和码率成本。整个流程端到端可微分，通过调节lambda实现灵活连续的码率-失真权衡。其设计目的不是单独增加一个模块，而是用上述表示、监督或渲染约束消除原流程在目标任务中的具体失配。首次将3D高斯压缩建模为端到端率失真优化问题，实现灵活连续的码率控制。三个核心创新：（1）动态剪枝和ECVQ同时优化率和失真；（2）自适应SH剪枝使不同高斯可拥有不同阶数的SH系数；（3）端到端学习避免了多阶段优化的局限。在真实和合成场景上将3D高斯大小减少40倍以上，率失真性能超越现有方法。适用时还应保留论文给出的边界：端到端优化的训练时间可能较长，对极高压缩率下的细节保持仍有改进空间。
\textbf{方法信息：} Microsoft 参与了这篇论文；代码：\url{https://github.com/USTC-IMCL/RDO-Gaussian}；GitHub 约 29 stars；引用很多，Semantic Scholar 约 85 次。



% METHOD: 2406.17074
\paragraph{Reducing the Memory Footprint（分辨率感知属性压缩）（2024.06）}
Reducing the Memory Footprint of 3D Gaussian Splatting~\cite{2406_17074} 从部署显存出发，而不仅是磁盘文件大小。它观察到远距离或低分辨率渲染中，高阶球谐和过小高斯贡献有限；因此压缩应与目标分辨率和可见贡献相关。

具体做法围绕三个存储来源逐一动刀。第一，resolution-aware primitive pruning 不只看 opacity，而是评估高斯在目标显示分辨率下的实际可见贡献和冗余程度；许多投影面积很小或被其他高斯遮挡的 primitive，在移动端或低分辨率显示中几乎不可见，因此可以被删除。第二，自适应球谐阶数为每个 primitive 判断需要多少 SH band：如果多视角一致性显示该高斯几乎没有视角相关外观，就退化为低阶或单色表示；只有镜面、高光或明显方向性区域保留更多系数。第三，剩余属性通过 codebook quantization 和 half-float 存储降低位宽。和 RDO-Gaussian 的端到端率失真训练不同，这篇更像对原始 3DGS 的工程化压缩清单：数量、颜色阶数和数值精度分别压缩，组合后论文报告约 $27\times$ 磁盘缩减和约 $1.7\times$ rasterization 加速；移动端示例中，下载时间从分钟级下降到数秒级。
\textbf{实现细节：} 该方法系统性地分析了3DGS高内存占用的原因并提出三个维度的解决方案：1）基元数量：扩展了现有剪枝方案，开发了分辨率感知的冗余分数，使用高效的两步算法消除约60
\textbf{方法信息：} 代码：\url{https://github.com/graphdeco-inria/reduced-3dgs}；GitHub 星标很高，约 237 stars；引用很多，Semantic Scholar 约 195 次。
\textbf{设计动机与边界：} 相对原始 3DGS 以独立高斯和图像重建误差为主的设定，Reducing the Memory Footprint of 3D Gaussian Splatting 把问题落在以下可实现的改动上：该方法系统性地分析了3DGS高内存占用的原因并提出三个维度的解决方案：1）基元数量：扩展了现有剪枝方案，开发了分辨率感知的冗余分数，使用高效的两步算法消除约60%的基元（相比基线）。该方法考虑了不同渲染分辨率下基元的实际贡献，在保持视觉质量的同时大幅减少基元数量。2）SH系数：观察到原始3DGS为所有基元使用三个SH带，即使在无视角依赖效果的情况下也如此。引入自适应调整方法，根据每个高斯基元的实际需要选择SH系数数量，许多基元仅需单一颜色值即可。3）量化精度：大多数基元属性不需要高精度或大动态范围。提出基于码本的量化方法，将属性值映射到共享码本中的索引，配合半浮点表示进一步减少内存。三个组件协同工作，在不损失视觉质量的前提下大幅降低存储需求。其设计目的不是单独增加一个模块，而是用上述表示、监督或渲染约束消除原流程在目标任务中的具体失配。从三个方面系统性地解决3DGS存储问题：提出分辨率感知的基元剪枝方法减少约60%基元；引入自适应SH带数选择方法；提出基于码本的量化方法和半浮点表示。三者结合实现27倍磁盘存储缩减和1.7倍渲染加速。在移动设备上展示了显著减少的下载时间。适用时还应保留论文给出的边界：码本量化的额外解码开销在某些硬件上可能影响渲染速度。极端压缩比下可能出现轻微的视觉质量损失。
% METHOD: 2406.18214
\paragraph{Trimming the Fat（梯度感知迭代剪枝）（2024.06）}
Trimming the Fat~\cite{2406_18214} 走的是简单可解释的后处理剪枝路线。它不重写表示，也不设计复杂码本，而是根据训练或微调过程中的梯度信号估计高斯重要性。

流程从一个已经训练好的 3DGS checkpoint 开始。方法先继续渲染/反传若干步，累计每个高斯相关参数的梯度幅值；梯度长期较小的高斯被视作对损失影响弱，进入候选删除集合。随后按目标 pruning level 删除一批低重要性高斯，再用短暂恢复训练让剩余高斯重新分配职责。这个过程可以多轮迭代，因此不会像一次性阈值剪枝那样突然制造空洞。论文还做了一个“lottery ticket”式实验：如果把剪枝后保留下来的高斯重新随机初始化再训练，模型无法正常收敛，说明被保留的不是一个可从零训练的万能子网络，而是已经编码了场景结构的 learned prior。相对 PUP 的二阶不确定性估计，它更轻量；相对 LightGaussian，它更集中在 pruning 本身而不是完整压缩流水线。
\textbf{方法信息：} 代码：\url{https://github.com/salmanali96/trimming-the-fat}；GitHub 约 15 stars；引用很多，Semantic Scholar 约 51 次。
\textbf{设计动机与边界：} 相对原始 3DGS 以独立高斯和图像重建误差为主的设定，Trimming the Fat 把问题落在以下可实现的改动上：方法核心是梯度感知的迭代剪枝（Gradient Aware Pruning）。在标准3DGS训练完成后，利用训练过程中累积的梯度信息来评估每个高斯对最终渲染结果的贡献重要性。具体而言，方法监控每个高斯在训练过程中的梯度幅值，梯度较小的高斯被认为对损失函数影响较小即冗余的。通过设定梯度阈值Q，将梯度低于阈值的高斯标记为可剪枝。剪枝过程是迭代进行的，每次迭代移除一部分最不重要的高斯，然后对剩余高斯进行微调以恢复质量。这种渐进式剪枝策略避免了一次性大幅剪枝导致的质量下降。方法不修改原始3DGS训练流程，作为后处理步骤应用，具有良好的通用性和兼容性。其设计目的不是单独增加一个模块，而是用上述表示、监督或渲染约束消除原流程在目标任务中的具体失配。提出梯度感知的迭代剪枝方法，可在移除高达75%高斯的情况下保持或提升基线性能，实现约50倍压缩且计算速度可达600FPS。为3DGS模型的可扩展性提供了初步解决方案。适用时还应保留论文给出的边界：剪枝后仍需要微调阶段来恢复质量。对梯度阈值选择较为敏感。极端压缩率下细节可能丢失。仅关注存储压缩，未涉及训练效率提升。
% METHOD: 2407.09510
\paragraph{3DGS.zip（压缩综述与路线图）（2024.07）}
3DGS.zip~\cite{2407_09510} 本身是综述，但在可信书中应保留，因为它给压缩路线建立了清晰分类。它把 3DGS 压缩拆成压缩方法和精简方法两大类：前者包括向量量化、熵编码、特征平面、上下文模型；后者包括剪枝、LoD、致密化控制和更紧凑的高斯生成。

对读者来说，3DGS.zip 的作用不是提供一个新 codec，而是帮助理解本章方法之间的关系：相对原始 3DGS 只给出显式高斯表示和 densification，它把后续工作按“属性怎么编码、基元怎么减少、层级怎么组织”三条锚点路线拆开。ContextGS/HAC++ 是上下文熵编码，MesonGS/Compact 3D Gaussian 是后训练变换和量化，Taming3DGS/Mini-Splatting 类方法控制高斯增长，Octree/LoD 方法则改变可见细节组织方式。
因此本章把 3DGS.zip 当作路线图使用：当一个方法主要减少每个高斯或 anchor 的属性比特时，应放在“编码与压缩”；当一个方法主要减少需要渲染的高斯数量时，应放在“剪枝与稀疏化”；当一个方法通过 octree、cluster 或连续距离函数改变可见细节组织时，应放在“层级与 LoD”。这一区分能避免把所有“小模型”工作混在一起，也能看清压缩、剪枝和层级表示之间的互补关系。
\textbf{方法信息：} 引用很多，Semantic Scholar 约 87 次。
\textbf{设计动机与边界：} 相对原始 3DGS 以独立高斯和图像重建误差为主的设定，3DGS.zip 把问题落在以下可实现的改动上：综述将3DGS压缩方法系统分为两大类。(1)压缩方法：ContextGS首次引入锚点级自回归模型用于3DGS压缩；HAC利用hash-grid辅助上下文框架压缩Scaffold-GS属性；CodecGS使用优化特征平面和标准视频编解码器；还有向量量化方法（如Compact3DGS将SH系数量化为码本）。(2)精简/致密化方法：Mini-Splatting引入Blur Split和深度重初始化；Octree-GS使用八叉树结构实现LOD；Taming3DGS采用全局评分指导高斯分裂/克隆；AtomGS通过原子高斯优先处理细节。综述还涵盖基础数学概念、实现细节和性能指标对比。其设计目的不是单独增加一个模块，而是用上述表示、监督或渲染约束消除原流程在目标任务中的具体失配。首次提供3DGS压缩和精简技术的全面综述，建立了统一的比较框架。系统分类了压缩方法（向量量化、熵编码、特征平面等）和精简方法（剪枝、LOD结构等），为该并行快速发展领域提供了清晰的路线图。适用时还应保留论文给出的边界：综述涵盖的方法发布时间密集，部分最新方法未纳入；不同方法的实验设置不统一，直接比较需谨慎
% METHOD: 2409.09756
\paragraph{MesonGS（后训练属性变换压缩）（2024.09）}
MesonGS~\cite{2409_09756} 是后训练压缩路线的代表。它的基线不是 Scaffold-GS 这类重训表示，而是已经完成 30k 迭代的原始 3DGS：输入是一个普通训练得到的高斯集合，输出是可直接解码渲染的压缩文件。它关注的实际问题是部署与传输：如果每个场景都要为了压缩重新训练一个特殊模型，训练成本会阻碍落地；因此 MesonGS 选择在训练后做剪枝、属性变换、量化和熵编码。

第一步是剪枝。MesonGS 的重要性分数同时包含视角相关项和视角无关项：前者统计高斯在训练视角的实际渲染贡献，考虑遮挡、透射率和像素覆盖；后者用体积等几何量避免只保留大而显眼的点。低重要性高斯被删除后，方法进入属性变换阶段：位置用 Octree 表示；旋转不再存四元数，而转成欧拉角以降低维度；不透明度、零阶球谐、尺度和角度等关键属性通过 region adaptive hierarchical transform（RAHT）降低空间相关属性的熵；高阶球谐则更适合用向量量化压缩。最后再做细粒度量化和 zlib/LZ77 一类通用无损压缩，并通过短暂 finetune 修复剪枝与量化带来的质量损失。

和 LightGaussian 相比，MesonGS 更像一个 codec：LightGaussian 把剪枝、蒸馏、VQ 串成学习型压缩流程，MesonGS 则更强调“拿到现成 3DGS 后如何转换成低熵数据”。和 RDO-Gaussian 这类率失真训练方法相比，MesonGS 的优势是后处理、易接入；代价是坐标通常不在 finetune 中更新，压缩选择也没有从原始训练开始共同优化。
\textbf{方法信息：} 代码：\url{https://github.com/ShuzhaoXie/MesonGS}；GitHub 约 28 stars；引用较多，Semantic Scholar 约 55 次。



% METHOD: 2501.12255
\paragraph{HAC++（哈希网格辅助上下文压缩）（2025.01）}
HAC++~\cite{2501_12255} 以 Scaffold-GS 为基线，进一步追求超过 100X 的压缩。Scaffold-GS 的 anchor 是无规则点集，直接熵编码时难以预测；HAC++ 的改变是学习一个结构化哈希网格作为上下文，把无序 anchor 放回可查询的空间结构中。

在编码每个 anchor 时，系统从 hash-grid 查询插值特征，用这些特征预测 anchor 属性的概率分布、量化步长和自适应 mask，而不是直接把 hash 特征当作新的外观表示。属性量化后进入算术编码，概率由上下文模型估计；AQM 决定不同属性在当前码率点上的量化精度，Adaptive Offset Masking 用可微 mask 把无效 offset 和 anchor 从码率计算中剔除。训练流程先稳定 Scaffold-GS anchor，再加入量化噪声，最后联合优化 hash-grid、概率模型和 mask；解码流程则按预测概率还原量化属性，再由 Scaffold-GS 的 MLP 生成实际高斯。Intra-anchor context 再沿特征通道做分块因果预测，利用同一 anchor 内部属性的相关性修正概率分布。相对 ContextGS 的层级自回归，HAC++ 更依赖“anchor 位置查询结构化网格”来预测属性分布，因此它把 Scaffold-GS 的无序锚点压缩问题转成了空间上下文建模问题。
\textbf{实现细节：} HAC++（2025.01） 将 Scaffold-GS 的锚点属性从"独立编码"转变为"条件编码"——不再逐锚点独立存储量化值，而是先利用结构化哈希网格提取空间上下文作为超先验（hyperprior），再以该先验为条件估计每个锚点属性的概率分布，从而大幅降低熵编码所需的比特数。
\textbf{设计动机与边界：} 相对原始 3DGS 以独立高斯和图像重建误差为主的设定，HAC++ 把问题落在以下可实现的改动上：HAC++基于Scaffold-GS构建，引入三大核心组件：(1) Hash-grid Assisted Context (HAC)：联合学习一个紧凑的结构化哈希网格，通过在锚点位置查询获得插值哈希特征，作为上下文预测锚点属性的值分布，用于后续的算术编码熵编码。同时输出自适应量化模块(AQM)的参数和高斯分布参数来建模值分布。(2) Intra-Anchor Context：消除锚点内部冗余，作为辅助信息提升HAC值分布的预测精度，增强整体压缩性能。(3) Adaptive Offset Masking：以可微分方式将掩码直接集成到码率计算中，无需额外的正则化损失项，使框架能在不同码率点达到最优掩码比率。为了消除原始 3DGS 在该目标任务中的具体失配，方法把上述表示、监督或渲染约束作为可学习和可执行的改动，而不是单独堆叠模块。提出HAC++框架，利用哈希网格辅助上下文建模和锚点内上下文关系，实现对3DGS的超100倍压缩。引入自适应量化模块提高属性量化精度，以及自适应掩码策略消除无效高斯。在所有数据集上平均实现超过100倍的大小缩减（相比原始3DGS），同时保持或提升渲染保真度。适用时还应保留论文给出的边界：依赖于Scaffold-GS的锚点结构，对于非锚点化的3DGS变体需要适配。哈希网格的学习增加了一定的训练开销。在极端低码率下可能存在细节损失。
\textbf{方法信息：} 引用很多，Semantic Scholar 约 50 次。



% METHOD: 2503.08511
\paragraph{PCGS（渐进式压缩）（2025.03）}
PCGS~\cite{2503_08511} 针对按需传输。固定码率压缩只能一次性下载一个完整模型，网络或存储变化时难以复用已有 bitstream；PCGS 的改动是让 3DGS 压缩同时具有数量渐进和质量渐进。

它基于 HAC++ 框架，数量上为每个 anchor 学习基础 mask 特征和层级特定 mask 特征，二者组合后产生 anchor 级与 Gaussian 级 mask；这些 mask 在层级之间被约束为单调非递减，因此已经解码的 anchor 不会在更高层级被丢弃，只会继续保留或细化。质量上，PCGS 不为每个码率重新训练模型，而是逐级减小量化步长，并用 trit-plane 方式把已有量化值继续细分。跨级上下文把低层已经解码的量化结果作为下一层概率预测的条件，高层只传新增 anchor、offset 或更精细属性所需的增量 bitstream。相对 LapisGS 的渲染层级，PCGS 的层级是压缩 bitstream 层级，目标是先下载低码率模型，再复用已有结果继续增强。
\textbf{实现细节：} 基于HAC++框架进行渐进扩展：(1) 渐进掩码策略（数量控制）：每个锚点学习一个基础掩码特征和一个级别特定掩码特征，组合生成锚点级和像素级掩码。掩码值在渐进级别间单调非递减，确保已解码的锚点不会被丢弃。每个级别可新增锚点或精化已有锚点。(2) 渐进量化策略（质量控制）：逐步减小量化步长，实现更精细的高斯属性建模。已有级别的量化结果作为下一级别的先验。(3) 跨级别上下文编码：利用已解码级别的量化结果改进下一级别的概率预测，提升熵编码效率。通过Hash grid捕获锚点间的上下文关系。 首次实现兼具高质量压缩和渐进性的3DGS压缩方法，通过从数量（渐进掩码）和质量（渐进量化）两个维度控制锚点，并利用跨级别上下文优化熵编码，一次训练即可获得完整的渐进RD曲线。 渐进级别数需预设；训练过程比单级压缩方法更复杂；在极低码率下的表现有提升空间。
\textbf{方法信息：} 代码：\url{https://github.com/yihangchen-ee/pcgs}。
\textbf{设计动机与边界：} 相对原始 3DGS 以独立高斯和图像重建误差为主的设定，PCGS 把问题落在以下可实现的改动上：基于HAC++框架进行渐进扩展：(1) 渐进掩码策略（数量控制）：每个锚点学习一个基础掩码特征和一个级别特定掩码特征，组合生成锚点级和像素级掩码。掩码值在渐进级别间单调非递减，确保已解码的锚点不会被丢弃。每个级别可新增锚点或精化已有锚点。(2) 渐进量化策略（质量控制）：逐步减小量化步长，实现更精细的高斯属性建模。已有级别的量化结果作为下一级别的先验。(3) 跨级别上下文编码：利用已解码级别的量化结果改进下一级别的概率预测，提升熵编码效率。通过Hash grid捕获锚点间的上下文关系。其设计目的不是单独增加一个模块，而是用上述表示、监督或渲染约束消除原流程在目标任务中的具体失配。首次实现兼具高质量压缩和渐进性的3DGS压缩方法，通过从数量（渐进掩码）和质量（渐进量化）两个维度控制锚点，并利用跨级别上下文优化熵编码，一次训练即可获得完整的渐进RD曲线。适用时还应保留论文给出的边界：渐进级别数需预设；训练过程比单级压缩方法更复杂；在极低码率下的表现有提升空间。
% METHOD: 2503.20221
\paragraph{TC-GS（三平面上下文压缩）（2025.03）}
TC-GS~\cite{2503_20221} 继承 Scaffold-GS 的 anchor 表示和 HAC 一类上下文熵编码思想，但把上下文来源从离散 hash-grid 换成三平面。它要解决的问题是：Scaffold-GS 的 anchor 在三维空间中稀疏、无序，直接用局部点邻域做概率建模时容易遇到哈希冲突、离散噪声和上下文不稳定；如果完全用神经网络预测属性分布，模型本身又会带来额外存储成本。TC-GS 的改变是把 anchor 坐标分别投影到 $xy/yz/xz$ 三个平面，在规则平面特征上读取空间上下文，再用这些上下文估计 anchor 特征、不透明度和局部尺度等属性的概率分布。

流程上，TC-GS 先用坐标 contraction 把不同尺度的 anchor 坐标映射到三平面有效范围内，再通过双线性插值得到平面特征；这些特征不直接替换 anchor 属性，而是作为 entropy model 的条件，预测每个待编码属性的均值、尺度或离散概率。直接用三平面替换 anchor 属性会在 anchor growing 阶段不稳定，因此论文仍保留 Scaffold-GS 的点结构，只让三平面承担“给概率模型提供上下文”的角色。解码时，方法进一步使用 $K=4$ 的邻近 anchor 信息，让位置相近的高斯在属性分布上互相提供先验；同时用 learnable mask 删除无效 anchor，并通过上采样/下采样模块降低三平面自身的存储量。

为了避免压缩后边缘和细纹理变软，TC-GS 还加入 adaptive wavelet loss：训练前期以重建稳定为主，后期逐渐提高高频小波项权重，使模型更关注压缩误差最容易破坏的细节区域。因此它相对 HAC/HAC++ 的差异不是简单换一个网络，而是把“无序点上下文”改写成“规则三平面上下文 + KNN 邻域修正 + 高频损失补偿”的编码管线；代价是训练时多了三平面和邻域解码开销，论文也指出 KNN 与 CPU 熵编码会拖慢实际编解码速度。
\textbf{方法信息：} 代码：\url{https://github.com/timwang2001/TC-GS}；GitHub 约 75 stars。
\textbf{设计动机与边界：} 相对原始 3DGS 以独立高斯和图像重建误差为主的设定，TC-GS 把问题落在以下可实现的改动上：TC-GS基于Scaffold-GS框架，核心创新包括：（1）三平面上下文模型：利用三平面（Tri-plane）表示利用锚点之间的固有空间相关性，将非结构化的高斯属性转化为结构化分布以进行压缩。通过三平面采样估计锚点属性分布。（2）KNN解码：在从三平面解码高斯分布时使用K最近邻方法利用相邻高斯之间的相关性。引入高斯位置信息作为位置敏感解码器的先验。（3）自适应小波损失：随着训练迭代增加逐渐关注高频细节，减轻压缩过程中高频信息的丢失。整体框架分为两部分：上层通过三平面采样估计锚点属性分布；下层使用带有熵编码的渲染管线实现压缩。还包括掩码方法减少过多高斯存储和三平面存储的进一步压缩。其设计目的不是单独增加一个模块，而是用上述表示、监督或渲染约束消除原流程在目标任务中的具体失配。提出了基于三平面的3DGS压缩方法TC-GS，利用三平面结构编码高斯属性以挖掘空间相关性，结合KNN解码和自适应小波损失，实现了与SOTA相当的压缩性能。适用时还应保留论文给出的边界：三平面分辨率的选择影响压缩率和质量的平衡；KNN计算增加了解码开销；对高斯分布高度非均匀的场景效果可能受限。
% METHOD: 2503.23162
\paragraph{NeuralGS（用神经场拟合高斯属性）（2025.03）}
NeuralGS~\cite{2503_23162} 的思路是把 NeRF 式神经场重新带回 3DGS 压缩中。它不改变最终渲染器，而是在存储层把大量高斯属性交给 MLP 拟合：位置作为输入，颜色、尺度、旋转等属性由网络解码。

它的关键观察是：如果直接用一个全局 MLP 从坐标预测所有高斯属性，误差会很大，因为相邻空间位置的高斯未必有相似尺度、旋转或颜色。NeuralGS 因此先做两件事。第一，用渲染贡献和归一化体积计算每个高斯的重要性，低贡献基元先被剪掉，剩余基元在拟合阶段按重要性加权。第二，不按空间位置聚类，而按高斯属性相似性做 K-means，使同一簇里的属性变化更接近一个小网络能拟合的函数。每个簇再训练一个轻量 MLP，以位置编码后的坐标为输入，输出不透明度、球谐颜色、尺度和旋转等属性；训练后只需保存高斯位置、簇分配和若干 MLP 参数。渲染前把属性一次性解码成普通 3DGS，再走标准 rasterization，因此它不会像逐视角神经解码那样拖慢每帧渲染。相对 F-3DGS 的低秩线性分解，NeuralGS 使用非线性神经场表达属性冗余；相对 Scaffold 类方法，它压缩的是已经训练好的原始 3DGS，而不是重新设计 anchor 表示。
\textbf{方法信息：} 项目仓库：\url{https://github.com/PKU-YuanGroup/NeuralGS}；GitHub 关注度较高，约 183 stars；截至 2026-07-09，仓库 README 仍把上传代码列为待办，当前主要提供项目说明和结果材料。
\textbf{设计动机与边界：} 相对原始 3DGS 以独立高斯和图像重建误差为主的设定，NeuralGS 把问题落在以下可实现的改动上：NeuralGS的核心思路：(1)全局重要性计算：基于每个高斯对所有训练视角像素的贡献计算重要性分数，结合归一化体积进行重要性排序；(2)聚类策略：按高斯属性分布进行聚类（而非简单的空间聚类），为每个簇拟合一个独立的MLP，解决高斯属性空间分布不均匀导致的拟合困难；(3)重要性感知拟合：使用重要性分数作为拟合权重，确保对渲染贡献大的高斯获得更好的拟合精度；(4)微调优化：在训练图像上使用光度损失和频率损失进行端到端微调，提升渲染质量。推理时只需将高斯位置通过MLP解码一次获取所有属性，保持原始3DGS的渲染效率。其设计目的不是单独增加一个模块，而是用上述表示、监督或渲染约束消除原流程在目标任务中的具体失配。提出将神经场与点云表示结合的压缩方案，通过聚类策略将高斯按属性分组并用不同MLP拟合，实现91倍平均模型大小缩减而不损害视觉质量，保持原始3DGS的渲染效率。适用时还应保留论文给出的边界：聚类数量需要预设；压缩后的渲染速度可能因MLP推理略有下降；对非常大规模场景的压缩效果待验证
% METHOD: 2504.13022
\paragraph{CompGS++（静态与动态统一预测压缩）（2025.04）}
CompGS++~\cite{2504_13022} 扩展了 CompGS 的“基元间预测”思想，使其同时覆盖静态和动态场景。静态 3DGS 的冗余主要在空间相邻基元之间，动态 3DGS 还多了时间冗余；CompGS++ 的改动是分别建立空间基元预测和时间基元预测，再用率约束优化减少基元内部属性冗余。

静态分支沿用“anchor primitives + coupled primitives”的混合基元结构，但把预测关系写得更系统：anchor primitive 保存完整参考信息，多个 coupled primitives 只保存紧凑残差嵌入；解码时由 anchor 的参考特征和 coupled 的残差特征共同预测后者的位置偏移、尺度、旋转、颜色和不透明度。实现材料中常见的设置是每个 anchor 派生若干 coupled primitives，参考嵌入维度明显高于残差嵌入维度；anchor 位置可交给 GPCC 这类点云几何编码器，其余连续属性由概率模型和算术编码处理。动态分支再引入 temporal primitive prediction：当前时刻的基元不再完全独立编码，而是由前一时刻或参考时刻的基元预测，编码器只保存无法由时间相关性解释的残差。率约束优化把重建误差和码率放进同一目标，量化步长可由先验特征自适应预测。相对 CompGS，它的增量不是换一个熵编码器，而是把“从 anchor 预测 coupled primitive”的空间关系扩展成“空间 + 时间”的统一预测框架。
\textbf{方法信息：} CompGS 相关代码：\url{https://github.com/LiuXiangrui/CompGS}；GitHub 约 46 stars；CompGS++ 论文说明实现将公开，需继续跟踪是否并入该仓库或另行发布。
\textbf{设计动机与边界：} 相对原始 3DGS 以独立高斯和图像重建误差为主的设定，CompGS++ 把问题落在以下可实现的改动上：CompGS++基于消除冗余原则设计：(1) 空间基元预测模块：建立场景基元间的预测关系，使大部分基元编码为紧凑残差，大幅减少空间冗余。利用高斯基元的空间相关性（相邻基元参数相似），通过预测编码替代直接存储。(2) 时间基元预测模块：处理动态场景，利用时间戳间的基元相关性减少时间冗余。跨时间戳预测基元变化。(3) 率约束优化模块：联合最小化重建误差和码率消耗，有效消除基元内参数冗余，增强场景表征的整体紧凑性。框架同时支持静态和动态场景的压缩，是首个完整的动态3DGS压缩方案。其设计目的不是单独增加一个模块，而是用上述表示、监督或渲染约束消除原流程在目标任务中的具体失配。提出消除基元间和基元内冗余的综合压缩框架，包含空间和时间基元预测模块以及率约束优化，在多个基准上显著超越现有压缩方法。适用时还应保留论文给出的边界：预测模块增加编码复杂度；对极端动态场景的压缩率可能下降；解码端需要预测模块的逆过程。
% METHOD: 2505.01938
\paragraph{HybridGS（生成式紧凑表示加标准点云编码）（2025.05）}
HybridGS~\cite{2505_01938} 针对的是压缩格式的实用性。许多 3DGS codec 使用定制神经熵模型，率失真好但编码慢、格式难标准化；HybridGS 的增量是把“生成紧凑 3DGS”和“标准点云编码”接起来。代码 README 也显示它建立在官方 3DGS 管线之上，并围绕 ICML 2025 论文实现双通道稀疏表示和下游点云编码。

具体做法分两步。第一步在 3DGS 生成阶段直接产出紧凑显式数据：位置、颜色、透明度、尺度和旋转仍保持可落盘的点云式属性，但颜色和旋转这类可压缩通道改用低维潜特征加一层轻量 MLP 解码，尺度这类对失真敏感的属性直接保留；训练时把量化器接入优化，让渲染损失看到量化后的误差，而不是等训练结束再硬量化。码率控制有两条手段：通过 progressive primitive pruning 调整高斯数量，通过位置和属性 bit-depth 调整每个基元的字节数。第二步把 \texttt{point\_cloud\_latent\_quan.ply}、反量化元数据和 latent decoder 交给 GPCC 做标准点云编码，输出可传输的 bitstream。相对 CompGS 的耦合基元和神经预测，HybridGS 更强调编码器兼容性和编解码速度；它不是单纯追求最低码率，而是在“可解释压缩程度、标准 codec、可部署速度”之间取平衡。
\textbf{设计动机与边界：} 相对原始 3DGS 以独立高斯和图像重建误差为主的设定，HybridGS 把问题落在以下可实现的改动上：HybridGS包含两个核心模块：(1) 双通道稀疏表示：通过特征通道降维(用低维潜特征+轻量MLP解码器替代PCA)和量化精度控制来生成紧凑的显式3DGS数据。具体地，颜色和旋转属性使用低维潜特征重建(共享潜特征和解码器)，缩放作为压缩敏感属性直接保留。量化采用训练感知的方式，受神经网络量化启发，在训练过程中即考虑量化影响；(2) 下游点云编码：利用标准GPCC编码器对紧凑3DGS数据进行进一步压缩，形成标准输出比特流。提出简单有效的码率控制方案来调节可解释的压缩程度。框架不含3DGS质量提升模块，但仍达到与SOTA可比的重建性能，且编解码速度显著更快。其设计目的不是单独增加一个模块，而是用上述表示、监督或渲染约束消除原流程在目标任务中的具体失配。提出HybridGS混合压缩框架，结合双通道稀疏表示和标准点云编码器，首次实现3DGS的生成式压缩与传统压缩的统一，在保持可比重建质量的同时显著提升编解码速度。适用时还应保留论文给出的边界：未集成3DGS质量提升模块，可能在某些复杂场景下质量略逊于生成式压缩SOTA；量化策略对超参数有一定敏感性；主要针对静态场景。
\textbf{方法信息：} Qualcomm 参与了这篇论文；代码：\url{https://github.com/qi-yangsjtu/hybridgs}。



% METHOD: 2506.01822
\paragraph{GSCodec Studio（模块化 3D/4D 高斯压缩框架）（2025.06）}
GSCodec Studio~\cite{2506_01822} 更像一个系统锚点。它不只提出单一 codec，而是把静态和动态 Gaussian Splat 压缩拆成可组合模块。项目框架建立在 gsplat 之上，包含 static/dynamic splats、重建、训练时压缩模拟和测试时压缩策略。

流程分为准备、带压缩模拟的训练、后训练压缩和查看四阶段。准备阶段解析 COLMAP 数据并初始化静态或动态 GS 表示；训练阶段把量化、熵约束和自适应 mask 等压缩模拟接进优化，使高斯在训练时就适应低码率；后训练阶段再用剪枝、量化、排序和熵编码生成紧凑 bitstream；查看阶段负责按用户视角实时解码和渲染。静态 GSCodec 通常以 gsplat/MCMC 训练为基础，把可微量化、熵约束、自适应掩码和训练后剪枝接入同一实验框架。动态 GSCodec 则采用时空高斯：位置可以由多项式运动描述，不透明度随时间变化，外观由特征分支预测，再复用静态分支中的压缩模拟和熵编码模块；在视频式设置中，它还支持按较短或较长 GOP 组织动态序列。相对只服务单篇论文的压缩方法，GSCodec Studio 的价值在于统一接口和可复现实验：读者可以把不同剪枝、量化、排序、熵编码策略作为模块比较，也能把静态 3DGS 压缩和 4DGS 压缩放在同一工程环境中评估。
\textbf{实现细节：} 框架设计包含四个阶段：(1) Preparation：解析COLMAP数据，初始化静态/动态GS表示。(2) Training with Compression Simulation：训练时集成压缩模拟（量化、熵约束、自适应掩码），使GS表示更紧凑和可压缩。(3) Post-Training Compression：使用GS编码器（剪枝、量化、排序、图像编码）生成紧凑比特流，解码器重建GS。(4) Viewing：基于用户视角实时渲染。Static GSCodec：基于gsplat+MCMC采样+可微量化+熵约束+自适应掩码+后训练剪枝和熵编码。Dynamic GSCodec：基于STG表示+多项式运动+时变不透明度+特征外观着色，复用Static GSCodec的压缩模块。 提出首个统一模块化的3D/4D高斯场压缩框架GSCodec Studio，整合多种重建和压缩技术为可组合组件，基于社区最佳实践提出Static GSCodec和Dynamic GSCodec，在静态和动态GS压缩上达到竞争力的率失真性能。 框架复杂度较高，不同模块组合的最优配置需大量实验；动态场景压缩效果仍有较大提升空间。
\textbf{设计动机与边界：} 相对原始 3DGS 以独立高斯和图像重建误差为主的设定，GSCodec Studio 把问题落在以下可实现的改动上：框架设计包含四个阶段：(1) Preparation：解析COLMAP数据，初始化静态/动态GS表示。(2) Training with Compression Simulation：训练时集成压缩模拟（量化、熵约束、自适应掩码），使GS表示更紧凑和可压缩。(3) Post-Training Compression：使用GS编码器（剪枝、量化、排序、图像编码）生成紧凑比特流，解码器重建GS。(4) Viewing：基于用户视角实时渲染。Static GSCodec：基于gsplat+MCMC采样+可微量化+熵约束+自适应掩码+后训练剪枝和熵编码。Dynamic GSCodec：基于STG表示+多项式运动+时变不透明度+特征外观着色，复用Static GSCodec的压缩模块。其设计目的不是单独增加一个模块，而是用上述表示、监督或渲染约束消除原流程在目标任务中的具体失配。提出首个统一模块化的3D/4D高斯场压缩框架GSCodec Studio，整合多种重建和压缩技术为可组合组件，基于社区最佳实践提出Static GSCodec和Dynamic GSCodec，在静态和动态GS压缩上达到竞争力的率失真性能。适用时还应保留论文给出的边界：框架复杂度较高，不同模块组合的最优配置需大量实验；动态场景压缩效果仍有较大提升空间。
\textbf{方法信息：} 代码：\url{https://github.com/JasonLSC/GSCodec_Studio}；GitHub 关注度较高，约 139 stars。



% METHOD: 2506.04174
\paragraph{FlexGS（一次训练，多尺寸部署）（2025.06）}
FlexGS~\cite{2506_04174} 解决的是固定压缩率的问题。许多剪枝方法要为某个目标大小重新剪枝和微调；FlexGS 的改变是训练一个可按目标比例选择高斯的弹性模型，让同一个训练结果可以部署到不同显存和带宽条件下。

方法包含两个轻量模块。自适应高斯选择器以目标压缩比例为条件，把每个高斯的属性和比例输入 GsNet，预测保留/丢弃的二值变量；训练时用 Gumbel-Softmax 把离散选择放进梯度优化，并在 1\%、5\%、10\%、15\% 等比例之间随机采样，让同一个模型学会多种预算。为了让选择器不只看局部透明度，FlexGS 定期更新 Global Importance，把空间覆盖、opacity 和透射贡献合成全局重要性信号。比例依赖的高斯变换场则只作用在被保留的高斯上，根据目标比例预测小位移，补偿删掉高斯后缺失的局部细节。训练通常先跑一段原始 3DGS，再联合优化选择器和变换场；推理时用户指定目标大小，模型直接输出对应子集，不需要为每个比例重新剪枝和微调。相对 MaskGaussian 的单一剪枝阈值，FlexGS 的增量是把“压缩率”变成可输入条件。
\textbf{实现细节：} FlexGS（2025.06） 与 ECO-GS 都关注 3DGS 在不同能力设备上的部署，但采取了截然不同的技术路线：ECO-GS 为每个设备预训练一个固定压缩比的本地小模型，部署时无法动态调整；FlexGS 则提出"训练一次，处处部署"的弹性推理范式——训练一个模型，推理时按任意比例动态选择高斯子集，无需额外微调即可适配不同设备的算力约束。
\textbf{设计动机与边界：} 相对原始 3DGS 以独立高斯和图像重建误差为主的设定，FlexGS 把问题落在以下可实现的改动上：FlexGS包含两个核心模块：(1)自适应高斯选择器：为每个高斯预测一个二值选择器，由轻量级MLP根据输入比率e决定是否保留，通过Gumbel-Softmax实现可微分选择，训练时使用随机采样比率增强泛化；(2)比率依赖高斯变换场：对被选中的高斯施加位移补偿，通过MLP根据比率输出位置偏移，弥补因减少高斯数量导致的细节损失。整体框架在训练阶段联合优化选择器和变换场，推理时根据目标比率直接输出压缩模型。其设计目的不是单独增加一个模块，而是用上述表示、监督或渲染约束消除原流程在目标任务中的具体失配。提出FlexGS，一种3DGS弹性推理方法，训练一次后无需微调即可根据目标模型大小自适应选择和变换高斯子集。通过轻量级可学习模块控制高斯选择，配合变换模块补偿减少高斯带来的性能损失，在ZipNeRF、MipNeRF和Tanks\&Temples场景上验证有效性。适用时还应保留论文给出的边界：选择器和变换场增加少量参数开销；极端压缩比率下质量下降明显；对训练时未覆盖的比率可能泛化不佳。
\textbf{方法信息：} NVIDIA 参与了这篇论文；代码：\url{https://github.com/liuhengyu321/flexgs}。



% METHOD: 2506.22973
\paragraph{Confident Splatting（基于置信度的可学习剪枝压缩）（2025.06）}
Confident Splatting~\cite{2506_22973} 把每个 splat 是否值得保留建模为可学习置信度，而不是用固定不透明度或梯度阈值。它为每个高斯引入 Beta 分布参数，平均置信度表示该高斯对场景重建的可信贡献。

具体来说，它为第 $i$ 个 splat 学习 Beta 分布参数 $\alpha_i,\beta_i$，用均值 $c_i=\alpha_i/(\alpha_i+\beta_i)$ 作为置信度，并在 rasterization 前把原始不透明度改成 $o_i c_i$。这一步很重要：置信度不是训练后另算的评分，而是在渲染时直接影响颜色合成，因此低置信度 splat 会在重建损失中承担较小贡献。训练目标由重建损失、稀疏正则、负熵项和显著性排序损失组成：稀疏正则鼓励少用 splat，负熵项避免置信度过早坍缩到 0/1，排序损失则抽样比较高低贡献 splat，使置信度排序与实际渲染贡献一致。训练结束后不需要重新训练压缩模型，用户只需调一个置信度阈值，就能在质量和数量之间扫出不同压缩点；同一机制也能接到 vanilla 3DGS 或 MCMC-GS 上。相对 LP-3DGS 和 MaskGaussian，它更强调置信度分布本身既是剪枝依据，也可以作为场景质量的诊断信号。
\textbf{方法信息：} 代码：\url{https://github.com/AmirHossein-razlighi/Confident-Splatting}；GitHub 约 9 stars。
\textbf{设计动机与边界：} 相对原始 3DGS 以独立高斯和图像重建误差为主的设定，Confident Splatting 把问题落在以下可实现的改动上：该方法为每个splat估计置信度分数，假设每个splat的置信度服从Beta分布（由alpha和beta两个参数化），通过联合优化重建损失和置信度参数。引入三种正则化器：稀疏性损失鼓励更少的活动splat；负熵损失避免过度置信；显著性排序损失采用相对方法而非绝对值来引导置信度优化。训练完成后，置信度低于阈值的splat被修剪。该方法可叠加在任何3DGS流程之上，无需修改原有管线。此外，平均置信度值可作为评估场景质量的新指标。其设计目的不是单独增加一个模块，而是用上述表示、监督或渲染约束消除原流程在目标任务中的具体失配。提出基于可学习Beta分布的置信度分数的有损压缩方法，每个splat的置信度通过重建感知损失优化，支持修剪低置信度splat同时保持视觉质量，方法与架构无关可应用于任何Gaussian Splatting变体。适用时还应保留论文给出的边界：Beta分布假设可能不适用于所有场景；修剪阈值需要人工设定；极端压缩率下视觉质量仍有下降。
% METHOD: 2507.05859
\paragraph{D-FCGS（动态高斯前馈压缩）（2025.07）}
D-FCGS~\cite{2507_05859} 面向自由视点视频中的动态高斯序列。静态 3DGS codec 只压缩一个场景，动态 3DGS 还要压缩帧间运动；D-FCGS 的增量是把动态高斯做成 I-P 结构的前馈压缩。

方法把连续 Gaussian frames 组织成 Group-of-Frames：I 帧保存或压缩完整高斯，P 帧参考前一帧。对一个 P 帧，D-FCGS 先用 farthest point sampling 从密集高斯中选出稀疏控制点，取当前帧和参考帧的控制点几何参数，经频率编码和 MLP 投影后得到 feature-domain motion tensor。这个稀疏 motion tensor 进入前馈压缩网络：训练时用均匀噪声模拟量化以保留梯度，测试时四舍五入量化，并由 arithmetic coding 输出 bitstream。熵模型同时看两类先验：hyperprior 估计全局 latent 分布，spatial-temporal context prior 则用参考帧和多分辨率 hash-grid 提供局部时空上下文。解码后，控制点 motion 通过邻域传播到全部高斯，再由轻量 color refinement network 修正跨视角外观不一致。相对 3DGStream 这类逐场景优化路线，D-FCGS 的重点是训练一个 scene-agnostic 的 P-frame codec；推理时可对未见过的 Gaussian sequence 做 zero-shot 压缩。
\textbf{方法信息：} 代码：\url{https://github.com/Mr-Zwkid/D-FCGS}；GitHub 约 20 stars。
\textbf{设计动机与边界：} 相对原始 3DGS 以独立高斯和图像重建误差为主的设定，D-FCGS 把问题落在以下可实现的改动上：D-FCGS框架包含三个关键组件：(1) 稀疏运动提取：通过最远点采样(FPS)从N个高斯中选取控制点，编码几何参数并提取运动张量（当前帧与参考帧的特征域差），大幅降低存储需求；(2) 前馈运动压缩：采用端到端压缩管线，包含数据编码、可微量化（训练时加均匀噪声，测试时四舍五入）和算术编码。提出双先验感知熵模型融合超先验和时空先验提升概率估计精度；(3) 运动补偿与细化：解码后通过控制点引导将稀疏运动传播到所有高斯帧，配合轻量级颜色细化网络提升视角一致保真度。采用GoF(Group-of-Frames)结构和I-P编码，I帧独立压缩，P帧参考前帧。训练在多视角视频生成的高斯帧上进行，可零样本泛化到未见场景。其设计目的不是单独增加一个模块，而是用上述表示、监督或渲染约束消除原流程在目标任务中的具体失配。提出D-FCGS，首个动态高斯泼溅的前馈压缩框架，采用GoF结构和I-P编码、双先验感知熵模型、控制点引导的运动补偿机制，实现零样本泛化的17倍以上压缩，匹配优化方法的率失真性能。适用时还应保留论文给出的边界：前馈模型对极端运动场景的泛化能力有限；控制点数量影响压缩效率和重建质量的平衡；长序列的误差累积问题。
% METHOD: 2508.14037
\paragraph{Distilled-3DGS（多教师蒸馏压缩）（2025.08）}
Distilled-3DGS~\cite{2508_14037} 把模型蒸馏引入 3DGS 压缩。普通分类模型可以对齐中间特征，但 3DGS 的高斯是无序显式基元，教师和学生很难逐点对应；该方法的改动是只在渲染结果和结构图上蒸馏。

它先训练教师集合：一个标准 3DGS 教师，一个在高梯度高斯上周期性加入随机扰动的 perturbation teacher，一个在训练中随机关闭部分基元的 dropout teacher。三类教师提供的不是同一个模型的重复预测，而是对几何细节、局部扰动和稀疏化鲁棒性的互补约束。学生模型先用重要性准则剪成更小的高斯集合，再用教师集合渲染出的伪图像做 RGB 监督；为了传递隐藏几何结构，论文还引入 spatial distribution distillation，把教师和学生的高斯分布投到体素统计或结构直方图中对齐，而不是强行做逐高斯匹配。这样压缩目标被写成“少量高斯复现教师渲染分布和空间分布”。相对纯剪枝方法，蒸馏允许学生重新组织高斯，而不仅仅是删掉原模型的一部分；代价是训练时需要多个教师和额外伪标签渲染，计算成本高于一次性剪枝。论文报告学生通常只保留原始 3DGS 约一成左右的高斯，仍能维持接近或优于若干压缩基线的渲染质量。
\textbf{方法信息：} 代码：\url{https://github.com/lt-xiang/Distilled-3DGS}；GitHub 约 82 stars。
\textbf{设计动机与边界：} 相对原始 3DGS 以独立高斯和图像重建误差为主的设定，Distilled-3DGS 把问题落在以下可实现的改动上：Distilled-3DGS是首个面向3DGS的知识蒸馏框架。核心设计：(1) 多教师训练：训练三个独立3DGS模型——标准3DGS、噪声增强变体（添加随机噪声到训练过程）和dropout正则化变体（训练时随机丢弃高斯），获得多样性的高质量教师模型。(2) 蒸馏流程：聚合教师模型输出生成伪图像，学生模型通过匹配伪图像进行训练。(3) 结构相似性损失：针对3DGS显式无序表示无法直接特征对齐的挑战，提出通过比较教师和学生的深度/不透明度图来传递隐藏几何结构信息。(4) 解决3DGS蒸馏的三大挑战：缺乏一致的潜在特征空间、高斯原语无序性、渲染输出对单个高斯不可微。方法简洁有效，无需复杂设计即可实现显著压缩。其设计目的不是单独增加一个模块，而是用上述表示、监督或渲染约束消除原流程在目标任务中的具体失配。首次将知识蒸馏应用于3DGS，提出多教师蒸馏框架和结构相似性损失，在保持渲染质量的同时大幅减少高斯数量和存储需求，为3DGS在资源受限设备上的部署提供新范式。适用时还应保留论文给出的边界：需要预先训练多个教师模型增加总体训练成本；学生模型容量下限受限于最小渲染质量要求。
% METHOD: 2509.03775
\paragraph{ContraGS（码本凝聚的可训练 3DGS）（2025.09）}
ContraGS~\cite{2509_03775} 关注的是“训练时就压缩”，而不是训练完再压缩。Compact 3D Gaussian、LightGaussian、RDO-Gaussian 等方法已经证明码本可以有效压缩训练好的 3DGS，但如果训练阶段仍保存每个高斯的完整属性，峰值显存和数据搬运开销并不会消失。ContraGS 的目标是在保持百万级高斯数量的同时，把球谐颜色、尺度和旋转放进共享码本，让训练、存储和渲染都使用压缩表示。

它把模型状态写成两部分：$G$ 是高斯集合，每个高斯保存位置、不透明度以及指向码本的整数索引；$C$ 是码本集合，包含 SH codebook 和 scale/rotation codebook。渲染时，高斯先根据索引查出颜色和几何向量，再按标准 3DGS/MCMC rasterization 生成图像。难点在于索引不可微，不能像普通连续参数那样直接用 SGD 更新。ContraGS 因此把参数估计改写成贝叶斯后验采样：后验概率由重建损失和码本规模正则共同决定，Metropolis--Hastings 采样器在“连续参数更新、码本向量 split、码本向量 merge”等状态转移之间切换，并根据 proposal 与 acceptance probability 决定是否接受新状态。这样，码本内容、码本大小和高斯到码本的映射都可以随训练推进改变。相对 Compact 3D Gaussian 的 VQ 后处理，ContraGS 更强调码本与 3DGS 训练同步，使压缩表示从一开始就是可训练的；论文报告平均降低训练峰值模型显存约 $3.49\times$，同时带来训练和渲染加速。
\textbf{实现细节：} ContraGS：码本直接训练。针对"先训练后压缩"范式带来的训练中期资源浪费问题——Compact 3DGS 的 R-VQ 码本仅在完整模型训练完毕后的最后 1000 次迭代中才以 K-means 初始化引入，导致训练中期所有高斯参数均以完整连续形式存在，需消耗大量 GPU 显存与计算资源，而这些本可在压缩表示上提前节省。Compact 3DGS 的 R-VQ 码本是在完整模型训练完毕后的最后 1000 次迭代中才引入的（以 K-means 初始化），这意味着先训练一个完整模型再压缩——浪费了大量存储和计算。ContraGS（2025.09） 的核心改动是将参数估计建模为贝叶斯推断问题，在码本压缩表示上直接训练，跳过"先训练完整模型再压缩"的两步流程。初始训练时每个高斯对应独立的码本向量（一对一映射），随着 MH 采样的进行，相似的向量被合并为共享码本条目，码本自然收缩。
\textbf{方法信息：} Intel 参与了这篇论文；截至 2026-07-09，未确认官方公开代码。
\textbf{设计动机与边界：} 相对原始 3DGS 以独立高斯和图像重建误差为主的设定，ContraGS 把问题落在以下可实现的改动上：ContraGS的核心创新：(1) 码本压缩状态空间——定义SH码本（球谐系数向量）和SR码本（缩放和四元数参数向量），每个高斯仅存储指向码本的整数指针和位置/不透明度，大幅减少参数量；(2) 贝叶斯推断框架——将码本压缩3DGS的学习建模为后验分布采样问题，后验正比于重建损失和码本正则化的指数；(3) MCMC采样——不同于3DGS-MCMC使用SGLD，ContraGS设计了新的状态转换、提议和接受概率用于码本指针和码本向量的联合采样，解决非可微参数的学习难题。这使得训练直接在压缩表示上进行，无需先训练完整模型再压缩。其设计目的不是单独增加一个模块，而是用上述表示、监督或渲染约束消除原流程在目标任务中的具体失配。提出首个在码本压缩表示上直接训练3DGS的方法，通过贝叶斯推断和MCMC采样解决码本压缩的非可微学习问题，实现3.49倍峰值内存降低、1.36倍训练加速和1.88倍渲染加速，同时保持接近SOTA的质量。适用时还应保留论文给出的边界：MCMC采样收敛需要调参；码本大小选择影响压缩率和质量的平衡。
% METHOD: 2512.07052
\paragraph{RAVE（连续码率自适应编码）（2025.12）}
RAVE~\cite{2512_07052} 处理的是压缩部署中的连续码率需求。许多方法需要为每个目标码率单独训练或单独选择参数；RAVE 的改变是在训练阶段学习上下界压缩上下文，推理时通过插值得到任意中间码率。

它可以作为插件接入不同 3DGS 表示和 codec。训练时先得到一个单码率 Scaffold-GS 压缩模型，再以此为起点训练多层 progressive anchor；每一层根据局部上下文重新排序新增高斯，而不是只按全局贡献排序。部署时，系统不必为 20 个目标码率分别训练 20 个模型，而是在相邻上下界 anchor 之间用插值系数生成中间模型：带宽低时选择更靠近低码率端的系数，带宽高时向高码率端移动。相对 FlexGS 主要控制渲染时使用哪些高斯，RAVE 更关注码率参数本身如何连续化，因此它服务的是可变带宽传输和率失真曲线的平滑控制。
\textbf{实现细节：} RAVE方法核心是模型无关的率自适应压缩：（1）动机：现有3DGS压缩方法在固定速率下操作，无法适应变化的带宽和设备约束。RAVE支持在预定义边界内任意速率插值。（2）核心方法：通过预定义的上下界压缩上下文训练模型，推理时可在两个极端之间任意插值获得中间速率。方案轻量级，仅需边际额外训练成本（少量额外迭代），无需为任何特定速率重新训练。（3）模型无关设计：可作为插件应用于不同3DGS表示和压缩方法。关键在于学习可插值的压缩参数化，使得在训练一次后通过简单的插值操作即可获得连续的率失真曲线。（4）应用场景：适合沉浸式应用中的实际部署，能根据网络带宽和设备能力动态调整压缩率。 首个支持3DGS在任意速率间灵活插值的压缩方案，无需重训练即可动态控制率失真权衡，计算轻量且保持广泛操作点的渲染质量。 需要训练上下界模型增加初始训练开销；极低比特率下质量可能仍受限；插值速率的率失真曲线可能非最优。
\textbf{实现补充：} 实现上还要核对输入、输出、参数、特征、初始化、训练损失、正则、渲染或推理接口，以及优化中的剪枝、分裂、采样、编码或解码步骤；这些抓手决定该方法到底改了 3DGS 管线的哪一环。
\textbf{方法信息：} 代码：\url{https://github.com/inspiros/rave}。
\textbf{设计动机与边界：} 相对原始 3DGS 以独立高斯和图像重建误差为主的设定，RAVE 把问题落在以下可实现的改动上：RAVE方法核心是模型无关的率自适应压缩：（1）动机：现有3DGS压缩方法在固定速率下操作，无法适应变化的带宽和设备约束。RAVE支持在预定义边界内任意速率插值。（2）核心方法：通过预定义的上下界压缩上下文训练模型，推理时可在两个极端之间任意插值获得中间速率。方案轻量级，仅需边际额外训练成本（少量额外迭代），无需为任何特定速率重新训练。（3）模型无关设计：可作为插件应用于不同3DGS表示和压缩方法。关键在于学习可插值的压缩参数化，使得在训练一次后通过简单的插值操作即可获得连续的率失真曲线。（4）应用场景：适合沉浸式应用中的实际部署，能根据网络带宽和设备能力动态调整压缩率。其设计目的不是单独增加一个模块，而是用上述表示、监督或渲染约束消除原流程在目标任务中的具体失配。首个支持3DGS在任意速率间灵活插值的压缩方案，无需重训练即可动态控制率失真权衡，计算轻量且保持广泛操作点的渲染质量。适用时还应保留论文给出的边界：需要训练上下界模型增加初始训练开销；极低比特率下质量可能仍受限；插值速率的率失真曲线可能非最优。
% METHOD: 2512.11186
\paragraph{Video Codec Compression（把高斯映射到视频块）（2025.12）}
Lightweight 3D Gaussian Splatting Compression via Video Codec~\cite{2512_11186} 追求标准硬件 codec 友好。它的核心问题是高斯点云无序，而视频编码器擅长处理规则二维块和时间连续帧；因此方法先把高斯属性重新排列成二维块，再交给成熟视频编码器。

LGSCV 的起点是 MPEG/I-3DGS 压缩实践：如果用 PLAS 把几十万高斯排成平滑 2D map，率失真不错，但 map 生成本身可能需要几十秒 GPU 时间，不适合轻量设备。它先用两阶段 Morton scan 替代重型 PLAS：3D Morton scan 按空间坐标排列高斯，2D Morton scan 再把一维序列按块映射到二维图像，使相邻空间高斯尽量落在视频 codec 的相邻 coding units 中。单靠 Morton map 在高码率下接近 PLAS，但中低码率会崩，因此论文再对 45 维高阶 SH AC 系数做 PCA，通常只保留少量主成分，把颜色属性能量集中到更少通道；最后用 MiniPLAS 在小块内做少量置换，既降低块内属性差异，也给 HEVC/HM 这类 codec 的 CU size 选择提供约束。相对 HybridGS 使用点云编码器，这条路线直接借用视频硬件生态；论文在 MPEG 测试序列上报告约 20\% RD gain，并把 2D map 生成时间降到秒级。
\textbf{实现细节：} LGSCV（2025.12） 沿着 GSICO"借用成熟编解码器生态"的思路继续推进，但将后端从图像编解码器升级为视频编解码器（如 HEVC），以利用帧间预测进一步挖掘多帧 3DGS 参数之间的冗余。然而，直接套用 GSICO 的方案会遇到一个实际瓶颈：GSICO 所依赖的 PLAS 排序虽然能产生高质量的 2D 映射，但其计算代价相当昂贵。打个比方，PLAS 排序就像是用"显微镜"逐一检查每对高斯之间的关系来确定最优排列——结果虽然精确，但耗时也成正比。LGSCV 的核心贡献正是用一个更轻量的替代方案来绕过这个瓶颈。两阶段 Morton 扫描替代 PLAS 排序。与 GSICO 的 PLAS 排序不同，LGSCV 提出了两阶段 Morton 扫描来建立 3D 高斯到 2D 图像的映射。Morton 码（也称 Z-order curve）的思想很简单：将三维坐标的二进制表示逐位交错（interleave），生成一个单一整数，使得空间上接近的点在排序后索引也接近。
\textbf{方法信息：} Qualcomm 参与了这篇论文；代码：\url{https://github.com/Qi-Yangsjtu/LGSCV}；GitHub 约 13 stars。
\textbf{设计动机与边界：} 相对原始 3DGS 以独立高斯和图像重建误差为主的设定，Lightweight 3D Gaussian Splatting Compression via Video Codec 把问题落在以下可实现的改动上：LGSCV包含三个核心创新：（1）两阶段Morton扫描：先用3D Morton扫描对高斯基元按空间坐标排序，再用2D Morton扫描将排序后的基元以块级方式映射到2D网格，生成的块级映射对标准视频编解码器的方形编码单元友好。（2）SH PCA降维：对球谐函数进行主成分分析，降低高维SH系数的维度，消除高阶交流分量的冗余。（3）MiniPLAS：灵活快速的局部排序方法，在特定块大小内对基元进行排列优化。MiniPLAS还可指导编解码器CU大小配置，进一步减少编码时间。其设计目的不是单独增加一个模块，而是用上述表示、监督或渲染约束消除原流程在目标任务中的具体失配。提出两阶段Morton扫描生成块级2D映射替代计算昂贵的PLAS，结合SH PCA降维和MiniPLAS，实现超过20%的RD增益，同时将2D映射生成时间降至约1秒，编码时间减少50%。适用时还应保留论文给出的边界：在极低比特率下质量仍有下降；MiniPLAS的块大小需要手动配置；主要针对静态场景压缩
% METHOD: 2512.17528
\paragraph{Voxel-GS（Octree 与游程编码）（2025.12）}
Voxel-GS~\cite{2512_17528} 反其道而行之：不依赖复杂神经熵模型，而是把 Scaffold-GS 属性离散为适合传统编码的数据。它的改动是将可微量化、轻量率代理和 Octree 加 run-length coding 组合在一起。

训练时，Voxel-GS 从 Scaffold-GS 的 anchor/attribute 表示出发，把连续属性通过可微量化转成整数型值；为了让训练目标知道这些整数最终会被 run-length coding 压缩，论文设计了基于 Laplacian 分布的 rate proxy，用它估计相邻量化值形成长游程时的实际码率，并把该估计作为熵约束加入重建损失。这样，模型不只是把属性四舍五入，而是倾向生成更容易被 RLC 表示的整数序列。编码阶段，几何位置通过 Octree 表示 occupied voxels，量化属性或索引再用游程编码无损压缩；解码时先还原 Octree 点云和属性序列，再恢复 Scaffold-GS 渲染所需的高斯。相对 HAC++ 的 learned context，Voxel-GS 更轻、更可解释，优势在于编码链路简单、速度快；代价是它依赖 Scaffold-GS 结构，压缩能力主要来自量化与游程规律，而不是复杂神经概率模型。
\textbf{方法信息：} 代码：\url{https://github.com/zb12138/VoxelGS}；GitHub 约 5 stars。
\textbf{设计动机与边界：} 相对原始 3DGS 以独立高斯和图像重建误差为主的设定，Voxel-GS 把问题落在以下可实现的改动上：Voxel-GS采用三步压缩流程：(1)可微量化：将Scaffold-GS的高斯属性(位置、颜色、尺度、旋转等)通过可微量化操作离散化为整数型值；(2)基于拉普拉斯的率代理：设计轻量级率代理准确估计游程编码的比特率，施加熵约束引导生成高保真且紧凑的重建；(3)Octree+游程编码：将整数型高斯点云通过Octree结构和游程编码进行无损压缩。关键创新是率代理能准确估计实际编码比特率，从而有效消除冗余并优化表示紧凑性。其设计目的不是单独增加一个模块，而是用上述表示、监督或渲染约束消除原流程在目标任务中的具体失配。提出Voxel-GS压缩框架，放弃复杂的神经熵模型，使用可微量化离散化Scaffold-GS高斯属性，设计基于拉普拉斯的率代理施加熵约束，最后使用Octree和游程编码进行无损压缩，实现显著压缩比和远超现有方法的编码速度。适用时还应保留论文给出的边界：依赖于Scaffold-GS的锚点结构；量化参数需要手动调节；对极高压缩率下质量保持有限。
% METHOD: 2604.26799
\paragraph{MesonGS++（后训练压缩与超参数搜索）（2026.04）}
MesonGS++~\cite{2604_26799} 继承 MesonGS 的后训练压缩路线，但把配置选择也系统化。后训练压缩常有许多耦合超参数：剪枝保留率、属性变换、位宽、量化组、熵编码设置。MesonGS++ 的改变是把这些选择建模为有约束的优化问题。

编码管线从已经训练好的 3DGS 出发，不重新训练完整模型。第一步用联合重要性分数剪枝：视角相关项衡量高斯在像素合成中的实际遮挡和贡献，视角无关项用体积等几何信息避免误删结构性基元。第二步把保留高斯的位置体素化，并用 octree/G-PCC 类几何编码器压缩坐标。第三步对属性分组：不透明度、尺度、旋转和 0 阶 SH 等重要属性经过 RAHT 变换以降低局部相关性；高阶 SH 系数更适合选择性向量量化；不同属性组再分配不同 bit-width。最后，量化符号进入熵编码。配置端是 MesonGS++ 的关键：它把保留率和各组位宽视为主要率失真旋钮，用线性 size estimator 估计目标文件大小，再用离散采样和 0-1 整数线性规划在存储预算下选择组合；CUDA 并行量化用于快速评估候选配置。相对 MesonGS，它不仅给 codec，也给“怎么选 codec 参数”的自动化方案，适合有明确目标大小的传输和归档场景。
\textbf{实现细节：} MesonGS++（2026.04） 是目前后训练压缩路线中最完整的管线，依次包含剪枝、属性变换、量化和熵编码四个阶段。其核心创新在于尺寸感知的超参数搜索：用户只需指定目标存储预算，MesonGS++ 自动搜索最优的高斯保留率（$\tau$）和分组量化位宽（$\mathbf{b}$），无需人工调参。压缩管线的四个阶段各有关键设计。在剪枝阶段，MesonGS++ 设计了联合重要性指标 $I_g = I_d \cdot I_i$，将视角依赖得分（考虑渲染时高斯的排序与遮挡关系）与视角无关得分（基于归一化体积）相乘，比仅用不透明度或梯度作为剪枝准则更准确地衡量每个高斯对最终渲染的贡献。在量化阶段，MesonGS++ 对 RAHT 变换后的 AC 系数进行分组量化（group-wise quantization），每个量化组 $c$ 按如下方式量化： 其中 $b$ 为位宽，$S_c$ 为缩放因子，$Z_c$ 为零点偏移——即先将属性值缩放到 $[0, 2^b{-}1]$ 的整数范围，超出范围的值被截断。不同属性组可以使用不同的位宽——例如实验发现不透明度和旋转对量化更敏感，通常被分配更高位宽，而 RGB 和缩放则使用较低位宽。最后在熵编码阶段，八叉树坐标用 G-PCC 压缩，其余整数元素通过算术编码（torchac）和 LZ77 进行压缩打包。直觉地说，压缩管线的超参数（保留率 $\tau$、各组的量化位宽 $\mathbf{b}$、码本大小）之间存在复杂的耦合关系——降低保留率意味着保留更多高斯但留给每个高斯的存储预算更少，需要更低位的量化来补偿；反之，提高保留率可以给每个高斯分配更多位宽，但高斯总数减少可能损失细节。MesonGS++ 将这一搜索问题形式化为混合整数非线性规划（Mixed Integer Nonlinear Programming, MINLP），目标是在满足存储预算约束 $S(\tau, \mathbf{b}) \leq \mathrm{Budget}$ 的条件下最小化质量损失 $Q(\tau, \mathbf{b})$。
\textbf{技术实现补充：} MesonGS++（2026.04） 是目前后训练压缩路线中最完整的管线，依次包含剪枝、属性变换、量化和熵编码四个阶段。其核心创新在于尺寸感知的超参数搜索：用户只需指定目标存储预算，MesonGS++ 自动搜索最优的高斯保留率（$\tau$）和分组量化位宽（$\mathbf{b}$），无需人工调参。压缩管线的四个阶段各有关键设计。在剪枝阶段，MesonGS++ 设计了联合重要性指标 $I_g = I_d \cdot I_i$，将视角依赖得分（考虑渲染时高斯的排序与遮挡关系）与视角无关得分（基于归一化体积）相乘，比仅用不透明度或梯度作为剪枝准则更准确地衡量每个高斯对最终渲染的贡献。在量化阶段，MesonGS++ 对 RAHT 变换后的 AC 系数进行分组量化（group-wise quantization），每个量化组 $c$ 按如下方式量化： 其中 $b$ 为位宽，$S_c$ 为缩放因子，$Z_c$ 为零点偏移——即先将属性值缩放到 $[0, 2^b{-}1]$ 的整数范围，超出范围的值被截断。不同属性组可以使用不同的位宽——例如实验发现不透明度和旋转对量化更敏感，通常被分配更高位宽，而 RGB 和缩放则使用较低位宽。最后在熵编码阶段，八叉树坐标用 G-PCC 压缩，其余整数元素通过算术编码（torchac）和 LZ77 进行压缩打包。直觉地说，压缩管线的超参数（保留率 $\tau$、各组的量化位宽 $\mathbf{b}$、码本大小）之间存在复杂的耦合关系——降低保留率意味着保留更多高斯但留给每个高斯的存储预算更少，需要更低位的量化来补偿；反之，提高保留率可以给每个高斯分配更多位宽，但高斯总数减少可能损失细节。MesonGS++ 将这一搜索问题形式化为混合整数非线性规划（Mixed Integer Nonlinear Programming, MINLP），目标是在满足存储预算约束 $S(\tau, \mathbf{b}) \leq \mathrm{Budget}$ 的条件下最小化质量损失 $Q(\tau, \mathbf{b})$。
\textbf{方法信息：} 代码：\url{https://github.com/mmlab-sigs/mesongs_plus}；GitHub 约 6 stars。
\textbf{设计动机与边界：} 相对原始 3DGS 以独立高斯和图像重建误差为主的设定，MesonGS++ 把问题落在以下可实现的改动上：MesonGS++包含编码端和配置端两部分。编码管线五阶段：(1) 联合重要性剪枝：结合视角依赖和视角无关的重要性分数，识别可剪枝高斯。(2) 八叉树几何编码：对体素化坐标使用八叉树结构。(3) 属性变换与选择性向量量化：RAHT变换重要属性，对高阶SH系数使用向量量化。(4) 分组混合精度量化：不同属性组使用不同位宽。(5) 熵编码。配置优化：将保留率和位宽分配建模为MINLP问题，通过离散采样和0-1整数线性规划求解。提出线性大小估计器和CUDA并行量化算子加速超参数搜索。其设计目的不是单独增加一个模块，而是用上述表示、监督或渲染约束消除原流程在目标任务中的具体失配。提出首个按目标存储预算精确控制压缩大小的3DGS后训练编解码器，实现34倍压缩同时保持渲染保真度，甚至在Stump场景上20倍压缩率下超越vanilla 3DGS的PSNR，无需任何重新训练。适用时还应保留论文给出的边界：后训练方法，无法利用训练时优化；超参数搜索需要额外时间；对极端压缩率的渲染质量损失较大。

\section{高斯剪枝与稀疏化}

剪枝路线的锚点是原始 3DGS 的 densification。3DGS 为了补细节会不断 clone/split 高斯，最终许多基元贡献很小、互相遮挡或只解释训练噪声。不同方法的差别在于“重要性”怎么定义：不透明度和梯度、二阶不确定性、可学习 mask、概率置信度、最优传输合并、语义 mask 或渲染无关几何近似，都对应一种不同的剪枝观。
% METHOD: 2311.17245
\paragraph{LightGaussian（剪枝、蒸馏与属性量化三步走）（2023.11）}
LightGaussian~\cite{2311_17245} 是 3DGS 压缩方向早期最重要的完整流水线之一。它的基线是原始 3DGS：每个高斯保存位置、尺度、旋转、不透明度和高阶球谐，点数多且属性冗余。LightGaussian 的改变不是只做单一剪枝，而是把压缩拆成三件事：先减少高斯数量，再降低外观阶数，最后压缩剩余属性。

第一步是 Gaussian pruning and recovery。方法为每个高斯计算全局显著性：一方面统计它在多视角渲染中的像素贡献和命中次数，另一方面结合体积、不透明度等几何/透明度因素，避免只凭单个视角或单个属性误删。低显著性高斯被删除后，剩余模型继续训练约 5000 步，让位置、尺度、不透明度和颜色重新适配新的点集。第二步是 SH distillation。原始 3DGS 依赖高阶球谐表达视角相关颜色，LightGaussian 通过伪视图渲染和蒸馏，把高阶球谐中的部分外观变化转移到低阶表示中，从而降低每个高斯的颜色参数量。第三步是 Gaussian vector quantization：对剩余属性建立 codebook，对贡献较低或冗余度较高的部分使用更激进的 VQ 比例，并配合熵编码存储索引和码本。

它的重要性在于给后续工作建立了一个可复用锚点：如果一个方法说自己改进“3DGS compression”，通常要说明它是在改显著性分数、恢复训练、球谐蒸馏、属性量化，还是把这些步骤合并进端到端训练。PUP 3D-GS 继承了 LightGaussian 的 prune-refine 框架但换成二阶敏感性分数；MaskGaussian 则把“删除”改成概率 mask，使训练中仍能评估被暂时屏蔽的高斯；RDO-Gaussian 则把属性保留与码率目标显式放进优化。
\textbf{实现细节：} 问题与动机。3DGS 的 ADC（自适应密度控制）策略在训练过程中不断克隆和分裂高斯以增强场景覆盖，但缺乏对高斯总量的有效约束。它只有"增加"机制（clone/split），没有"删除"机制（prune），导致训练后高斯数量从数千膨胀到数百万。一个直接的想法是"按不透明度剪枝"：移除不透明度最低的高斯。然而，不透明度只是高斯的局部属性，一个不透明度较低但覆盖面积很大的高斯可能对大量像素有不可忽视的颜色贡献。LightGaussian（2023.11） 受深度学习中网络剪枝（network pruning，通过移除不重要的神经元或滤波器来压缩网络的技术）的启发，提出需要一个全局视角来评估每个高斯的真实重要性。这要求不仅看它"自身有多透明"，还要看它"在所有训练视角的渲染中实际贡献了多少颜色"。其核心思想是：一个高斯的重要性不取决于它在单张图像中的表现，而取决于它对所有训练视角渲染结果的累积影响。全局显著性分数。具体而言，LightGaussian 的定义从 3DGS 的 $\alpha$-混合渲染方程出发。直觉上，3DGS 的渲染过程是一个逐层叠加的 $\alpha$-混合： \[ C(\mathbf{p}) = \sum\_{i=1}\^{}{N} \alpha\_i T\_i \mathbf{c}\_i \] 其中 $\alpha_i$ 为第 $i$ 个高斯在像素 $\mathbf{p}$ 处的不透明度，$T_i = \prod_{j=1}^{i-1}(1-\alpha_j)$ 为透射率（即该高斯前方所有高斯"漏过来"的光线比例），$\mathbf{c}_i$ 为颜色。LightGaussian 的关键洞察是：第 $i$ 个高斯对像素 $\mathbf{p}$ 的有效颜色贡献为 $\alpha_i T_i \mathbf{c}_i$，贡献权重 $\alpha_i T_i$ 同时取决于高斯自身的不透明度和它前方所有高斯的遮挡。因此，LightGaussian 在所有训练视角的所有像素上累积每个高斯的贡献权重，得到全局显著性分数： 其中 $\mathcal{V}$ 为训练视角集合，$\mathcal{P}_v$ 为视角 $v$ 中的像素集合，$\mathbb{1}(g_i, r_{\mathbf{p}})$ 是指示函数（当高斯 $g_i$ 与射线 $r_{\mathbf{p}}$ 相交时为 1，否则为 0），$V_{\mathrm{norm}}(g_i)$ 为归一化体积项。用白话解释这个公式：对每个高斯，数一数"它被多少条渲染射线命中"、每次命中时"它的不透明度有多大贡献"、以及"它的物理体积有多大"。该分数越高，说明高斯 $g_i$ 对渲染结果的贡献越大，越不应被剪枝。
\textbf{方法信息：} 代码：\url{https://github.com/VITA-Group/LightGaussian}；GitHub 星标很高，约 803 stars；引用很多，Semantic Scholar 约 568 次。



% METHOD: 2312.04564
\paragraph{EAGLES（轻量编码与粗到细量化）（2023.12）}
EAGLES~\cite{2312_04564} 是另一条早期轻量化路线：它不把主要矛盾只归结为“高斯太多”，而是指出每个高斯内部的颜色球谐和旋转参数占据了大量存储。它的基线仍是原始 3DGS，但改动发生在属性表示上：位置、尺度和低频颜色这类对几何和基础外观敏感的属性尽量保留，而高维颜色和旋转被压入可量化 latent vector，再由轻量 MLP 解码回渲染需要的属性。

训练时，EAGLES 为被压缩属性学习连续 latent，并在前向中量化这些 latent；反向传播用 straight-through estimator 近似离散量化的梯度。这样模型不是训练完再被动压缩，而是在训练期间就适应“属性必须通过有限码字表达”的限制。不透明度还会做量化处理，论文认为这能缓解不透明度饱和和漂浮点问题。与此同时，EAGLES 引入 coarse-to-fine 训练：先在较低分辨率或较粗尺度上稳定优化点云，再逐步提高分辨率，使量化属性、剪枝和几何调整不会在早期互相干扰。

完整流程包括三块：量化 latent 编码降低单点属性成本，粗到细训练提高优化稳定性，周期性剪枝移除低贡献高斯。相对 LightGaussian 的后处理 VQ，EAGLES 更强调“编码表示参与训练”；相对 Compact-3DGS 的 mask/网格式压缩，它更像一个面向单高斯属性的轻量编码器。
\textbf{方法信息：} 代码：\url{https://github.com/Sharath-girish/efficientgaussian}；GitHub 星标较高，约 176 stars；引用很多，Semantic Scholar 约 246 次。
\textbf{设计动机与边界：} 相对原始 3DGS 以独立高斯和图像重建误差为主的设定，EAGLES 把问题落在以下可实现的改动上：EAGLES针对3DGS的内存瓶颈（每场景>1GB存储、训练时约20GB GPU内存），提出三个核心创新：（1）属性量化：使用量化潜向量表示高斯体属性，通过Straight-Through Estimator（STE）实现端到端训练。对颜色（球谐系数）和旋转属性进行量化编码（使用MLP解码器恢复），不量化位置、缩放和基带颜色SH系数（因对初始化敏感，量化后性能下降大）。不透明度量化不仅减少内存还改善优化过程，避免梯度过大导致的不透明度饱和为0或1引起的伪影；（2）粗到细训练策略：先在低分辨率下训练粗粒度模型，再逐步提升分辨率细化，提升训练稳定性和收敛速度，获得更好的重建质量；（3）剪枝阶段：针对3DGS频繁密化（通过克隆和分裂）产生的冗余高斯体，引入剪枝机制移除不必要的高斯体，减少表示大小，加快训练和渲染速度。训练完成后，将潜向量取整为整数并使用熵编码高效存储。其设计目的不是单独增加一个模块，而是用上述表示、监督或渲染约束消除原流程在目标任务中的具体失配。提出EAGLES方法，通过量化嵌入大幅减少3DGS的存储内存（10-20倍压缩），粗到细训练策略提升训练稳定性和收敛速度，剪枝阶段减少冗余高斯体。在保持视觉质量的同时实现更快的训练和推理速度。适用时还应保留论文给出的边界：量化可能对极端细节场景有轻微质量损失；MLP解码器引入少量额外推理开销；部分属性（位置、缩放）未量化限制了进一步压缩
% METHOD: 2405.18784
\paragraph{LP-3DGS（学习式剪枝 mask）（2024.05）}
LP-3DGS~\cite{2405_18784} 是学习剪枝路线的简洁代表。它的出发点是：RadSplat、Mini-Splatting 和 LightGaussian 这类方法通常要手动设定剪枝比例或重要性阈值，而不同场景的最优比例并不一致；为了找到合适阈值，研究者常常需要对同一场景做多轮训练和 sweep。LP-3DGS 的改变是把“剪枝比例”变成训练中自动学习的二值 mask。

它并不重新设计一个新的重要性分数，而是在已有分数之上加一个可训练 mask。每个高斯先根据 RadSplat、Mini-Splatting 或其他方法得到重要性；LP-3DGS 再用 Gumbel-Sigmoid 产生近似二值的保留/删除决策，使离散 mask 在反向传播中可优化。训练目标仍以渲染误差为主，同时通过 mask loss 推动模型减少高斯数量。论文中 mask 只在训练后段短时间学习，例如从约 19500 到 20000 迭代更新重要性和 mask，随后按学习到的 mask 剪掉高斯，再继续优化剩余点集。

相对手工阈值剪枝，LP-3DGS 的价值在于一次训练内为每个场景找到较合适的模型大小；相对 Compact-3DGS 的 STE mask，它用 Gumbel-Sigmoid 提供更平滑的梯度；相对 MaskGaussian，它的 mask 更直接，主要解决“阈值不想手调”，没有 masked rasterization 那种让未采样高斯持续接收潜在贡献梯度的反悔机制。
\textbf{方法信息：} 代码：\url{https://github.com/ASU-ESIC-FAN-Lab/LP-3DGS}；GitHub 约 62 stars；引用较多，Semantic Scholar 约 58 次。
\textbf{设计动机与边界：} 相对原始 3DGS 以独立高斯和图像重建误差为主的设定，LP-3DGS 把问题落在以下可实现的改动上：LP-3DGS的核心是将剪枝比例从手动超参数转变为可学习参数。（1）问题分析：现有3DGS剪枝方法（LightGaussian、RadSplat、MiniSplatting）需手动设定剪枝比例或阈值，且最优值因场景而异，需多轮训练试错。（2）可学习二值掩码：为每个高斯引入一个重要性分数和二值掩码。掩码=1保留高斯，=0剪枝。掩码通过sigmoid函数参数化，用Gumbel-Sigmoid技术使离散的二值决策在训练中可微分。（3）Gumbel-Sigmoid：相比传统直通估计器（STE），Gumbel-Sigmoid提供更平滑的梯度估计，与3DGS的联合训练更稳定。训练过程中掩码逐渐从连续值收敛到二值决策。（4）兼容性：LP-3DGS可与不同类型的重要性评分（基于不透明度、覆盖范围、颜色贡献等）结合使用。训练结束后掩码自然收敛到最优剪枝比例。其设计目的不是单独增加一个模块，而是用上述表示、监督或渲染约束消除原流程在目标任务中的具体失配。提出LP-3DGS，首次将可学习二值掩码应用于3DGS剪枝。使用Gumbel-Sigmoid替代直通估计器（STE）使掩码可微分且兼容现有3DGS训练流程。自动为每个场景找到最优剪枝比例，无需手动调参和多轮训练。适用时还应保留论文给出的边界：训练收敛速度可能比固定比例方法慢；Gumbel-Sigmoid的温度参数仍需设定；极端压缩比下效果有限。
% METHOD: 2406.10219
\paragraph{PUP 3D-GS（二阶不确定性剪枝）（2024.06）}
PUP 3D-GS~\cite{2406_10219} 继承 LightGaussian 的“剪枝后再恢复”框架，但把剪枝分数从启发式显著性换成二阶敏感性。它的问题意识很明确：不透明度、可见次数或一阶梯度只能说明高斯在当前渲染中的局部表现，不能可靠判断它的空间参数被扰动或删除后会不会显著增加训练视角重建误差。特别是在高压缩率下，启发式剪枝容易误删前景细节，导致细杆、边缘和小物体消失。

方法把不确定性定义为“高斯空间参数可以被扰动而不改变训练视图重建损失的程度”。直接计算所有高斯参数的 Hessian 不可行，所以 PUP 在模型收敛后用 Fisher 信息矩阵近似 $L_2$ 重建误差的 Hessian，并只保留每个高斯自己的块对角项。进一步地，它发现只用位置均值和尺度参数就足以构成有效的 spatial sensitivity pruning score；该分数可理解为当前高斯在投影几何中的局部曲率体积，曲率小、可扰动空间大的高斯更不敏感，更适合删除。

实际管线是后处理式的：输入一个已训练好的 3DGS，先在训练视角上计算 Fisher/敏感性分数，按分数删除低敏感高斯，再微调剩余高斯恢复质量；随后可以再次计算分数并做第二轮剪枝。论文采用 80\% 后再从剩余部分剪 50\% 的两轮策略，总计删除约 90\% 高斯，并报告平均渲染速度提升约 $3.56\times$。相对 Trimming the Fat 的一阶梯度积累，PUP 更重视二阶曲率；相对 MaskGaussian，它不学习可反悔 mask，而是给后训练剪枝提供更强的统计判据。
\textbf{实现细节：} 与 LightGaussian 的一阶累积贡献不同，PUP 3D-GS（2024.06） 改用二阶灵敏度分析来评估每个高斯的剪枝影响。LightGaussian 的全局显著性分数回答的是"每个高斯对渲染贡献了多少颜色"，但这个问题隐含一个盲点：贡献大不代表移除后误差就大。如果周围有其他高斯能补偿它的贡献，那移除它的实际损失可能很小。PUP 3D-GS 转而回答一个更精确的问题："如果扰动某个高斯的参数，重建误差会增加多少？" 从 Fisher 信息矩阵获取灵敏度。考虑 $L_2$ 重建损失关于高斯参数的 Hessian 矩阵（即二阶导数矩阵）$\nabla_G^{}2 L_2$。对 Hessian 做展开后会发现两项：一项是渲染图像梯度 $\nabla_G I_G$ 的外积，另一项含重建残差 $(I_G - I_{\mathrm{gt}})$ 的二阶导数。关键洞察是：在已收敛的 3DGS 模型上，残差趋近于零，因此第二项消失，Hessian 退化为： 这正是统计学中的Fisher 信息矩阵——它只依赖于输入相机位姿（不依赖输入图像），因此计算成本远低于完整的 Hessian。注意，这一近似的前提是模型已经充分收敛。如果残差仍然较大，则 Fisher 近似会低估真实的灵敏度。分块灵敏度与空间参数选择。完整的 Fisher 矩阵对所有 $G$ 个高斯的参数求二阶导，规模为二次方，计算代价极高。PUP 3D-GS 发现只需保留每个高斯自身的参数块（即块对角近似 $\mathbf{H}_i = \nabla_{G_i} I_G\, \nabla_{G_i} I_G^{}\top$，忽略高斯间的交叉项）就足以衡量单个高斯的灵敏度。这相当于假设各高斯之间的影响是可分离的。进一步，论文发现\textbf{只需使用空间均值 $\mathbf{x}_i$ 和缩放 $\mathbf{s}_i$ 参数}即可获得有效的灵敏度分数，不需要旋转和不透明度参数。作者推测，这是因为场景到输入视角的投影几何不变性主要由位置和尺度决定，旋转在不变量存在时不会改变 3D 几何。最终的灵敏度分数定义为该参数块 Fisher 矩阵的对数行列式： 直觉上，对数行列式衡量的是该高斯参数周围"误差盆地"（error basin）的体积——值越小，说明该高斯的参数可以在很大范围内变化而不影响重建误差，因此它很可能是冗余的。Patch-wise 计算与多轮管线。
\textbf{方法信息：} 代码：\url{https://github.com/j-alex-hanson/gaussian-splatting-pup}；GitHub 星标较高，约 145 stars；引用较多，Semantic Scholar 约 87 次。
\textbf{设计动机与边界：} 相对原始 3DGS 以独立高斯和图像重建误差为主的设定，PUP 3D-GS 把问题落在以下可实现的改动上：PUP 3D-GS提出了一种原则性的3DGS剪枝方法。核心是空间敏感度分数，通过二阶近似重建误差对高斯空间参数的Hessian矩阵来量化每个高斯的重要性。直接计算完整Hessian不可行，因此采用Fisher信息矩阵近似：在收敛的3DGS模型上，L2损失的Hessian矩阵中的二阶项趋近于零，剩下的梯度外积即为Fisher信息矩阵。仅需块对角近似即可量化高斯敏感度，且基于图像块计算已足够。基于此敏感度分数进行剪枝，并设计多轮剪枝-细化管线：每轮先按敏感度分数剪枝低重要性高斯，然后微调剩余高斯。这种多轮策略允许在第一轮剪枝80%后进行第二轮再剪枝50%（总计90%），而单次剪枝90%会导致质量严重下降。该方法是后处理的，可应用于任何预训练3DGS模型而不改变训练管线。其设计目的不是单独增加一个模块，而是用上述表示、监督或渲染约束消除原流程在目标任务中的具体失配。提出了一种基于二阶近似的原则性敏感度剪枝分数，通过Fisher信息矩阵近似Hessian矩阵来评估每个高斯对重建误差的空间敏感度。设计了多轮剪枝-细化管线，可在不改变训练管线的情况下应用于任何预训练3DGS模型。剪枝90%高斯后渲染速度提升3.56倍，同时保持更好的图像质量。适用时还应保留论文给出的边界：Fisher信息矩阵近似假设模型已收敛，对于未充分收敛的模型效果可能下降。多轮剪枝增加了总处理时间。块对角近似可能遗漏高斯间的交互信息。
% METHOD: 2410.23213
\paragraph{ELMGS（GAP 剪枝、QAT 与熵编码）（2024.10）}
ELMGS~\cite{2410_23213} 目标是同时降低内存和计算。它的基线是训练完成的原始 3DGS，问题是最终模型里仍有大量不透明度很低、梯度也不活跃的冗余高斯；这些高斯既占存储，又增加每帧需要排序和混合的 primitive 数。ELMGS 把部署压缩拆成三步：Gradient and Opacity Aware Pruning（GAP）、量化感知微调（QAT）和熵编码。

GAP 阶段同时看不透明度和梯度的分位数阈值：只有当一个高斯的不透明度低、对应梯度也低时，才认为它对当前重建贡献小且不值得继续优化。剪枝不是一次性完成，而是在若干轮中逐步达到目标稀疏度；每轮剪枝后都继续 finetune，使剩余高斯调整协方差、颜色和不透明度来填补被删点的贡献。论文强调慢速迭代剪枝优于一次性剪枝，因为后者容易让优化落入质量较差的局部状态。

QAT 阶段使用 Learned Step Size Quantization（LSQ）让量化步长也参与优化：球谐等属性被压到低比特表示，训练过程显式感知量化误差，而不是训练后直接四舍五入。最后，压缩后的属性序列再通过熵编码进一步减少文件大小，README 中也把命令拆成 baseline 训练、\texttt{prunning\_play.py}、\texttt{QAT.py} 和 \texttt{compress.py} 几个入口。相对单一剪枝方法，ELMGS 的定位是部署型组合：先减 primitive 数，再减每个 primitive 的比特数，最后让编码器利用量化后的低熵分布。
\textbf{方法信息：} 代码：\url{https://github.com/salmanali96/elmgs}；GitHub 约 3 stars。
\textbf{设计动机与边界：} 相对原始 3DGS 以独立高斯和图像重建误差为主的设定，ELMGS 把问题落在以下可实现的改动上：ELMGS包含三个核心步骤：(1) Gradient and Opacity Aware Pruning (GAP)：基于不透明度和梯度的分位数函数设定阈值，当高斯的不透明度和梯度均低于阈值时进行剪枝，并周期性微调以恢复剪枝带来的质量损失。(2) 量化感知微调：将高斯参数从32位浮点数压缩到低位整数表示，通过可微量化层在训练中感知量化误差。(3) 熵编码：利用可微熵编码估计器在优化过程中考虑压缩率，使模型参数分布更适合熵编码压缩。三步策略可叠加使用，显著减少模型大小。其设计目的不是单独增加一个模块，而是用上述表示、监督或渲染约束消除原流程在目标任务中的具体失配。提出ELMGS框架，通过梯度/不透明度感知剪枝、量化感知微调和熵编码三步策略，实现3DGS模型的内存和计算可扩展性，在保持渲染质量的同时大幅减少存储需求适用时还应保留论文给出的边界：压缩过程需要额外的迭代剪枝和微调步骤，增加了总训练时间；极端压缩比下仍会有一定质量损失
% METHOD: 2412.20522
\paragraph{MaskGaussian（概率 mask 与可反悔剪枝）（2024.12）}
MaskGaussian~\cite{2412_20522} 把剪枝从确定性删除改成概率采样。它的基线包括 LightGaussian/RadSplat 这类手工重要性剪枝，也包括 Compact-3DGS、LP-3DGS 这类可学习 mask。核心问题是：一旦某个高斯在某个训练时刻被确定性删除，后续优化状态变化后，系统没有机会重新判断它是否应该回来；如果它解释的是细小前景或细杆结构，这种早删会造成不可恢复的细节损失。

它为每个高斯学习一个由 logits 参数化的二分类 mask 分布，并用 Gumbel-Softmax/Gumbel-Sigmoid 采样得到当前迭代是否使用该高斯。关键不在于“有 mask”本身，而在于 masked rasterization 的位置设计：MaskGaussian 不把 mask 乘到尺度或不透明度上，因为那样会让高斯在 $\alpha$ filtering 之前被过滤掉，无法收到梯度；它让所有高斯先正常 splat 并通过过滤，再在颜色累积和透射率演化中应用 mask。当前 mask 为 0 的高斯不会影响最终颜色，也不会消耗透射率，但仍能通过被零掉的潜在颜色贡献得到关于“如果我存在会不会减少误差”的梯度。

因此 MaskGaussian 的训练流程是动态的：如果其他高斯已经补偿了某个 masked 高斯的缺失，它得到的梯度会让存在概率继续降低；如果缺失导致重建误差，它的存在概率会升高，之后更可能被采样并参与渲染。论文既支持从零训练，也支持对已有 3DGS 做 post-training pruning，并报告在 Mip-NeRF 360、Tanks and Temples、Deep Blending 上平均剪掉 60\% 以上高斯且 PSNR 下降约 0.02。相对 LP-3DGS，它更重视“被暂时屏蔽的高斯仍可被重新评估”。
\textbf{方法信息：} 代码：\url{https://github.com/kaikai23/MaskGaussian}；GitHub 约 99 stars。



% METHOD: 2505.05587
\paragraph{SteepGS（最陡密度控制）（2025.05）}
Steepest Descent Density Control for Compact 3D Gaussian Splatting~\cite{2505_05587} 从优化理论重新解释 densification。原始 3DGS 的 Adaptive Density Control 主要按位置梯度和尺度阈值决定 split/clone：它能补细节，但并不判断一个高斯是否真的处在“梯度推不动、但分裂后可以下降”的位置。结果是许多已经能被普通梯度继续优化的高斯也被分裂，点数增长中相当一部分可能只是冗余复制。SteepGS 的改变，是把“何时分裂”从启发式阈值改成局部二阶下降判据。

具体做法是为每个候选高斯构造 splitting matrix，把渲染损失对该高斯参数的一阶梯度和局部 Hessian 近似联系起来。如果这个矩阵是半正定的，分裂不能带来额外下降，优化器继续更新原高斯即可；如果存在负曲率方向，说明高斯位于鞍点附近，两个子高斯沿最小特征值对应的正负方向对称移动，能够让后续梯度下降重新发挥作用。子代的不透明度不再简单继承，而是用解析归一化保持局部密度守恒，避免分裂本身改变总贡献。

训练流程仍接在原始 3DGS 管线里：从 SfM 初始化开始，每隔固定步数检查候选高斯，CUDA 侧计算 splitting matrix 和特征方向，只对满足负曲率条件的高斯执行分裂；未满足条件的高斯留给常规参数优化。这样它与 LightGaussian、REFINE 这类训练后剪枝不同，不是先长大再删，而是在增长发生前就减少无效 split。相对 Taming 3DGS 的预算化工程控制，SteepGS 更像给 densification 提供理论选择器：Taming 控制“最多长多少”，SteepGS 判断“谁值得长”。实验中它报告约 50\% 的高斯点数减少，同时保持接近原始 ADC 的渲染质量。
\textbf{设计动机与边界：} 相对原始 3DGS 以独立高斯和图像重建误差为主的设定，Steepest Descent Density Control for Compact 3D Gaussian Splatting 把问题落在以下可实现的改动上：SteepGS的理论贡献和方法为：(1) 引入分裂矩阵概念，连接每个高斯的一阶梯度和对应Hessian，证明分裂操作对逃离鞍点的关键作用；(2) 建立密度化的必要条件：确定何时分裂有帮助、子代高斯的最少数量、最优参数更新方向；(3) 提供子代不透明度归一化的解析解；(4) 基于以上理论洞察，提出最陡密度控制策略：在最小化损失的同时保持紧凑点云，实现硬件高效实现。其设计目的不是单独增加一个模块，而是用上述表示、监督或渲染约束消除原流程在目标任务中的具体失配。从非凸优化理论角度揭示3DGS密度控制机制，证明分裂对逃离鞍点的必要性，提出最陡密度控制策略SteepGS，在保持渲染质量的同时减少约50%高斯点数。适用时还应保留论文给出的边界：理论分析基于一定简化假设；对动态场景的适用性未验证；实际实现中的超参数选择仍需经验调整。
\textbf{方法信息：} Meta Reality Labs 参与了这篇论文；代码：\url{https://github.com/facebookresearch/SteepGS}；GitHub 关注度较高，约 106 stars。



% METHOD: 2506.09534
\paragraph{Gaussian Herding（最优传输视角的高斯合并）（2025.06）}
Gaussian Herding across Pens（GHAP）~\cite{2506_09534} 不把压缩看作删点，而是看作高斯混合模型归约。它针对的是剪枝方法的一个根本缺陷：LightGaussian、PUP 或 Trimming the Fat 都要在已有高斯里选择“留下谁、删除谁”，但原始 3DGS 中多个邻近高斯往往共同描述一个局部分布，删除其中一部分并不等价于用更少高斯重新近似这个分布。GHAP 的改变是把位置、协方差和不透明度解释为一个未归一化 Gaussian mixture，把压缩改写成“用更少混合组件近似原混合分布”。

具体流程分两阶段。第一阶段做几何压缩：用 KD-tree 将场景划分为多个局部区域，在每个区域里最小化 composite transport divergence（CTD），把若干原高斯的质量传输到更少的新高斯上。这个过程类似在“高斯空间”中做 $k$-means：分配步骤根据两个高斯分布之间的传输代价把原组件分到目标组件，更新步骤重新计算目标高斯的均值、协方差和权重。它不是保留某个原高斯，而是生成新的、更能代表局部分布的高斯组件。第二阶段再接回 3DGS 的标准训练，只微调颜色和不透明度，让几何归约后的点集恢复照片级外观。

GHAP 可以作为插件接到 vanilla 3DGS、Mini-Splatting、AtomGS 或 3DGS-MCMC 的不同阶段：既可以在构建过程中压缩，也可以对已训练 checkpoint 做压缩。论文重点展示了只保留约 10\% 高斯时仍能接近 vanilla 3DGS 的渲染质量，并在同等 primitive 数下优于多种剪枝式压缩。相对 PUP 的“判断谁不敏感”，GHAP 的局部锚点是“用最优传输合并出新的几何代表”。
\textbf{实现细节：} 第二，高效代价函数。Zhang 等（2025.06） 指出，直接用 KL 散度作为 $c(\cdot,\cdot)$ 虽在理论上最优，但需要 $O(s^{}3 n)$ 次协方差矩阵求逆且存在数值不稳定问题。他们设计了如下高效替代： 即代价等于均值向量的 $L^{}2$ 距离加上协方差矩阵的 Frobenius 范数距离。该代价函数无需矩阵求逆，且更新步骤退化为加权平均：$\tilde{\mu}_j^{(t)} = \frac{\sum_i \pi_{ij} w_i \mu_i}{\sum_i \pi_{ij} w_i}$，$\tilde{\Sigma}_j^{(t)} = \frac{\sum_i \pi_{ij} w_i \Sigma_i}{\sum_i \pi_{ij} w_i}$。整个 GMR 求解过程本质上是高斯空间中的 $k$-means 聚类变体：分配步骤将每个高斯分到最近的聚类中心，更新步骤将中心替换为聚类内高斯的加权 barycenter。
\textbf{实现补充：} 实现上还要核对输入、输出、参数、特征、初始化、训练损失、正则、渲染或推理接口，以及优化中的剪枝、分裂、采样、编码或解码步骤；这些抓手决定该方法到底改了 3DGS 管线的哪一环。
\textbf{方法信息：} 代码：\url{https://github.com/DrunkenPoet/GHAP}；GitHub 约 48 stars。



% METHOD: 2601.00913
\paragraph{Clean-GS（语义 mask 引导的伪高斯清理）（2026.01）}
Clean-GS~\cite{2601_00913} 不是只追求压缩率，而是清理 floaters 和背景伪影。它的基线是普通训练后的 3DGS 点云：在户外纪念碑、雕像或文物场景中，原始 3DGS 往往会把背景、游客、弱纹理区域和重建空洞解释成大量漂浮高斯。这些高斯不一定低不透明度，也不一定低梯度，所以 LightGaussian/PUP 这类全局重要性剪枝不能可靠地区分“目标物体”和“背景伪影”。Clean-GS 的改变是引入稀疏语义 mask，把剪枝目标从“保留渲染贡献”改成“只保留属于目标对象的高斯”。

流程分三步。第一步是 whitelist filtering：给少量训练视角提供二值语义 mask，例如论文中 302 个视角只用 3 个 mask；系统把每个高斯中心投影到这些视角，只要它至少一次落在目标区域内，就进入白名单，否则直接删除。第二步是 depth-buffered color validation：只渲染白名单高斯，并对每个 mask 视角保留前景深度最近的高斯，再比较该高斯由 DC/SH 颜色得到的 RGB 与 mask 图像中对应像素的颜色；如果一个可见高斯颜色长期不匹配，就更可能是投影到目标附近的伪影。第三步是邻域异常值移除：对剩余高斯做 kNN 分析，删除空间上孤立、邻域距离异常或只在极少视角偶然通过白名单的点。

这条路线的适用场景与一般压缩剪枝不同：它会主动改变场景内容，因此论文不把 PSNR/SSIM 当作唯一指标，而是关注对象隔离、文件减小和视觉清洁度。Temple 示例中，模型可从约 526K 高斯、125 MB 清理到约 198K 高斯、47 MB；Isha Statue 示例从约 1.11M 高斯、263.6 MB 到约 223K 高斯、52.8 MB。它的局限也直接来自语义监督：mask 视角覆盖不足会误删被遮挡部位，背景颜色接近目标时 color validation 可能保留一些伪影。
\textbf{方法信息：} 代码：\url{https://github.com/smlab-niser/clean-gs}；GitHub 约 24 stars。



% METHOD: 2603.09277
\paragraph{Shorter Gaussian Lists（缩短每像素高斯列表）（2026.03）}
Speeding Up the Learning of 3D Gaussians with Much Shorter Gaussian Lists~\cite{2603_09277} 关注渲染和训练中的每像素高斯列表长度。3DGS 的计算并不只取决于总高斯数，也取决于每个像素需要 alpha blending 的高斯数量；该方法的改变是让高斯更小、更集中，减少长列表。

它包含三项设计。周期性 scale reset 按固定因子缩小所有高斯尺度，使每个高斯覆盖更少像素；由于覆盖面积变小后仍要解释同一辐射贡献，训练会相应提高有效不透明度，从而让高斯更集中。alpha blending 权重的熵约束直接作用在每条射线的混合权重上，压低大量次要权重、强化少数主导权重，减少“许多低权重高斯共同解释一个像素”的长列表。渲染分辨率调度器则从低分辨率到高分辨率逐步训练，并按阶段调整 scale reset 和熵约束强度。相对普通剪枝，它不只是少高斯，而是少“每个像素要处理的高斯”，因此加速点落在 rasterizer 的 per-pixel work 上。
\textbf{设计动机与边界：} 相对原始 3DGS 以独立高斯和图像重建误差为主的设定，Speeding Up the Learning of 3D Gaussians with Much Shorter Gaussian Lists 把问题落在以下可实现的改动上：方法包含三个核心组件：(1)缩放重置——定期（如每N轮迭代）将所有高斯的缩放按固定比例缩小（如乘以0.8），鼓励更小高斯覆盖更少像素，缩短高斯列表。重置也鼓励更大的不透明度，因为较小的高斯需要更高不透明度来维持相同渲染效果。相比体积惩罚损失，缩放重置更简单且超参数更易调节；(2)熵约束——对alpha混合过程中的权重分布施加熵约束，驱动主导权重更大而次要权重更小。使每个高斯更集中于其主导像素，减弱对邻近像素的影响，导致更短的高斯列表。增加建模效率同时改善渲染质量；(3)渲染分辨率调度器——将缩放重置和熵约束集成到渐进分辨率增加的训练策略中，低分辨率阶段快速建立大致结构，高分辨率阶段精化细节，进一步提升整体训练效率。其设计目的不是单独增加一个模块，而是用上述表示、监督或渲染约束消除原流程在目标任务中的具体失配。提出两种互补策略缩短渲染像素的高斯列表以加速3DGS学习：(1)缩放重置——定期按比例缩小高斯尺寸，鼓励更小高斯覆盖更少像素；(2)alpha混合的熵约束——锐化沿射线的高斯权重分布，使主权重更大、次权重更小，每个高斯更集中于主导像素。结合渲染分辨率调度器进一步提升效率。在Mip-NeRF 360上保持渲染质量的同时显著提升效率。适用时还应保留论文给出的边界：缩放重置比例需要调参；熵约束权重过大可能影响渲染质量；与某些高斯密集化策略的兼容性需验证。
\textbf{方法信息：} NVIDIA 参与了这篇论文；代码：\url{https://github.com/machineperceptionlab/shortersplatting}。
\section{结构化与锚点表示}

结构化表示路线的共同目标是给无序高斯建立层次或锚点。Scaffold-GS 是这条线的核心锚点：它不再把所有高斯都当作独立基元，而是用 anchor 和特征生成视角自适应高斯。后续 SAGS、ProtoGS、V3DG、CLoD-GS 将结构化思想分别推向稀疏重建、原型共享、大规模 LoD 和连续距离衰减。


% METHOD: 2312.00109
\paragraph{Scaffold-GS（锚点生成式高斯）（2023.12）}
Scaffold-GS~\cite{2312_00109} 是结构化表示路线的核心锚点。原始 3DGS 把每个高斯都当成独立原语：位置、尺度、旋转、不透明度和球谐颜色分别存储，视角变化主要靠 SH 系数解释。这种方式直接、快速，但冗余也明显：相邻高斯常共享局部结构和外观规律，却仍各自保存完整属性。Scaffold-GS 的改变是把“每个高斯独立存属性”改成“稀疏 anchor 存特征，再按视角生成局部 neural Gaussians”。

具体实现从初始化开始。Scaffold-GS 仍以 COLMAP/SfM 点云为入口，但不把每个 SfM 点直接变成一个自由高斯，而是先按体素大小 $v_s$ 量化点云，取每个非空体素中心作为 anchor 位置 $\mathbf{x}_v$。每个 anchor 存三类状态：32 维局部上下文特征 $\mathbf{f}_v$、3 维缩放因子 $\mathbf{l}_v$，以及 $k$ 个可学习偏移 $\{\mathbf{O}_i\}_{i=0}^{k-1}$，论文实验默认 $k=10$。这一步相当于把原始 3DGS 的“每个基元都是独立个体”改成“每个局部区域先有一个支架，再由支架生成若干子高斯”。读者要注意，anchor 本身不是最终用于 $\alpha$-混合的高斯；它更像局部生成器的坐标和特征缓存。

为了让同一个局部区域在远近视角下使用不同细节，Scaffold-GS 还把 anchor 特征扩展成三级特征库 $\{\mathbf{f}_v,\mathbf{f}_{v/2},\mathbf{f}_{v/4}\}$。渲染某一相机 $\mathbf{x}_c$ 时，先计算 anchor 到相机的距离 $d_{vc}=\|\mathbf{x}_v-\mathbf{x}_c\|_2$ 和观察方向 $\boldsymbol{\delta}_{vc}=(\mathbf{x}_v-\mathbf{x}_c)/\|\mathbf{x}_v-\mathbf{x}_c\|_2$，再用小型权重网络 $F_w(d_{vc},\boldsymbol{\delta}_{vc})$ 预测三层特征的融合权重，得到视角相关特征 $\tilde{\mathbf{f}}_v$。因此，Scaffold-GS 的视角自适应不是简单给颜色多加几个 SH 系数，而是在生成高斯属性之前就让 anchor 特征感知距离和方向：近处可以偏向细节特征，远处可以偏向更平滑的上下文。

真正参与渲染的 neural Gaussian 由 anchor 按需解码。对视锥内每个可见 anchor，第 $i$ 个子高斯的位置写成
\[
\boldsymbol{\mu}_i=\mathbf{x}_v+\mathbf{O}_i\odot\mathbf{l}_v ,
\]
其中偏移 $\mathbf{O}_i$ 决定子高斯在局部支架周围的位置，缩放因子 $\mathbf{l}_v$ 控制这个局部坐标系的空间范围。其余属性，包括不透明度 $\alpha_i$、颜色 $\mathbf{c}_i$、旋转四元数 $\mathbf{q}_i$ 和尺度 $\mathbf{s}_i$，由独立的两层 MLP 从 $\tilde{\mathbf{f}}_v$、$d_{vc}$ 和 $\boldsymbol{\delta}_{vc}$ 解码；这些 MLP 的隐藏层维度为 32，使用 ReLU。这样，原始 3DGS 中长期存储在每个高斯上的大部分属性，变成了 anchor 特征和轻量网络的输出。渲染时只对视锥内 anchor 做前向计算，并且把预测不透明度低于阈值的 neural Gaussian 过滤掉，再送入原始 3DGS 风格的 tile rasterizer 和 $\alpha$-混合管线，所以额外 MLP 开销被限制在可见且有贡献的局部区域。

anchor 的 grow/prune 是它相对原始 3DGS 致密化最重要的实现差异。原始 3DGS 根据屏幕空间梯度对自由高斯做 clone/split，再周期性剪掉低不透明度或过大的高斯；Scaffold-GS 不直接复制自由高斯，而是在 neural Gaussian 的空间量化网格中统计一段训练窗口内的梯度贡献。论文默认每 $N=100$ 次迭代累积一次梯度均值 $\bar{g}$，若某个体素的贡献超过阈值 $\tau_g$，就在该体素中心新增 anchor；为了支持多粒度覆盖，体素大小和阈值还按层级调整，例如 $v_g^{(m)}=v_g/4^{m-1}$、$\tau_g^{(m)}=\tau_g\cdot 2^{m-1}$。剪枝则聚合一个 anchor 生成的 neural Gaussians 的不透明度，如果累计贡献低于 0.5，就移除该 anchor。直观地说，原始 3DGS 是“哪个高斯梯度大就复制哪个高斯”，Scaffold-GS 是“哪个局部区域的支架不够用，就补一个支架；哪个支架长期生成不出有用高斯，就删掉支架”。

训练目标仍围绕图像重建展开，通常继承原始 3DGS 的 $L_1$ 与 D-SSIM 组合，同时加入轻量体积或正则项来稳定 anchor 表示。论文在 30k 迭代设定下与原始 3DGS 对齐比较，并报告在 Mip-NeRF 360 上质量相当或略优，在 Tanks and Temples、Deep Blending、BungeeNeRF 等包含弱纹理、反射或多尺度观察的场景上更稳；它的优势不是单纯“更少高斯”，而是通过 anchor 特征共享减少冗余，并把视角相关变化放到解码阶段处理。对后续压缩方法来说，这一步非常关键，因为无序高斯点云变成了可建模的 anchor 级结构。ContextGS 可以让粗层 anchor 预测细层 anchor，HAC++ 可以在 anchor 位置查询 hash-grid 上下文，TC-GS 可以把 anchor 属性投影到 tri-plane，Voxel-GS 也能围绕 anchor feature 做量化传输。读这些方法时，先抓住 Scaffold-GS 的基线：anchor 负责存空间和特征，MLP 负责生成视角相关高斯；后续方法通常不是推翻这件事，而是压缩 anchor、预测 anchor 属性、组织 anchor 层级或减少需要生成的 offset。
\textbf{锚点细化：} Scaffold-GS：~是锚点表示的奠基性工作，也是后续~HAC++、ContextGS~等熵编码方法的基础架构。其核心思想是将场景组织为由锚点定义的局部高斯簇。每个锚点存储一个紧凑的特征向量，根据当前视角和距离动态预测其管辖范围内子高斯的完整属性，而非像原始~3DGS~那样为每个高斯静态存储所有参数。
\textbf{实现细节：} Scaffold-GS（2023.12） 是锚点表示的奠基性工作，也是后续 HAC++、ContextGS 等熵编码方法的基础架构。其核心思想是将场景组织为由锚点定义的局部高斯簇。每个锚点存储一个紧凑的特征向量，根据当前视角和距离动态预测其管辖范围内子高斯的完整属性，而非像原始 3DGS 那样为每个高斯静态存储所有参数。直觉地说，Scaffold-GS 将 3DGS 从"每个高斯独立行动"的组织模式转变为"层级管理"模式。在原始 3DGS 中，每个高斯独立存储 59 维参数（位置 3 维 + 四元数 4 维 + 缩放 3 维 + 不透明度 1 维 + 球谐系数 48 维），相邻高斯往往具有相似的颜色和几何，却各自保存一份完整副本，这正是冗余的根源。锚点则扮演管理者的角色：它知道"我的辖区内应该有什么样的高斯"，但具体的属性值根据当前视角动态生成，而非静态存储。这种设计带来三个优势：（1）参数共享，同一锚点下的 $k$ 个子高斯共享锚点特征，无需重复存储相似属性；（2）视图自适应，高斯属性根据视角动态调整，近处取高分辨率细节、远处取低分辨率概要，自然支持多尺度渲染；（3）冗余控制，锚点的增删策略比逐高斯的 ADC 更稳定，从结构层面抑制过重建。锚点初始化与视图相关特征融合。Scaffold-GS 的锚点并非手动指定，而是从 SfM 稀疏点云自动生成。具体而言，将 SfM 点云体素化（体素大小为 $v_s$），取每个非空体素的中心作为锚点位置。每个锚点配备一个 32 维局部上下文特征 $\mathbf{f}_v$、一个 3 维缩放因子 $\mathbf{l}_v$ 和 $k$ 个可学习偏移量 $\{\mathbf{O}_i\}_{i=0}^{k-1} \in \mathbb{R}^{k \times 3}$（默认 $k=10$）。为了赋予锚点多尺度感知能力，$\mathbf{f}_v$ 被扩展为一个三级特征库 $\{\mathbf{f}_v, \mathbf{f}_{v/2}, \mathbf{f}_{v/4}\}$，分别对应原始分辨率和 $2\times$、$4\times$ 下采样。渲染时，根据相机 $\mathbf{x}_c$ 与锚点 $\mathbf{x}_v$ 之间的相对距离 $d_{vc}$ 和视角方向 $\boldsymbol{\delta}_{vc}$，通过一个小型权重网络 $F_w$ 预测融合权重，动态加权求和生成视图相关的锚点特征： 其中 $d_{vc} = \|\mathbf{x}_v - \mathbf{x}_c\|_2$，$\boldsymbol{\delta}_{vc} = (\mathbf{x}_v - \mathbf{x}_c) / \|\mathbf{x}_v - \mathbf{x}_c\|_2$。
\textbf{设计动机与边界：} 相对原始 3DGS 以独立高斯和图像重建误差为主的设定，Scaffold-GS 把问题落在以下可实现的改动上：Scaffold-GS的核心方法包含四个部分。第一，锚点初始化：从COLMAP获取的SfM点云出发，构建规则稀疏网格的锚点，每个锚点关联一个可学习的特征向量。第二，神经高斯预测：每个锚点根据视角信息（观察方向和距离）生成一组神经高斯，通过MLP从锚点特征和视角信息动态预测高斯的属性（不透明度、颜色、旋转、缩放），而非像原始3D-GS那样直接优化每个高斯的固定属性。第三，锚点生长与剪枝：基于神经高斯的重要性（覆盖度和贡献度）开发生长和剪枝策略，可靠地提升场景覆盖。不重要的锚点被剪枝，覆盖不足的区域添加新锚点。第四，推理时只处理视锥体内的锚点，根据不透明度过滤不重要的神经高斯，保持实时渲染速度。为了消除原始 3DGS 在该目标任务中的具体失配，方法把上述表示、监督或渲染约束作为可学习和可执行的改动，而不是单独堆叠模块。提出Scaffold-GS，使用锚点构建层次化区域感知的3D场景表示。锚点从SfM点初始化形成稀疏网格，每个锚点绑定一组带可学习偏移的神经高斯，其属性（不透明度、颜色、旋转、缩放）根据锚点特征和观察位置动态预测。有效减少冗余高斯，增强对视角变化和细节层次变化的适应能力。适用时还应保留论文给出的边界：锚点网格分辨率需要手动设置；MLP预测引入额外计算开销；对极大规模场景的扩展性待验证。
\textbf{方法信息：} 代码：\url{https://github.com/city-super/Scaffold-GS}；GitHub 关注度高，约 1260 stars；引用很多，Semantic Scholar 约 841 次。



% METHOD: 2404.19149
\paragraph{SAGS（结构感知高斯生成）（2024.04）}
SAGS~\cite{2404_19149} 针对的是稀疏或资源受限重建中“高斯长得太多却结构不稳”的问题。公开 README 显示它是 Huawei Noah's Ark Lab 与高校合作的官方实现，代码目录保留 gaussian renderer、train、render 等完整管线。

它的改变有三层。第一层是 curvature-aware densification：先在 COLMAP 点云上用局部 PCA 估计曲率，对低曲率但点稀少的区域做 kNN 中点插值，补足弱纹理表面和大平面上的初始化空洞。第二层是 local-global graph encoder：局部图由 k 近邻连接点云，全局图由下采样点提供场景级信息，图网络在两级结构上传播特征，使相邻高斯共享几何和外观上下文，而不是各自独立漂移。第三层是 refinement layer：图编码后再预测高斯属性的残差偏移，保证结构约束之外仍有足够表达能力。轻量版 SAGS-Lite 去掉完整图网络，只用中点插值和较少属性存储取得紧凑模型。相对原始 3DGS 的梯度阈值 densification，SAGS 让高斯增长受几何结构约束；相对 Scaffold-GS，它更强调结构感知的高效重建而不是压缩编码。
\textbf{实现细节：} SAGS包含三个核心组件：(1)曲率感知密集化：利用局部PCA估计点云高斯曲率，对低曲率区域（弱纹理）进行中点插值上采样，增强COLMAP初始化的不足。(2)结构感知编码器：构建局部-全局图表示，其中局部图基于k近邻建立，全局图通过下采样获取。通过图卷积学习每个高斯的特征嵌入，强制邻近高斯共享属性，保持场景拓扑。(3)细化层：在图编码后添加细化层预测高斯属性的残差偏移。轻量版SAGS使用中点插值方案替代完整图网络，在保持渲染质量的同时大幅减少存储。 提出基于局部-全局图表示的结构感知GS方法，隐式编码场景几何；引入曲率感知密集化策略增强弱纹理区域；提出轻量版SAGS使用中点插值方案，实现最高24倍尺寸缩减无需压缩策略；有效减轻浮动伪影和不规则扭曲，获得更精确的深度图。 图构建增加计算开销；曲率估计依赖初始点云质量；大规模场景的图处理效率有待提升。
\textbf{实现补充：} 实现上还要核对输入、输出、参数、特征、初始化、训练损失、正则、渲染或推理接口，以及优化中的剪枝、分裂、采样、编码或解码步骤；这些抓手决定该方法到底改了 3DGS 管线的哪一环。
\textbf{方法信息：} Huawei 参与了这篇论文；代码：\url{https://github.com/eververas/sags}。



% METHOD: 2503.17486
\paragraph{ProtoGS（高斯原型共享）（2025.03）}
ProtoGS~\cite{2503_17486} 观察到邻近区域，尤其是天空、墙面、地面等均匀区域中的高斯属性常高度相似，因此不必让每个高斯都独立保存完整参数。它的基线是原始 3DGS 和 Scaffold-GS 一类压缩表示：原始 3DGS 渲染快但点多，Scaffold-GS 通过 anchor 和 MLP 生成局部高斯但会引入神经解码开销。ProtoGS 的改变是学习 Gaussian prototypes，用一个原型代表同一局部组内多个相似高斯，渲染时直接使用原型化后的显式基元，而不是通过 MLP 临时生成属性。

具体做法不是把空间均匀切块，因为简单切块可能把同一表面或纹理区域切到不同组里，导致原型破坏几何连续性。ProtoGS 从 SfM 点云中随机选取一组 anchor，把每个高斯分配给最近 anchor 形成 tile；这样同一 tile 中的高斯在空间上接近，也更可能共享结构。随后在每个 tile 内做迭代 K-means，推导若干 prototype；每个原型承载该组内多个高斯的共享位置、形状和外观特征。关键是聚类不是离线固定后结束，ProtoGS 直接用原型参与渲染，并用图像重建损失反传来同时优化原型和 anchor，使聚类结果受最终渲染质量约束。

相对 LightGaussian 的“删低贡献点”，ProtoGS 更接近“把多个相似点合并成可渲染代表”；相对 GHAP 的最优传输合并，它没有显式求解混合分布的传输距离，而是用 SfM anchoring 加 K-means 得到工程上更直接的原型组；相对 Scaffold-GS，它避免 MLP 解码，保留显式渲染效率。论文报告在 Mip-NeRF 360、Tanks and Temples、Deep Blending、DTU 等场景中减少高斯数并提高渲染速度，均匀区域收益尤其明显。
\textbf{实现细节：} 与 GaussianForest 通过层次结构共享隐式属性不同，ProtoGS（2025.03） 从聚类角度解决冗余问题。其核心动机源于一个简单观察：场景中相邻区域的同类高斯基元（尤其是纹理均匀区域）往往高度相似，却各自独立存储 59 维参数，造成大量冗余。与其像 GaussianForest 那样设计树形结构让叶节点沿路径共享特征，不如直接将相似基元"合并"为少数代表性原型。每个原型本质上就是一组相似高斯的质心向量，直接承载完整的 3DGS 参数（位置、协方差、颜色、不透明度），而非需要 MLP 解码的隐式特征。具体而言，ProtoGS 采用三步流水线。第一步是 SfM 锚点分片：从 SfM 点云中随机采样 $M$ 个锚点，按欧氏距离将所有高斯基元 $G \in \mathbb{R}^{N \times d}$ 分配到最近的锚点所属分片。这种基于空间位置的分组确保每组分片内的基元在几何和纹理上具有局部一致性，避免将空间上不相邻但恰好落入同一网格的基元强行聚类——这是朴素网格划分的主要缺陷。第二步是 渲染引导的 K-means 聚类：在第 $m$ 个分片内，对 $n_m$ 个基元执行 K-means 聚类，推导出 $K^m$ 个原型 $\bm{\mu}_{m,k}$。这意味着聚类分配 $r_{ik}^m$ 的更新受渲染误差的监督——如果某次聚类重分配导致渲染质量下降，梯度信号会修正原型参数和分配关系。第三步是 渐进衰减：在训练过程中以 0.5 的衰减率逐步减半原型数量，实现自适应压缩。这一设计的核心优势在于：原型本身就是标准 3DGS 参数向量，可以直接输入 3DGS 的光栅化渲染管线，无需任何 MLP 解码步骤。这使得 ProtoGS 在渲染效率上显著优于 Scaffold-GS 和 GaussianForest 等混合表示方法——后者需要为每个高斯运行 MLP 前向传播来解码协方差和颜色，引入了不可忽视的计算开销。实验中，ProtoGS 在 Mip-NeRF 360、Tanks and Temples、Deep Blending、DTU 和 Synthetic Blender 等 27 个场景上，以远少于基线方法的原型数量实现了可比甚至更优的渲染质量，存储开销大幅降低的同时保持了 3DGS 级别的高速渲染速度。
\textbf{方法信息：} Alibaba 参与了这篇论文；截至 2026-07-09，未确认官方公开代码。



% METHOD: 2505.06523
\paragraph{Virtualized 3D Gaussians（集群 LoD 系统）（2025.05）}
Virtualized 3D Gaussians~\cite{2505_06523} 借鉴游戏引擎中的虚拟化几何思想，处理由大量 3DGS 资产组合成的大场景。问题不是单个小场景文件太大，而是上千万高斯同时渲染时，必须按视角只加载和绘制需要的细节。

离线阶段，每个资产先按空间位置切成固定大小的高斯集群，例如每个最细集群约 4096 个高斯；然后把相邻集群成组，用 local splatting 在少量伪视角上蒸馏出更粗集群，每一层大致把集群数量减半，并保存父子层级 bundle。这个过程不是简单删点，而是优化粗层高斯，使其在局部视角下尽量复现细层渲染。在线阶段，渲染器先把所有候选集群投影到屏幕，计算 footprint，并按用户给定的容忍阈值选择足够细但不过度细的层级；最终只把被选集群送入原始 splatting rasterizer。相对 FLoD 面向单场景多级模型，V3DG 更像资产级 cluster LoD，适合人群、树木、建筑等组合场景。
\textbf{设计动机与边界：} 相对原始 3DGS 以独立高斯和图像重建误差为主的设定，Virtualized 3D Gaussians 把问题落在以下可实现的改动上：V3DG包含两个阶段：(1) 离线构建：每个3DGS资产先被分割为固定大小的集群(nGC=4096个高斯/集群)作为最细层级，使用二分中值分割进行聚类。然后通过局部splatting简化策略迭代构建更粗层级，每层将集群数量减半，建立父子关系的树结构，最小化不同粒度间的视觉差异；(2) 在线选择：渲染时将所有集群投影到2D屏幕空间计算足迹(footprint)，衡量每个集群的细节水平。根据预定义的足迹容忍度选择合适层级的集群，仅光栅化必要的高斯，实现与渲染条件匹配的细节水平和抗锯齿效果。创建了包含合成和真实场景的数据集(物体、树木、人物、建筑)，每个需要0.1亿高斯捕捉细节。其设计目的不是单独增加一个模块，而是用上述表示、监督或渲染约束消除原流程在目标任务中的具体失配。提出V3DG，受UE5 Nanite启发的集群LOD系统，通过离线构建层次化高斯集群和在线足迹评估选择，实现包含0.1亿高斯的大规模组合场景实时渲染。适用时还应保留论文给出的边界：离线构建过程耗时；足迹容忍度需手动调节；对快速移动相机可能产生LOD切换跳变。
\textbf{方法信息：} 代码：\url{https://github.com/city-super/v3dg}；GitHub 关注度较高，约 103 stars。



% METHOD: 2510.09997
\paragraph{CLoD-GS（连续距离衰减 LoD）（2025.10）}
CLoD-GS~\cite{2510_09997} 解决离散 LoD 切换的跳变问题。许多 LoD 方法在距离阈值处切换高斯子集，可能造成闪烁；CLoD-GS 的改变是给每个高斯学习随距离衰减的不透明度或可见性函数。

方法为高斯引入 distance decay 和 virtual distance scale，使远处高斯自然降低贡献或被动态 mask，而近处保留细节。训练采用 coarse-to-fine，使模型先学远距离稳定结构，再逐步加入近距离细节。相对 V3DG 的集群层级选择，CLoD-GS 的 LoD 是连续的、逐高斯的，更适合避免层级切换伪影。
具体机制接在标准 3DGS 的不透明度项上：每个高斯额外学习一个非负距离衰减因子，渲染时先计算它到当前相机的距离，并用当前视锥内高斯的平均距离做归一化；用户给定的 virtual distance scale 会放大或缩小这种距离感。归一化距离越大，衰减因子越强，该高斯的有效不透明度越低，低于阈值时就不送入 rasterizer。训练阶段随机改变虚拟距离，并加入渲染点数正则，使同一个模型既能在近距离保留细节，也能在远距离自动变稀。这样 CLoD-GS 不需要存多份离散 LoD 模型，而是在单一高斯集合中学习连续的可见性调度。
\textbf{实现细节：}
\textbf{方法信息：} Alibaba 参与了这篇论文。
\textbf{设计动机与边界：} 相对原始 3DGS 以独立高斯和图像重建误差为主的设定，CLoD-GS 把问题落在以下可实现的改动上：CLoD-GS在标准3DGS基础上引入三个创新：（1）距离自适应不透明度：为每个高斯基元i引入可学习距离衰减因子γ\_{d,i}。渲染时计算归一化距离d'\_i = d\_i / mean(d\_j)，结合用户可控虚拟距离缩放因子，使用指数衰减函数动态调整有效不透明度：α'\_i = α\_i · exp(-γ\_{d,i} · max(0, d'\_i - 1))。距离远的高斯自然淡出，避免离散移除的视觉跳跃。（2）动态掩码机制：基于视距自动屏蔽贡献极小的高斯，提高渲染效率。（3）训练策略：虚拟距离缩放机制和从粗到精训练策略配合渲染点数正则化，确保模型在所有距离下的鲁棒性。整个框架使用单一统一模型，无需存储多个LoD版本。其设计目的不是单独增加一个模块，而是用上述表示、监督或渲染约束消除原流程在目标任务中的具体失配。提出CLoD-GS框架，为每个高斯基元引入可学习的距离衰减参数，基于视点距离动态调整不透明度实现连续LoD，结合虚拟距离缩放和从粗到精训练策略，在单一模型中实现平滑质量缩放，消除离散LoD的存储开销和视觉跳跃伪影。适用时还应保留论文给出的边界：距离衰减参数增加了每高斯的存储开销；对极端视角变化的适应性有限；训练中需要额外的虚拟距离缩放增加训练时间。

\section{训练加速与高效表示}

训练加速路线不一定改变最终文件格式，它问的是“哪些计算没有必要发生”。原始 3DGS 训练初期高斯少但常用高分辨率渲染，后期高斯多但许多已经收敛，densification 还可能无预算地增长。EDGS、Taming 3DGS、Trick-GS、GSta、DashGaussian 和 FLoD 分别从初始化、预算控制、技巧组合、冻结收敛高斯、分辨率调度和多级表示处理这一问题。
% METHOD: 2406.15643
\paragraph{Taming 3DGS（预算受控 densification）（2024.06）}
Taming 3DGS~\cite{2406_15643} 是训练预算控制路线的代表性锚点。原始 3DGS 的最终高斯数不是输入，而是 densification 过程的副产品：同样从 SfM 点云开始，不同场景可能长出相差一个数量级的高斯，训练时间和显存峰值也随之波动。Taming 的改变，是把 primitive count 从“训练后才知道的结果”改成“训练前给定的预算”，并且尽量在这个预算内把新增高斯分给真正缺细节的位置。

它的 densification 是纯构建式的：不依赖先大量增长再剪枝，而是在每个致密化阶段先根据预设增长曲线计算还允许新增多少高斯。论文观察到原始 3DGS 的高斯增长量通常随训练推进而减小，因此用从 SfM 初始点数走向目标预算 $B$ 的调度曲线模拟这种增长趋势；若低不透明度 pruning 使当前数量低于目标轨迹，下一轮就补足差额。这样峰值数量和最终数量都可控，不会像增长-剪枝循环那样先产生很大的中间模型。

新增谁由 score-based densification 决定。系统在若干训练视角上收集候选高斯的误差相关信号、覆盖信息和不透明度等统计量，对候选排序，然后只在排名靠前的高斯附近 split/clone；高不透明度高斯可在后期打开，以增强少量基元的表达能力。公开代码的训练入口保留了 \texttt{--budget}、\texttt{--mode}、\texttt{--cams}、\texttt{--ho\_iteration} 和 \texttt{--sh\_lower} 等参数，对应预算模式、评分用相机数、高不透明度启用时刻和低频 SH 更新。实现层面，Taming 还把若干梯度计算和属性更新改成数值等价或近似等价的更快形式，并调整反向传播并行化，使“少长点”和“少算不必要更新”同时发生。

相对 LightGaussian 的训练后剪枝，Taming 更早介入：它让模型一开始就不长出大量冗余高斯，因此训练时间、显存和最终大小都会下降。相对 SteepGS，Taming 不试图证明哪一个高斯处在鞍点，而是给工程部署提供可预测预算控制。论文报告在预算设置下模型大小和训练时间约缩减 4--5$\times$，质量仍与原始 3DGS 有竞争力；预算放宽时，它也可以接近或超过原始 3DGS 的质量。
\textbf{锚点细化：} Taming 3DGS：~针对~Gaussians on a Diet~增长-剪枝循环每轮剪枝移除的高斯可能恰好在下一轮增长中被重新生成、造成冗余振荡的局限，用\textbf{引导性的纯粹建设性增密策略}替代增长-剪枝循环，核心思想是\textbf{从不执行任何剪枝，模型大小以受控方式单调增长直至精确命中预算}——利用~3DGS~增密的自然减速特性，用一条从~SfM~初始点数出发精确终止于预算~$B$~的抛物线调度曲线控制每步新增数量，配合八指标分数引导的可控增密（位置梯度、像素覆盖、显著性等）和低频增密（降至原始~1/5~给优化留出自我纠错时间），以~4--5$\times$~更小模型和~4--5$\times$~更快训练达到竞争质量。
\textbf{实现细节：} 与 Gaussians on a Diet 的迭代增长-剪枝循环不同，Taming 3DGS（2024.06） 用引导性的、纯粹建设性的增密策略替代原始 ADC，在训练全过程中只增不减——从不执行任何剪枝操作，模型大小以受控方式单调增长直至精确命中用户预设的预算上限。从增长-剪枝到纯粹建设性：为什么舍弃循环。Gaussians on a Diet 的增长-剪枝循环允许频繁试错，但也引入了不稳定性：每轮剪枝移除的高斯可能恰好在下一轮增长中被重新生成，造成冗余的振荡。Taming 3DGS 基于对 3DGS 增密行为的实证分析发现，原始方法在 Mip-NeRF 360 室外场景中，每步新增高斯数量呈现二次递减趋势——即增密本身具有自然的减速特性。利用这一规律，Taming 3DGS 用一条从 SfM 初始点数出发、精确终止于预算 $B$ 的抛物线调度曲线 $A(x) = -\frac{S-2N+N^{}2x^{}2}{N^{}2} + 2x + B$（$N$ 为增密步数、$S$ 为初始点数）控制每步新增高斯数量，使模型增长过程完全可预测且无峰值溢出。由于不执行剪枝，框架通过累积差值补偿自然剪枝（低不透明度高斯的逐渐消失）：每步计算当前数量与累积目标之间的差额，补增相应数量的新高斯，确保最终数量精确命中预算。分数引导的可控增密。原始 ADC 仅依赖位置梯度阈值做增密决策，无法区分"真正需要更多采样"与"高斯本身被错误放置"两种情况，容易导致重复复制错误高斯。Taming 3DGS 提出一种多指标评分函数 $F$，为每个高斯计算综合得分 $S_g$ 作为增密采样权重。
\textbf{设计动机与边界：} 相对原始 3DGS 以独立高斯和图像重建误差为主的设定，Taming 3DGS 把问题落在以下可实现的改动上：Taming 3DGS系统性地解决了3DGS资源消耗问题：1）可预测模型增长：分析了原始3DGS在MipNeRF360上的致密化行为，发现模型大小难以预测且波动剧烈。提出替代原始启发式致密化的方法：使用基于分数的排名引导致密化，每次致密化阶段添加的新高斯数量有上限，用户可完全控制最终模型大小。模仿原始3DGS的增长曲线但峰值和最终高斯数量完全可控。2）高质量每高斯：通过分数排名确保新增高斯贡献最大，使用高不透明度高斯增加表达力。3）训练加速：对原始3DGS管线进行仔细分析，推导梯度计算和属性更新的更快数值等价解，包括高效反向传播的替代并行化方案。在适当的地方提出保持质量的近似以进一步减少训练时间。4）整体方案：优化具有高可控性、灵活性和性能，模型大小确定性地增长到用户指定预算。其设计目的不是单独增加一个模块，而是用上述表示、监督或渲染约束消除原流程在目标任务中的具体失配。提出引导式纯构建性致密化过程，用基于分数的高斯基元致密化精确控制模型大小向用户预算增长。推导了梯度计算和属性更新的更快数值等价解，提出保持质量的近似以进一步减少训练时间。实现4-5倍模型大小和训练时间的缩减，在预算设置下保持竞争力质量。适用时还应保留论文给出的边界：较小的模型大小可能限制对极复杂场景的重建能力。某些近似可能在特定场景中引入轻微质量损失。
\textbf{方法信息：} 代码：\url{https://github.com/humansensinglab/taming-3dgs}；GitHub 关注度高，约 352 stars；引用很多，Semantic Scholar 约 250 次。



% METHOD: 2408.12894
\paragraph{FLoD（可定制多级细节训练）（2024.08）}
FLoD~\cite{2408_12894} 在训练和部署之间建立多级模型。它的基线是普通 3DGS 和 LightGaussian/CompactGS 一类单模型压缩：压缩模型能在低端 GPU 上跑，但在高端 GPU 上无法利用更多显存换更高质量；高质量模型则可能在低端设备上跑不动。FLoD 的改变是把一个场景训练成多级 Level of Detail 表示，每一级都是能独立重建完整场景的 3DGS，而不是只提供局部补丁。

核心机制是 level-specific 3D scale constraint。第 $l$ 级训练时，每个高斯的 3D scale 受到一个下界约束：低级别的下界更大，只能用较大高斯表示粗结构；高级别逐步减小下界，允许更小高斯捕获细节。训练采用 level-by-level 过程：第 1 级从 SfM 初始化开始训练，完成后保存为该级输出，同时用其 checkpoint 初始化下一等级；下一等级继承前一级结构，再放松尺度下界并继续 densification/optimization。这样不同级别之间结构连续，不会像从头训练多个模型那样互相不一致。

FLoD 还提供两种渲染模式。单级渲染时，设备显存低就选较低 level，显存高就选最高 level。选择性渲染时，系统根据高斯到相机的距离选择级别：远处使用低级大高斯，近处使用高级小高斯，从而在感知差异较小的区域节省显存。它与 V3DG 的资产级 cluster LoD 不同，FLoD 更像单场景多级 3DGS；与 Octree-GS 这类层级结构不同，它的重点不是只追求最高质量层级，而是让低、中、高端硬件都有可用质量档位。
\textbf{方法信息：} 代码：\url{https://github.com/3DGS-FLoD/flod}；GitHub 约 29 stars。
\textbf{设计动机与边界：} 相对原始 3DGS 以独立高斯和图像重建误差为主的设定，FLoD 把问题落在以下可实现的改动上：FLoD构建L\_max级3D Gaussian表示，每级通过不同的3D尺度约束s\_min\^{}(l)控制细节水平。低级别用更少更大的Gaussian重建粗结构，高级别用更多更小Gaussian捕捉细节。关键技术：(1)3D尺度约束而非2D约束，避免因相机位置变化导致不一致训练信号；(2)逐级训练(coarse-to-fine)：每级训练完成后保存检查点并克隆作为最终输出，然后用检查点初始化下一级；(3)重叠剪枝消除Gaussian过度重叠导致的伪影；(4)选择性渲染：根据Gaussian到相机的距离选择不同级别的Gaussian，远处用低级别、近处用高级别，实现内存高效渲染。FLoD可泛化到不同3DGS框架。其设计目的不是单独增加一个模块，而是用上述表示、监督或渲染约束消除原流程在目标任务中的具体失配。提出FLoD（Flexible Level of Detail），首个真正遵循LoD核心原则的3DGS方法，通过3D尺度约束构建多级表示，每级独立重建完整场景，支持单级渲染和选择性渲染，可适配从低端到高端的各种GPU设置。适用时还应保留论文给出的边界：多级表示增加了总存储需求；训练时间随级别数线性增加；选择性渲染的距离阈值设定需要根据场景调整。
% METHOD: 2501.14534
\paragraph{Trick-GS（效率技巧组合）（2025.01）}
Trick-GS~\cite{2501_14534} 不是单一新模块，而是系统性组合多种有效策略。它的基线是原始 Gaussian Splatting，以及 CompactGS、EAGLES、LightGaussian、Mini-Splatting 等效率改进。论文的问题意识是：单个技巧往往只改善一端，例如只减少模型大小、只加快训练、或只加快渲染；面向手机和资源受限设备时，训练速度、渲染速度、磁盘大小和质量必须一起权衡。Trick-GS 因此把它定位成一组经过消融筛选的效率配方。

它包含四类策略。第一类是 progressive training：训练从较低图像分辨率开始，逐步回到原分辨率；同时使用逐渐减弱的模糊/噪声和尺度约束，让模型先拟合低频结构，再学习细节，避免一开始就用高分辨率渲染浪费计算。第二类是 volume mask pruning：为每个高斯学习硬 mask，把 mask 作用在不透明度和尺度上，并用 straight-through estimator 让二值决策可反传；这些 mask 在 densification 阶段和后续固定间隔中用于删除低收益高斯。第三类是 significance pruning：参考 LightGaussian，用射线命中频率、体积和不透明度计算显著性，进一步删掉整体影响较小的高斯。第四类是 SH band masking：对球谐波段学习 mask，删除不需要的高阶视角相关颜色，降低磁盘存储。

相对 Taming 3DGS 的预算化 densification，Trick-GS 不强调一个全局目标高斯数，而是把多种已知有效技巧组合成可调管线；相对 EAGLES 和 LightGaussian，它更关注训练过程中的渐进调度与多轴平衡。Samsung Research 的技术博客中总结其效果为训练时间约降至 $1.7\times$ 更快、存储约小 $23\times$、FPS 约提升 $2\times$；论文摘要也报告最高约 $2\times$ 训练加速、$40\times$ 磁盘压缩和 $2\times$ 渲染加速。
\textbf{方法信息：} Samsung R\&D Institute UK 参与了这篇论文；截至 2026-07-09，未确认官方公开代码。



% METHOD: 2503.18402
\paragraph{DashGaussian（渲染分辨率与基元增长调度）（2025.03）}
DashGaussian~\cite{2503_18402} 把 3DGS 优化理解为逐步拟合更高频图像成分。它的基线包括原始 3DGS、Taming 3DGS、Mip-Splatting 等可插拔 backbone；问题是 3DGS 的训练复杂度主要由两个变量决定：当前渲染分辨率和当前高斯数量。训练初期高斯稀疏，直接渲染高分辨率图像是在拟合模型还表达不了的高频；训练后期高斯数量庞大，继续高强度增长又收益有限。DashGaussian 的增量是把渲染分辨率和高斯增长作为同一个 optimization complexity 来调度。

方法先从低分辨率训练开始，让高斯集合拟合训练视图中的低频结构；随着训练推进，resolution scheduler 逐步提高渲染分辨率，使模型开始接触更高频细节。为了避免低分辨率阶段过早 densification 造成过多 primitive，DashGaussian 又同步调度高斯增长：当前分辨率较低时抑制高斯数量增长，分辨率提高后再允许更多高斯补足细节。换句话说，它不是简单把图片缩小训练，而是让“图片频率需求”和“高斯表达能力”匹配，避免低分辨率 aliasing 和过度致密化。

实现上，官方代码将 DashGaussian 作为插件接入 3DGS，并切换到更快的 backward/rasterization 分支；配置中可通过 \texttt{--dash} 打开调度，通过 \texttt{--fast} 使用 Sparse Adam，并可设置 primitive number upper bound。项目页报告 DashGaussian 在不同 backbone 上平均减少约 45.7\% 训练时间；官方仓库示例显示，在 Mip-NeRF 360 上 3DGS-Dash 可将训练时间从约 12.70 min 降到 6.21 min，结合 Sparse Adam 后可到 3.69 min。相对 Taming 3DGS 控制总预算，DashGaussian 更关注“什么时候用多少分辨率和多少高斯”；相对工程 rasterizer 加速，它是训练调度层面的复杂度分配。
\textbf{实现细节：} 核心洞察：渐进频率拟合。DashGaussian（2025.03） 从频率分析的角度重新审视 3DGS 优化。根据 Nyquist 采样定理，将高分辨率图像下采样为低分辨率图像等价于丢弃高频分量、保留低频分量。因此，3DGS 优化可以重新表述为：先拟合训练视角的低频分量（低分辨率即可监督），再逐步拟合高频分量（需要全分辨率监督）。这一重新表述直接回答了一个关键问题——低分辨率训练视图能在多大程度上替代全分辨率？答案是：只要低频分量已被充分拟合，就应切换到更高分辨率以捕获残余的高频细节。频率引导的分辨率调度。基于上述洞察，DashGaussian 提出自适应的分辨率调度策略。对于训练视图的某个分辨率 $(h, w)$，定义其频率显著性为 其中 $F^n = \mathrm{DFT}(I^n)$ 为第 $n$ 个训练视图的离散傅里叶变换（Discrete Fourier Transform, DFT）。该度量强调高频分量——高分辨率视图的 $\mathcal{X}(F)$ 包含更多高频能量，低分辨率视图的 $\mathcal{X}(F_r)$ 仅包含低频能量。直觉上，如果 $\mathcal{X}(F) - \mathcal{X}(F_r)$ 远大于 $\mathcal{X}(F_r)$，说明高频信息尚未被充分优化，应分配更多迭代给全分辨率渲染；反之则低分辨率足以完成大部分优化。据此，全分辨率渲染应分得的迭代比例为 切换分辨率的时间点为第 $s_r = S \cdot \mathcal{X}(F_r) / \mathcal{X}(F)$ 次迭代（$S$ 为总迭代数）。该策略可推广至多个中间分辨率 $r_1 < r_2 < \cdots < r_m$，其中最大下采样因子 $r_m$ 满足 $\mathcal{X}(F_{r_m}) = \frac{1}{a}\mathcal{X}(F)$（$a$ 默认为 4），其余因子在 $\mathcal{X}(F)$ 和 $\mathcal{X}(F_{r_m})$ 之间线性采样。在相邻切换点之间，分辨率通过 $r_i$ 和 $r_{i-1}$ 的逆平方插值平滑过渡。
\textbf{方法信息：} Huawei Noah's Ark Lab 参与了这篇论文；代码：\url{https://github.com/YouyuChen0207/DashGaussian}；GitHub 星标较高，约 151 stars。



% METHOD: 2504.06716
\paragraph{GSta（午休高斯与动态冻结）（2025.04）}
GSta~\cite{2504_06716} 的观察是：训练中不同高斯收敛速度不同，许多高斯很早稳定，却仍在后续迭代中反复计算梯度和更新参数。它的基线是原始 3DGS 训练：每轮渲染、反传、densification 和属性更新默认面向全部活跃高斯。GSta 的改变不是删掉高斯，而是把已经收敛的高斯放入 siesta 状态，暂时冻结位置和颜色等参数，从而减少无收益的反向传播和优化器更新。

具体判据来自梯度范数。GSta 同时看位置参数和 diffuse color 参数的梯度，只有二者都低于动态阈值时，才认为该高斯当前已足够稳定，可以冻结。阈值不是固定常数，而是随训练迭代逐步提高：早期少冻结，避免模型还没成形就停止学习；中后期随着高斯数量增长，冻结更多低收益点。冻结 map 周期性更新，也会在特定间隔重置，让部分高斯有机会重新参与优化；如果高斯在 densification/pruning 阶段受到影响，也可以被解冻或重新评估。

GSta 还加入两个配套设计。第一是基于训练图像子集 PSNR 的 early stopping：它不用测试集调参，而是在训练子集上观察 PSNR 变化，变化小于阈值时进入后续 finetune。第二是学习率调度和 opacity 激活调整，使冻结策略不破坏最终收敛。相对剪枝方法，GSta 的收益来自“少算已经稳定的参数”；相对 Trick-GS，它是正交插件，可以与 Trick-GS 的 progressive training、mask 和显著性剪枝叠加。论文报告二者结合时可达到约 $5\times$ 训练加速、$16\times$ 磁盘减小，并把峰值内存降到约一半。
\textbf{方法信息：} Samsung R\&D Institute UK 和 Samsung Electronics 参与了这篇论文；项目仓库显示代码尚未放出：\url{https://github.com/anilarmagan/SRUK-GSta}。
\textbf{设计动机与边界：} 相对原始 3DGS 以独立高斯和图像重建误差为主的设定，GSta 把问题落在以下可实现的改动上：核心观察是训练中每个高斯有不同的收敛速率，大部分高斯很早收敛但致密化阶段仍继续添加新高斯。GSta基于位置和颜色梯度范数动态识别已收敛高斯，将其冻结（进入午休状态）停止更新。冻结的高斯在后续训练中应在适当时候解冻，因为其他高斯的更新可能影响整体拟合。方法还包括：(1) 基于训练图像子集PSNR的早停机制，避免过度训练；(2) 学习率调度器适应训练动态；(3) 与其他效率改进方法的正交互补性。与Trick-GS结合时，训练速度提升5倍，磁盘大小缩小16倍，峰值内存减半。其设计目的不是单独增加一个模块，而是用上述表示、监督或渲染约束消除原流程在目标任务中的具体失配。提出GSta方法动态识别训练中已收敛的高斯基元，基于位置和颜色梯度范数将它们进入'午休'状态（冻结更新），提高训练速度同时保持准确度。引入基于训练图像子集PSNR的早停机制。结合学习率调度器等改进，GSta在收敛速度、内存和存储需求上实现改进的Pareto前沿。与Trick-GS结合可实现5倍更快训练和16倍更小磁盘占用，同时准确度可比，峰值内存仅为一半。适用时还应保留论文给出的边界：冻结策略的阈值需要调参；对于高度动态场景可能需要更灵活的冻结-解冻策略；主要在单目场景验证。
% METHOD: 2504.13204
\paragraph{EDGS（消除迭代致密化的密集初始化）（2025.04）\cite{2504_13204}}
EDGS 针对的是原始 3DGS 训练中最核心、也最耗时的一步：densification。标准 3DGS 从 SfM 稀疏点云开始，先用少量高斯拟合粗结构，再根据图像梯度和尺度阈值反复 clone/split 高斯；这条路径的缺点是优化必须等待新高斯逐步出现、重新调整颜色和尺度，很多细节区域在训练早期没有足够基元承载误差。Taming 3DGS 控制“最多长多少”，SteepGS 判断“谁值得分裂”，而 EDGS 更激进：它试图在训练开始前就用图像对应关系把足够密的高斯放到合理位置，从而去掉迭代致密化本身。

具体流程从每张参考图像的密集像素对应开始。系统先用预训练匹配/对应网络在相邻视图之间建立 dense correspondences，再把这些像素对应三角化到 3D 空间。不是每个三角化点都等价可信，EDGS 会结合语义置信度和几何一致性，在每个参考图像上构建采样分布 $p^i$：纹理、边缘和高频区域会获得更多候选，高几何一致性的对应更容易被采样。随后从这个分布中一次性采样高斯，并初始化中心、尺度、颜色和球谐项。这样，高斯不是靠光度误差一点点“长出来”，而是在优化开始时就覆盖图像中可由多视角对应支持的细节。

进入训练后，EDGS 仍使用 3DGS 的可微 splatting 和重建损失，只是去掉或大幅弱化后续 densification。可以把它理解为把“增长策略”前移到初始化阶段：原始 3DGS 用若干千次迭代发现哪里缺点，EDGS 用密集对应和三角化提前给这些位置分配基元。这个改变对训练路径影响很大，因为每个高斯从一开始就能并行接受颜色、尺度和位置优化，不需要经历新生高斯的冷启动期。官方项目页和代码 \url{https://github.com/CompVis/EDGS} 显示这是 CompVis 的 CVPR 2026 工作，仓库约 720 stars；论文报告达到原始 3DGS LPIPS 水平只需约 25\% 训练时间，并使用约 60\% splats。它应作为“密度控制从迭代增长转向密集初始化”的重要锚点，而不是只作为普通加速技巧被顺带提及。

\textbf{实现细节：} 最激进的策略是 EDGS（2025.04） ——完全消除致密化过程。其核心动机是：标准 3DGS 的坐标位移量（初始位置到收敛位置的 $L_2$ 距离）极大，说明高斯在优化过程中发生了大幅位移才到达正确位置。如果能一开始就把高斯放在接近正确的位置，优化器只需微调即可。EDGS 的实现分三步。第一步，密集对应估计：对参考图 $I^i$，用预训练匹配网络 RoMa 估计它与邻域图 $I^j$ 之间的逐像素正向扭曲场 $W^{i \to j} \in \mathbb{R}^{2 \times H \times W}$ 和对应置信度 $c^{ij} \in \mathbb{R}^{H \times W}$。第二步，三角化：将匹配的 2D 像素对 $(u_k^i, v_k^i)$、$(u_k^j, v_k^j)$ 通过投影矩阵 $P^i, P^j$ 提升到 3D。具体地，对每个匹配点 $k$，从 $4 \times 3$ 投影矩阵的各列构造线性方程组 $A\,\mathbf{g}_k^x = \mathbf{b}$，用最小二乘求解高斯的 3D 齐次坐标： \begin{equation*} \mathbf{g}\_k\^{}x := \arg\min\_{\mathbf{x}} \|A\,\mathbf{x} + \mathbf{b}\|\^{}2, \quad \mathbf{g}\_k\^{}x = [x\_k, y\_k, z\_k, 1]\^{}T. \end{equation*} 第三步，采样分布：三角化产生的 3D 点数量庞大（每对图可能产生数十万个匹配），直接全部用作初始高斯既不现实也不必要——其中包含大量误匹配。
\textbf{设计动机与边界：} 相对原始 3DGS 以独立高斯和图像重建误差为主的设定，EDGS 把问题落在以下可实现的改动上：EDGS的核心思想是从2D图像中预先提取所有可用信息，而非通过光度损失逐步添加。首先选择参考图像，通过预训练网络计算密集像素对应关系。然后通过三角化将密集像素对应提升到3D空间，结合语义置信度和几何一致性聚合相邻视图的信息，为每个参考图像构建采样分布。从该分布中采样高斯，并初始化球谐函数。这种密集初始化确保了场景中均匀的细节覆盖，特别是在其他方法难以处理的高频区域。由于所有高斯在优化开始时并行初始化，消除了等待密集化步骤调整新高斯的需求。其设计目的不是单独增加一个模块，而是用上述表示、监督或渲染约束消除原流程在目标任务中的具体失配。提出从根本上消除3DGS密集化过程的方法，通过密集图像对应的三角化像素实现一步几何近似，为每个高斯提供初始化良好的颜色、尺度和位置，大幅缩短优化路径，在保持质量的同时显著加快训练速度，并可与其他加速技术兼容。适用时还应保留论文给出的边界：依赖密集对应网络的精度；对无纹理区域的对应估计可能不准确；内存开销可能因密集初始化而增加。
\textbf{方法信息：} 代码已公开，仓库为 compvis/edgs (720 stars)；GitHub 关注度很高，compvis/edgs 约 720 stars。入选依据：企业:否；代码:是；星标:是；引用:否。


\section{视频与流式渲染}

流式路线把 3DGS 看成网络传输对象。静态场景只需下载一次，动态自由视点视频则需要持续传输高斯帧；XR 场景还要求低延迟和渐进显示。LapisGS 做层级渐进 3DGS，V$^3$ 把动态高斯映射成 2D 视频，GaussianVideo 则更偏视频表示本身，用连续运动和层级学习降低动态表示成本。


% METHOD: 2408.14823
\paragraph{LapisGS（分层渐进流式 3DGS）（2024.08）}
LapisGS~\cite{2408_14823} 处理带宽受限环境下的渐进显示。它的基线是普通 3DGS 和单模型压缩：普通 3DGS 要一次性传完整场景，用户在弱网下等待时间长；单模型压缩只能在某个固定质量点上工作，不适合网络带宽动态变化。LapisGS 的改变是把 3DGS 做成“基础层 + 增强层”的累计表示，让用户先收到可渲染的粗场景，再随带宽逐步接收更高质量细节。

训练从最低分辨率图像金字塔开始，得到基础层 $G_0$；随后逐步提高训练图像分辨率，每一级只新增一个增强层 $\Delta G_i$，并把已有层作为 foundation。训练第 $i$ 层时，前面各层的位置、尺度、颜色等结构参数基本固定，只允许不透明度调整，使旧层能与新层平滑组合；新增层则完整优化，用来补充更高频细节。这样任意前缀 $G_0+\Delta G_1+\cdots+\Delta G_i$ 都是可渲染的质量层，而不是必须等全部层下载完。

为了减小冗余，LapisGS 用 occupancy map 统计每个 splat 对渲染的实际贡献，移除透明或长期无效的 splat；为了层间切换平滑，它在相邻质量层之间做 opacity interpolation。代码中还提供 dynamic-lapis-gs：首帧按静态 LapisGS 训练多层，后续帧继承前一帧各层高斯，只优化位置和旋转并保持颜色/尺度等属性稳定，用物理约束维护时间一致性。相对 FLoD 的硬件质量档，LapisGS 更面向网络传输；相对 V$^3$ 的 2D 视频化传输，它仍保持 3DGS 层结构本身。
\textbf{实现细节：} 上述 L3GS 采用"基础层+增强层"的叠加结构，通过冻结式增量训练逐层构建分层场景，再由调度器按视口优先级分批下载基元。这一方案有效地将 3DGS 场景组织为可渐进传输的层次结构，但其训练过程中各层参数完全冻结——包括不透明度在内，这意味着低层基元对渲染的贡献权重无法随高层内容的引入而动态调整。当高层基元捕获了更丰富的细节后，低层中那些仅表征粗略特征的基元仍然以原始不透明度参与渲染，既造成视觉冗余（低频基元与高频基元叠加产生伪影），又导致模型体积膨胀（需要更多基元来补偿未优化的低层贡献）。LapisGS（2024.08） 正是在这一背景下提出的改进方法，它继承了 L3GS 的分层渐进训练框架，但引入了三项关键设计来克服上述局限。与 L3GS 的核心差异：累积表示的分层结构。L3GS 将排序后的 $d$ 个基元均匀切分为 $L$ 层，每层包含独立的一组基元，层与层之间互不共享参数。LapisGS 则采用不同的构造路径：它从多尺度图像金字塔出发（将全分辨率图像集 $D_N$ 逐级下采样为 $D_0, D_1, \ldots, D_N$），从最低分辨率开始逐级训练。形式化地，分层模型表示为 $\{G_0, \Delta G_1, \Delta G_2, \ldots, \Delta G_N\}$，其中 $G_0$ 为基础层，$\Delta G_k$ 为第 $k$ 个增强层。第 $i$ 级的完整表示为 $G_i = G_0 + \sum_{k=1}^{i} \Delta G_k$——即所有低层基元的累积叠加，而非 L3GS 中各层独立的基元集合。这种累积结构使得不同分辨率层之间可以共享视觉信息：低层基元捕获场景的粗略布局和低频特征，高层基元则专注于补充高频细节，避免了跨层的信息冗余。动态不透明度优化。这是 LapisGS 与 L3GS 最本质的设计差异。在 L3GS 的冻结式增量训练中，训练第 $l$ 层时完全冻结前 $l-1$ 层的所有参数（包括位置、协方差、球谐系数和不透明度）。LapisGS 则选择了一条中间路径：训练第 $i$ 级增强层 $\Delta G_i$ 时，冻结前 $i-1$ 级基元的所有参数\textit{除了不透明度}。
\textbf{方法信息：} 代码：\url{https://github.com/nus-vv-streams/lapis-gs}；GitHub 约 51 stars。
\textbf{设计动机与边界：} 相对原始 3DGS 以独立高斯和图像重建误差为主的设定，LapisGS 把问题落在以下可实现的改动上：LapisGS构建N+1级质量的分层3DGS模型。核心方法：(1)分层渐进表示：从最低分辨率图像集开始优化基础层G\_0，逐步增加分辨率训练增强层DeltaG\_i，每层在上一层级基础上构建，低频特征在早期学习，高频特征在后期补充；(2)多级优化：训练每级时固定前级参数仅优化不透明度值，新增强层的Gaussian完全可优化；(3)不透明度插值：在层级间平滑过渡，通过插值相邻层的不透明度值实现；(4)占用图：追踪每个splat的贡献度，移除透明splat减小模型体积；(5)流媒体渲染时根据带宽选择适当的层级组合。其设计目的不是单独增加一个模块，而是用上述表示、监督或渲染约束消除原流程在目标任务中的具体失配。提出LapisGS，首个支持自适应流媒体和渐进渲染的分层3DGS，通过多分辨率训练创建累积层级表示，动态不透明度优化维持视觉保真度，占用图高效管理Gaussian。适用时还应保留论文给出的边界：多级训练增加总训练时间；层级间不透明度优化可能引入微小视觉不一致；对极大规模场景的可扩展性待验证。
% METHOD: 2409.13648
\paragraph{V$^3$（把动态高斯当作 2D 视频传输）（2024.09）}
V$^3$~\cite{2409_13648} 的锚点是硬件视频 codec。动态 3DGS 质量高，但每帧高斯属性都要传输和解码，移动端很难承担；NeRF/VideoRF 类方法可流式传输但细节偏软。V$^3$ 的改变是把动态 3DGS 的属性序列重新排列成 2D Gaussian video，让位置、旋转、尺度、不透明度、颜色和球谐等属性变成视频像素通道，从而直接利用 H.264/H.265 等成熟硬件编解码器。

训练采用分组和两阶段策略。对于一个长序列，V$^3$ 先按帧组处理拓扑变化；每个组的首帧用静态 3DGS 优化并剪枝，得到适合流式传输的紧凑高斯集合。后续帧的第一阶段用 hash encoding 和浅 MLP 从上一帧估计运动，把上一帧高斯 warp 到当前帧；第二阶段不再重新训练全部结构，而是微调高斯属性，并加入 residual entropy loss 和 temporal loss，使属性序列更容易被视频 codec 压缩，也减少帧间闪烁。这样运动和外观被解耦：运动先负责把点带到正确位置，外观微调再修复衣服纹理、头发、表情等细节。

播放器端先用平台硬件解码 2D 高斯视频，再把像素中的属性恢复为 3D 高斯并用 shader/SIBR/iOS 等播放器实时渲染。官方仓库还分出主训练代码、SIBR viewer 和 iOS viewer。相对 D-FCGS 的专用高斯运动补偿与熵模型，V$^3$ 的工程优势是拥抱现有视频生态；相对 LapisGS，它处理的是动态体视频并把数据布局改成 codec 友好的 2D 网格。
\textbf{方法信息：} 代码：\url{https://github.com/AuthorityWang/VideoGS}；GitHub 星标较高，约 107 stars；引用较多，Semantic Scholar 约 52 次。
\textbf{设计动机与边界：} 相对原始 3DGS 以独立高斯和图像重建误差为主的设定，V\^{}3 把问题落在以下可实现的改动上：核心创新：(1) 动态3DGS转2D视频：将每帧的高斯属性（位置、颜色等）排列为2D图像格式，利用标准硬件视频编解码器（H.264/H.265）压缩传输；(2) 两阶段训练：阶段一用哈希编码+浅MLP学习运动，通过剪枝减少高斯数量满足流式需求；阶段二用残差熵损失和时间损失微调其他高斯属性，提升时间连续性；(3) 运动与外观解耦：分离运动和外观表示，保持高渲染质量同时存储紧凑；(4) 多平台播放器：解码2D高斯视频并渲染，支持平滑滚动和即时分享。其设计目的不是单独增加一个模块，而是用上述表示、监督或渲染约束消除原流程在目标任务中的具体失配。首次实现动态3DGS在移动设备上的流式传输和渲染，将动态高斯视为2D视频利用硬件编解码器，两阶段训练实现紧凑存储，附带多平台播放器适用时还应保留论文给出的边界：长序列存储仍有压力；对网络带宽有要求；动态场景复杂度限制
% METHOD: 2501.04782
\paragraph{GaussianVideo（层级高斯视频表示）（2025.01）}
GaussianVideo~\cite{2501_04782} 面向动态视频场景的高效表示。它不是传统自由视点视频的逐帧 3DGS，也不是把每一帧当成独立图像拟合；它继承 3DGS 的显式高斯渲染优势，但把视频看成一个连续的时空对象，用连续运动模型来减少逐帧独立高斯的冗余。它的基线包括 GaussianImage、NeRV/HNeRV 和 Splatter-a-Video：前两类难以保持高频细节或时序一致性，Splatter-a-Video 依赖光流、深度、语义等额外监督；GaussianVideo 希望让高斯的语义运动从表示本身中自然出现。

方法用三次 B-spline 表示每个高斯位置随时间变化的轨迹：一个高斯不再只有固定中心，而是有一组控制点，任意时间 $t$ 的位置由 B-spline 基函数插值得到。协方差、颜色和不透明度没有同等自由度：论文发现如果让 opacity 或 scale 强烈随时间变化，优化会让大量高斯变得无效；如果让颜色由 MLP 随时间变化，模型容易“改颜色”而不是“跟着物体运动”。因此 GaussianVideo 让位置承担主要动态，用较低阶多项式或常量描述其他属性。相机运动也不是预先给定，而是用 Neural ODE 学习连续相机轨迹，使单目或弱相机信息视频也能进入统一优化。

训练采用时空层级策略。空间上，先用较粗高斯集合拟合主体结构，再增加高斯补高频细节；时间上，先学习低复杂度运动，再逐步提高时间表达能力。训练中还会把完全没有被任何帧渲染到的高斯 warp 到高误差区域，避免参数长期闲置。相对 V$^3$ 更强调移动端视频 codec，GaussianVideo 更像一种连续神经视频表示：它支持帧插值、任意空间重采样和视频风格化，优势是运动语义和时间一致性；代价是仍需为每个视频训练显式时空模型。
\textbf{方法信息：} Adobe Research 参与了这篇论文；项目页的代码链接截至 2026-07-09 仍为占位链接，未确认官方完整代码公开。
\textbf{设计动机与边界：} 相对原始 3DGS 以独立高斯和图像重建误差为主的设定，GaussianVideo 把问题落在以下可实现的改动上：GaussianVideo用3D高斯表示场景，引入B样条运动表示建模场景元素平滑运动轨迹。对于高斯位置使用三次B样条（p=3）建模时间依赖函数，每个高斯有N个控制点。对于协方差、颜色和不透明度使用较低表达力的模型（低阶多项式），因为重建损失倾向于推动高斯调整位置而非外观变化。引入Neural ODE学习连续相机轨迹。时空分层学习策略渐进精化空间和时间特征：空间上从粗到细分层训练，时间上逐步增加时间复杂度。这种分层方法实现高效训练和内存友好的优化。其设计目的不是单独增加一个模块，而是用上述表示、监督或渲染约束消除原流程在目标任务中的具体失配。提出结合3DGS与连续相机运动建模的神经视频表示方法。引入基于B样条的运动表示和时空分层学习策略。Neural ODE学习平滑相机轨迹，分层策略渐进精化时空特征，实现高质量渲染和快速收敛。适用时还应保留论文给出的边界：B样条控制点数量影响内存占用；长视频序列的建模能力受限；对复杂拓扑变化的适应性有限。

\section{其他压缩表示}

少数方法不直接压缩标准 3DGS，而是借 3DGS 的问题引出替代表示或跨域编码。它们不一定会成为主线 3DGS codec，但能提示压缩的边界：如果目标是极简传输，可以换成 triangle soup；如果目标是超高分辨率图像压缩，可以把 Gaussian sampling 放到 2D；如果目标是视角相关外观表达，也可能尝试量子编码式函数。
% METHOD: 2512.20377
\paragraph{SmartSplat（面向超高分辨率图像的高斯采样）（2025.12）}
SmartSplat~\cite{2512_20377} 属于 2D 图像压缩旁支，不是核心 3D 场景压缩，但它对本章有参考价值：它把 Gaussian splat 的“可变位置、尺度和覆盖范围”用于超高分辨率图像压缩。基线是 GaussianImage、ImageGS、LIG 等 2D Gaussian image representation；这些方法在 2K 左右图像上可行，但到 8K/16K 时，如果给定极高压缩比，可用高斯数量很少，随机或均匀初始化会让边缘、纹理和颜色过渡区域表达不足。

SmartSplat 的核心是初始化和采样。第一，Gradient-Color Guided Variational Sampling 用图像梯度和颜色方差作为结构特征，在边缘、细纹理和颜色变化大的区域分配更多高斯，在平坦区域少分配。第二，Exclusion-based Uniform Sampling 在采样时排除已经被覆盖的邻域，避免很多高斯挤在同一局部，保证有限 primitive 对整个图像有非重叠覆盖。第三，Scale-Adaptive Gaussian Color Sampling 根据高斯尺度选择颜色初始化方式：小尺度高斯更接近局部像素颜色，大尺度高斯则需要聚合更大邻域的颜色统计。

训练时，位置、尺度、旋转和颜色一起优化；实现上用 tile-based sampling 和 CUDA KNN 处理 16K 图像初始化，避免全图级内存爆炸。论文报告在 DIV8K 和自建 DIV16K 上，SmartSplat 在同压缩比下超过 3DGS/GaussianImage/LIG 等 2D 高斯基线，并能在高压缩比下保持稳定收敛。它没有多视角几何、相机位姿或可见性排序问题，但说明 Gaussian primitive 也可以作为“按内容分布的二维采样单元”，这对理解 3DGS 压缩里的采样、覆盖和属性初始化仍有借鉴意义。
\textbf{方法信息：} 代码：\url{https://github.com/lif314/SmartSplat}；GitHub 约 19 stars。
\textbf{设计动机与边界：} 相对原始 3DGS 以独立高斯和图像重建误差为主的设定，SmartSplat 把问题落在以下可实现的改动上：SmartSplat的核心创新包括三个关键策略：(1)梯度-色彩引导变分采样(Gradient-Color Guided Variational Sampling)，利用图像梯度和色彩方差作为结构感知特征，引导高斯原语在像素空间中的非重叠覆盖优化，在结构复杂区域分配更多高斯；(2)排除式均匀采样(Exclusion-based Uniform Sampling)，确保高斯不重叠覆盖图像空间，通过排除已采样区域实现均匀分布；(3)尺度自适应高斯颜色采样(Scale-Adaptive Gaussian Color Sampling)，根据高斯尺度自适应地初始化颜色参数，提升不同尺度下的颜色精度。三个策略联合优化空间布局、尺度和颜色初始化，使用有限数量的高斯高效捕获局部结构和全局纹理。支持DIV8K和新构建的16K数据集。其设计目的不是单独增加一个模块，而是用上述表示、监督或渲染约束消除原流程在目标任务中的具体失配。提出SmartSplat，一种特征感知的自适应2D高斯图像压缩框架，通过梯度-色彩引导变分采样、排除式均匀采样和尺度自适应颜色采样策略，支持任意分辨率和压缩比，在8K和16K图像上超越SOTA。适用时还应保留论文给出的边界：训练时间可能较长；对极端高压缩比下的细节保持仍有改进空间；仅在自然图像上测试。
% METHOD: 2602.05047
\paragraph{QuantumGS（量子电路式视角相关外观编码）（2026.02）}
QuantumGS~\cite{2602_05047} 的目标不是减少高斯数量，而是替代球谐或 MLP 的视角相关颜色调制。原始 3DGS 用低阶球谐表示视角相关颜色，速度快但对尖锐反射、玻璃透明和高频高光表达有限；VDGS 一类方法用经典 MLP 扩展视角相关外观，但在低参数预算下表达力仍受限。QuantumGS 的改变是把观察方向映射到 Bloch 球上的量子比特状态，再用 variational quantum circuit（VQC）输出颜色调制。

它有两个控制管线。Pipeline I 是 per-Gaussian：一个经典 hypernetwork 为每个高斯生成 VQC 的旋转角等参数，因此每个高斯都有局部自适应的量子外观函数，适合 synthetic specular object 等复杂反射场景。Pipeline II 是 global：多个高斯共享 VQC，并用 hash-conditioned 参数提供场景上下文，适合 Mip-NeRF 360、Tanks and Temples、Deep Blending 这类大场景。两者都保留原始 3DGS 的 rasterization，只替换“给定视角方向如何得到颜色/透明调制”的函数族。

从压缩角度看，QuantumGS 是旁支但不无关：许多压缩方法试图减少 SH 系数、用 MLP 或 codebook 表达颜色；QuantumGS 则探索在低参数量下用另一种函数族获得更强视角相关表达。论文在经典 GPU 上模拟量子电路，并报告 10--16 FPS 的交互速度，但瓶颈也在这里：量子电路模拟比普通 SH 慢，真实量子硬件部署仍是未来假设。因此本书把它作为“外观编码函数替代”的探索，而不是当前主流 3DGS 压缩路线。
\textbf{方法信息：} 代码：\url{https://github.com/gwilczynski95/QuantumGS}；GitHub 约 9 stars。
\textbf{设计动机与边界：} 相对原始 3DGS 以独立高斯和图像重建误差为主的设定，QuantumGS 把问题落在以下可实现的改动上：核心创新是Bloch球方向编码：将3D观看方向映射到量子比特状态，通过单量子比特旋转自然表示各向异性反射。双VQC控制策略：(1)管线I（Per-Gaussian）——超网络为每个高斯生成空间自适应VQC参数（旋转角度），最大化局部表达能力；(2)管线II（Global）——共享VQC配合哈希条件化全局参数，适用于大规模场景。VQC替代经典MLP用于视相关颜色调制（类似VDGS但使用量子电路替代）。VQC的参数化量子门序列提供比等参数量经典网络更高的希尔伯特空间表达能力。整体保持3DGS的实时渲染优势。其设计目的不是单独增加一个模块，而是用上述表示、监督或渲染约束消除原流程在目标任务中的具体失配。提出QuantumGS，首个将变分量子电路（VQC）集成到3DGS管线的混合框架。创新性地将观看方向映射到Bloch球面，利用量子比特的天然几何表示3D方向数据，通过超网络或全局条件机制生成VQC参数，在低参数量下实现更高表达能力。适用时还应保留论文给出的边界：依赖量子计算硬件/模拟器的可用性；当前仅在模拟量子环境中验证；大规模场景的可扩展性有待验证。
% METHOD: 2603.27151
\paragraph{DiffSoup（三角形汤的极简辐射场）（2026.03）}
DiffSoup~\cite{2603_27151} 并非标准 3DGS 压缩，而是提出一种极端简化的替代表示：不透明 triangle soup 加神经纹理。它的出发点是 3DGS 半透明基元需要排序和 alpha blending，传输与渲染仍有成本；三角形可以直接使用标准深度测试。

方法用一组三角形和逐三角形神经纹理函数表示颜色与不透明度：三角形顶点给几何，纹理函数在重心坐标上输出颜色与 alpha，另一个小 MLP 可结合视角方向做延迟着色。渲染时不做半透明排序，而是先产生每个像素的候选片段，再用二值 opacity 阈值和标准深度测试选择最前片段。为了让这种离散可见性可训练，DiffSoup 引入 stochastic opacity mask：前向阶段随机采样阈值决定哪些片段可见，反向阶段根据阈值附近的片段和像素损失给 opacity 与几何分配梯度；训练还采用粗到细优化，从粗特征网格过渡到细网格并重新初始化 alpha 特征。相对 Shorter Gaussian Lists，它彻底绕开高斯列表；相对 mesh 重建，它仍以辐射场训练目标优化外观。
\textbf{设计动机与边界：} 相对原始 3DGS 以独立高斯和图像重建误差为主的设定，DiffSoup 把问题落在以下可实现的改动上：DiffSoup方法包含：(1)三角形汤表示：使用三角形集合配合逐三角形神经纹理函数（颜色和不透明度），每个三角形由三个顶点定义。(2)二值不透明度的深度测试光栅化：与半透明基元不同，采用标准光栅化生成片段数组，通过不透明度阈值>0.5选择最前片段，无需3DGS中的显式排序。(3)随机不透明度掩码(SOM)：受隐式表面重建启发，通过随机采样阈值使二值不透明度决策可微。前向传播采样阈值确定可见片段，反向传播基于阈值位置分配梯度。(4)粗到细优化：从粗特征网格过渡到细网格，期间重新初始化alpha特征。(5)延迟着色：使用共享MLP处理连接颜色特征和视图方向的16维输入向量。整体流程无需mollifier，在训练中产生精确像素颜色。其设计目的不是单独增加一个模块，而是用上述表示、监督或渲染约束消除原流程在目标任务中的具体失配。提出一种极端模型简化方法，用少量不透明三角形（三角形汤）配合神经纹理和二值不透明度表示辐射场。核心创新是随机不透明度掩码(SOM)使二值不透明度可直接微分，无需平滑化光栅化(mollifier)即可稳定训练。支持标准深度测试，可无缝集成到传统图形管线，在消费级笔记本和移动设备上实现交互式渲染。适用时还应保留论文给出的边界：三角形数量限制可能影响复杂场景的细节表现，神经纹理增加了一定存储开销，粗到细优化需要调参。
\textbf{方法信息：} 代码：\url{https://github.com/kenji-tojo/diffsoup}；GitHub 关注度较高，约 69 stars。



\section{剪枝与重要性评估}

这一节收纳 2026 年后更专注“重要性估计”的方法。它们与前面剪枝、量化和压缩方法相连，但关注点更窄：如何判断一个高斯是否值得保留。原始 3DGS 主要按透明度和训练启发式调整数量，而压缩场景需要更可靠的贡献判断，最好不经过昂贵候选渲染，也能在训练中异步决定哪些高斯保留、微调或删除。


% METHOD: 2606.02346
\paragraph{VEDAL（变分误差驱动异步剪枝）（2026.06）}
VEDAL~\cite{2606_02346} 把高斯是否保留建模为潜变量，并用变分自由能优化。它针对的是同步剪枝的两个问题：一是启发式重要性分数可能过早删除还没收敛的复杂区域，二是全场景统一更新会在已经稳定的区域继续花计算。VEDAL 的改变，是让剪枝压力异步激活：只有当某个高斯的预测误差信号已经足够稳定，才让保留/删除决策进入强约束阶段。

具体流程包含两层。第一层是 prediction-error gate：系统维护每个高斯的重建误差或 leave-one-out 重要性估计的滑动统计，判断该高斯的贡献是否已经稳定；未稳定的高斯暂不承受强 KL 剪枝压力，避免慢收敛区域被提前删除。第二层是 variational uncertainty head：为每个高斯预测保留概率，把保留变量看成 Bernoulli/Concrete 风格潜变量，用 free energy 目标同时平衡重建似然和模型复杂度。训练结束后，根据保留概率剪掉低价值高斯，再保留可渲染的普通 3DGS 表示。

相对 MaskGaussian 的逐高斯概率 mask，VEDAL 的重点不是“mask 怎么反传”，而是“什么时候允许 mask 产生剪枝压力”；相对 REFINE 的后处理重要性代理，它把不确定性和复杂度惩罚放进训练过程。论文报告在 Mip-NeRF 360、Tanks and Temples 和 Deep Blending 上达到约 5.2$\times$ primitive 压缩、约 0.31 dB PSNR 下降，并保持实时渲染。需要注意的是，其官方仓库 README 目前说明核心 C++/CUDA 剪枝实现仍待完整发布，因此本书把它作为技术路线和后续代码跟踪项，而不把仓库状态等同于完整实现已经可复现。
\textbf{实现细节：} VEDAL。VEDAL（2026.06） 延续 3DGS pruning 在"剪枝与稀疏化"上的基本设定，主要改变是：VEDAL：Variational Error-Driven Asynchronous Learning，将高斯剪枝表述为变分自由能最小化，使用误差门控和不确定性头。3DGS 的显式高斯表示和可微渲染能力保持为基础。其设计动机是：把3DGS pruning从启发式重要性评分推进到信息论式变分目标，引入逐高斯预测误差门控和可学习先验的不确定性剪枝决策。具体而言，方法为高斯是否保留建模潜变量，通过free energy目标同时惩罚重建误差和复杂度；异步机制只在高误差/不确定区域触发剪枝更新。适用边界是：引入不确定性头和变分目标增加训练复杂度；超参数对压缩-质量平衡敏感。
\textbf{方法信息：} 官方仓库：\url{https://github.com/AsakaTigar/VEDAL}；核心实现仍待发布，需要持续关注。
\textbf{设计动机与边界：} 相对原始 3DGS 以独立高斯和图像重建误差为主的设定，VEDAL 把问题落在以下可实现的改动上：方法为高斯是否保留建模潜变量，通过free energy目标同时惩罚重建误差和复杂度；异步机制只在高误差/不确定区域触发剪枝更新。其设计目的不是单独增加一个模块，而是用上述表示、监督或渲染约束消除原流程在目标任务中的具体失配。把3DGS pruning从启发式重要性评分推进到信息论式变分目标，引入逐高斯预测误差门控和可学习先验的不确定性剪枝决策。适用时还应保留论文给出的边界：引入不确定性头和变分目标增加训练复杂度；超参数对压缩-质量平衡敏感。
% METHOD: 2606.09074
\paragraph{REFINE（无需渲染候选的 primitive 重要性）（2026.06）}
REFINE~\cite{2606_09074} 针对剪枝成本。许多 pruning 方法要评估“删掉某个高斯会让图像差多少”，最直接的做法是生成候选模型并前向渲染，再比较误差；这对大场景、多个压缩比例或端侧部署搜索都太慢。REFINE 的改变，是构造 rendering-free primitive importance metric：不对每个候选删除重新渲染，而是用解析近似估计删除该 primitive 的感知误差。

具体流程是先从已训练 3DGS 模型和相机集合出发，为每个高斯计算可见性、投影几何和内容自适应权重。它用解析近似的 rendering-aware Hessian field 表示删除高斯后对像素误差的局部敏感性，再把这种敏感性转成 anisotropic perceptual weight field。这个权重场同时考虑高斯在不同视角下是否可见、投影 footprint 如何覆盖图像区域，以及局部图像内容对误差是否敏感。排序时，低重要性的高斯可直接剪掉；README 中的使用方式也是输入预训练 \texttt{point\_cloud.ply} 和 \texttt{cameras.json}，由 \texttt{REFINE\_pruning.py} 输出剪枝后的点云，再交回原始 3DGS 渲染和评测脚本。

相对 PUP 的二阶不确定性，REFINE 更强调“无需候选 forward rendering”的快速代理；相对 Gaussian Herding 的局部合并，它仍是删除低重要性 primitive；相对 VEDAL 的训练内概率保留，它更像一个即插即用的后处理剪枝器。论文报告 pruning 相关计算复杂度约降低 3000$\times$，端侧延迟相对渲染式剪枝约有 20$\times$ 实际加速。
\textbf{实现细节：} REFINE。REFINE（2026.06） 延续 3DGS pruning 在"剪枝与重要性评估"上的基本设定，主要改变是：REFINE：rendering-free primitive importance，用解析近似的rendering-aware Hessian field估计删除高斯的感知误差。3DGS 的显式高斯表示和可微渲染能力保持为基础。其设计动机是：提出渲染无关但渲染感知的primitive importance metric，建模可见性、投影几何和内容自适应权重，绕开昂贵的逐候选forward rendering。具体而言，方法解析近似每个高斯移除对感知误差的影响，形成各向异性感知权重场，并据此排序/剪枝primitive。适用边界是：Hessian近似可能在复杂遮挡/视角依赖外观下失准；剪枝阈值仍需控制。
\textbf{技术实现补充：} REFINE。REFINE（2026.06） 延续 3DGS pruning 在"剪枝与重要性评估"上的基本设定，主要改变是：REFINE：rendering-free primitive importance，用解析近似的rendering-aware Hessian field估计删除高斯的感知误差。3DGS 的显式高斯表示和可微渲染能力保持为基础。其设计动机是：提出渲染无关但渲染感知的primitive importance metric，建模可见性、投影几何和内容自适应权重，绕开昂贵的逐候选forward rendering。具体而言，方法解析近似每个高斯移除对感知误差的影响，形成各向异性感知权重场，并据此排序/剪枝primitive。适用边界是：Hessian近似可能在复杂遮挡/视角依赖外观下失准；剪枝阈值仍需控制。
\textbf{方法信息：} 代码：\url{https://github.com/ZhangChen2022/REFINE}。
\textbf{设计动机与边界：} 相对原始 3DGS 以独立高斯和图像重建误差为主的设定，REFINE 把问题落在以下可实现的改动上：方法解析近似每个高斯移除对感知误差的影响，形成各向异性感知权重场，并据此排序/剪枝primitive。其设计目的不是单独增加一个模块，而是用上述表示、监督或渲染约束消除原流程在目标任务中的具体失配。提出渲染无关但渲染感知的primitive importance metric，建模可见性、投影几何和内容自适应权重，绕开昂贵的逐候选forward rendering。适用时还应保留论文给出的边界：Hessian近似可能在复杂遮挡/视角依赖外观下失准；剪枝阈值仍需控制。

\section{移动端渲染与辐射压缩}

移动端方法强调的不只是文件大小，还包括片上内存、能耗和高阶外观计算。这里的代表方法把高阶球谐辐射残差压成更紧凑的能量聚合形式，使移动 GPU 不必为每个高斯执行完整高阶外观计算。
% METHOD: 2401.02436
\paragraph{Compressed 3D Gaussian Splatting for Accelerated Novel View Synthesis（压缩、量化与加速）（2023.11）}
Compressed 3D Gaussian Splatting for Accelerated Novel View Synthesis~\cite{2401_02436} 关注的问题是：压缩或加速 3DGS 表示与渲染 Compressed 3D Gaussian Splatting for Accelerated Novel View Synthesis 是宽召回补全中发现的论文，主要根据标题、摘要和 arXiv 类别判定其与 3DGS 的关系。
相对原始 3DGS，Compressed 3D Gaussian Splatting for Accelerated Novel View Synthesis 的变化不是只替换损失函数，而是把表示、初始化、约束或网络预测路径重新组织到任务流程中：该工作被召回是因为论文文本包含 Gaussian/3D Gaussian 相关表述；当前分类先确定其任务域和与 3DGS 的相关性，后续证据流水线再补代码、企业和引用信息。
\textbf{实现细节：} 具体实现可以按输入、核心模块和渲染/优化路径来读。输入端首先取得任务所需的图像、相机、文本、时间序列、语义或几何线索，并把这些线索转换成可以约束高斯中心、尺度、旋转、不透明度、颜色或附加特征的条件。核心结构中，方法使用的主要机制是：从摘要层面看，Compressed 3D Gaussian Splatting for Accelerated Novel View Synthesis 围绕“压缩或加速 3DGS 表示与渲染”展开。若其正文或代码后续可读，应进一步核对它具体改动的是表示、初始化、渲染、优化、损失、动态建模、压缩或应用接口中的哪一环。 训练或推理时，这些模块通常先给出初始化或中间特征，再通过可微 Gaussian rasterization、重建损失、感知/语义约束、几何正则或任务损失回传到高斯参数；如果方法包含前馈网络，则网络输出的是高斯属性、变形场、压缩码、语义特征或控制变量，而不是只输出最终 2D 图像。渲染路径仍以显式高斯投影、深度排序和 alpha 合成为基础，只是在采样、属性预测、局部编辑、动态变形、压缩编码或下游任务接口处加入额外约束。这样写入本章时，它与前后方法的关系是：相对原始 3DGS 和本章已有锚点方法，它可以理解为 3DGS extends\_or\_uses 宽召回分类阶段的关系判断，后续需要按论文方法部分复核。
\textbf{证据与边界：} 本轮召回分类尚未逐项核对实验表格；后续证据补全阶段需要读取论文源码/PDF或项目页确认数据集、指标和对比方法。 该条来自宽召回批次；源码下载受 arXiv e-print 网络限制，部分论文暂时只能依据标题、摘要和类别分类，需要在入可信书前做人工/LLM精读复核。 本书只采用本地 paper pack、数据库分类、源码/引用登记和论文摘要中可核验的信息；没有本地证据支撑的速度、排行榜或开源完整性不在正文中强化。\textbf{方法信息：} 代码：\url{https://github.com/darglein/adop}；引用信号：291 citations (9.2/month, domain-relative cutoff 4.2/month, semantic\_scholar)；代码关注：keksboter/c3dgs has 399 stars (12.6/month, domain-relative cutoff 13.9/month)。
\textbf{设计动机与边界：} 相对原始 3DGS 以独立高斯和图像重建误差为主的设定，Compressed 3D Gaussian Splatting for Accelerated Novel View Synthesis 把问题落在以下可实现的改动上：从摘要层面看，Compressed 3D Gaussian Splatting for Accelerated Novel View Synthesis 围绕“压缩或加速 3DGS 表示与渲染”展开。若其正文或代码后续可读，应进一步核对它具体改动的是表示、初始化、渲染、优化、损失、动态建模、压缩或应用接口中的哪一环。其设计目的不是单独增加一个模块，而是用上述表示、监督或渲染约束消除原流程在目标任务中的具体失配。该工作被召回是因为论文文本包含 Gaussian/3D Gaussian 相关表述；当前分类先确定其任务域和与 3DGS 的相关性，后续证据流水线再补代码、企业和引用信息。适用时还应保留论文给出的边界：该条来自宽召回批次；源码下载受 arXiv e-print 网络限制，部分论文暂时只能依据标题、摘要和类别分类，需要在入可信书前做人工/LLM精读复核。
% METHOD: 2311.18159
\paragraph{CompGS（压缩、量化与加速）（2023.11）}
CompGS~\cite{2311_18159} 关注的问题是：压缩或加速 3DGS 表示与渲染 CompGS 是宽召回补全中发现的论文，主要根据标题、摘要和 arXiv 类别判定其与 3DGS 的关系。
相对原始 3DGS，CompGS 的变化不是只替换损失函数，而是把表示、初始化、约束或网络预测路径重新组织到任务流程中：该工作被召回是因为论文文本包含 Gaussian/3D Gaussian 相关表述；当前分类先确定其任务域和与 3DGS 的相关性，后续证据流水线再补代码、企业和引用信息。
\textbf{实现细节：} 具体实现可以按输入、核心模块和渲染/优化路径来读。输入端首先取得任务所需的图像、相机、文本、时间序列、语义或几何线索，并把这些线索转换成可以约束高斯中心、尺度、旋转、不透明度、颜色或附加特征的条件。核心结构中，方法使用的主要机制是：从摘要层面看，CompGS 围绕“压缩或加速 3DGS 表示与渲染”展开。若其正文或代码后续可读，应进一步核对它具体改动的是表示、初始化、渲染、优化、损失、动态建模、压缩或应用接口中的哪一环。 训练或推理时，这些模块通常先给出初始化或中间特征，再通过可微 Gaussian rasterization、重建损失、感知/语义约束、几何正则或任务损失回传到高斯参数；如果方法包含前馈网络，则网络输出的是高斯属性、变形场、压缩码、语义特征或控制变量，而不是只输出最终 2D 图像。渲染路径仍以显式高斯投影、深度排序和 alpha 合成为基础，只是在采样、属性预测、局部编辑、动态变形、压缩编码或下游任务接口处加入额外约束。这样写入本章时，它与前后方法的关系是：相对原始 3DGS 和本章已有锚点方法，它可以理解为 3DGS extends\_or\_uses 宽召回分类阶段的关系判断，后续需要按论文方法部分复核。
\textbf{证据与边界：} 本轮召回分类尚未逐项核对实验表格；后续证据补全阶段需要读取论文源码/PDF或项目页确认数据集、指标和对比方法。 该条来自宽召回批次；源码下载受 arXiv e-print 网络限制，部分论文暂时只能依据标题、摘要和类别分类，需要在入可信书前做人工/LLM精读复核。 本书只采用本地 paper pack、数据库分类、源码/引用登记和论文摘要中可核验的信息；没有本地证据支撑的速度、排行榜或开源完整性不在正文中强化。\textbf{方法信息：} 代码：\url{https://github.com/liuxiangrui/compgs}；引用信号：151 citations (4.8/month, domain-relative cutoff 4.2/month, semantic\_scholar)。
\textbf{设计动机与边界：} 相对原始 3DGS 以独立高斯和图像重建误差为主的设定，CompGS 把问题落在以下可实现的改动上：从摘要层面看，CompGS 围绕“压缩或加速 3DGS 表示与渲染”展开。若其正文或代码后续可读，应进一步核对它具体改动的是表示、初始化、渲染、优化、损失、动态建模、压缩或应用接口中的哪一环。其设计目的不是单独增加一个模块，而是用上述表示、监督或渲染约束消除原流程在目标任务中的具体失配。该工作被召回是因为论文文本包含 Gaussian/3D Gaussian 相关表述；当前分类先确定其任务域和与 3DGS 的相关性，后续证据流水线再补代码、企业和引用信息。适用时还应保留论文给出的边界：该条来自宽召回批次；源码下载受 arXiv e-print 网络限制，部分论文暂时只能依据标题、摘要和类别分类，需要在入可信书前做人工/LLM精读复核。
% METHOD: 2312.03461
\paragraph{HiFi4G（压缩、量化与加速）（2023.12）}
HiFi4G~\cite{2312_03461} 关注的问题是：压缩或加速 3DGS 表示与渲染 HiFi4G 是宽召回补全中发现的论文，主要根据标题、摘要和 arXiv 类别判定其与 3DGS 的关系。
相对原始 3DGS，HiFi4G 的变化不是只替换损失函数，而是把表示、初始化、约束或网络预测路径重新组织到任务流程中：该工作被召回是因为论文文本包含 Gaussian/3D Gaussian 相关表述；当前分类先确定其任务域和与 3DGS 的相关性，后续证据流水线再补代码、企业和引用信息。
\textbf{实现细节：} 具体实现可以按输入、核心模块和渲染/优化路径来读。输入端首先取得任务所需的图像、相机、文本、时间序列、语义或几何线索，并把这些线索转换成可以约束高斯中心、尺度、旋转、不透明度、颜色或附加特征的条件。核心结构中，方法使用的主要机制是：从摘要层面看，HiFi4G 围绕“压缩或加速 3DGS 表示与渲染”展开。若其正文或代码后续可读，应进一步核对它具体改动的是表示、初始化、渲染、优化、损失、动态建模、压缩或应用接口中的哪一环。 训练或推理时，这些模块通常先给出初始化或中间特征，再通过可微 Gaussian rasterization、重建损失、感知/语义约束、几何正则或任务损失回传到高斯参数；如果方法包含前馈网络，则网络输出的是高斯属性、变形场、压缩码、语义特征或控制变量，而不是只输出最终 2D 图像。渲染路径仍以显式高斯投影、深度排序和 alpha 合成为基础，只是在采样、属性预测、局部编辑、动态变形、压缩编码或下游任务接口处加入额外约束。这样写入本章时，它与前后方法的关系是：相对原始 3DGS 和本章已有锚点方法，它可以理解为 3DGS extends\_or\_uses 宽召回分类阶段的关系判断，后续需要按论文方法部分复核。
\textbf{证据与边界：} 本轮召回分类尚未逐项核对实验表格；后续证据补全阶段需要读取论文源码/PDF或项目页确认数据集、指标和对比方法。 该条来自宽召回批次；源码下载受 arXiv e-print 网络限制，部分论文暂时只能依据标题、摘要和类别分类，需要在入可信书前做人工/LLM精读复核。 本书只采用本地 paper pack、数据库分类、源码/引用登记和论文摘要中可核验的信息；没有本地证据支撑的速度、排行榜或开源完整性不在正文中强化。\textbf{方法信息：} ByteDance 参与；引用信号：96 citations (3.1/month, domain-relative cutoff 4.2/month, semantic\_scholar)。
\textbf{设计动机与边界：} 相对原始 3DGS 以独立高斯和图像重建误差为主的设定，HiFi4G 把问题落在以下可实现的改动上：从摘要层面看，HiFi4G 围绕“压缩或加速 3DGS 表示与渲染”展开。若其正文或代码后续可读，应进一步核对它具体改动的是表示、初始化、渲染、优化、损失、动态建模、压缩或应用接口中的哪一环。其设计目的不是单独增加一个模块，而是用上述表示、监督或渲染约束消除原流程在目标任务中的具体失配。该工作被召回是因为论文文本包含 Gaussian/3D Gaussian 相关表述；当前分类先确定其任务域和与 3DGS 的相关性，后续证据流水线再补代码、企业和引用信息。适用时还应保留论文给出的边界：该条来自宽召回批次；源码下载受 arXiv e-print 网络限制，部分论文暂时只能依据标题、摘要和类别分类，需要在入可信书前做人工/LLM精读复核。
% METHOD: 2401.04099
\paragraph{AGG（压缩、量化与加速）（2024.01）}
AGG~\cite{2401_04099} 关注的问题是：压缩或加速 3DGS 表示与渲染 AGG 是宽召回补全中发现的论文，主要根据标题、摘要和 arXiv 类别判定其与 3DGS 的关系。
相对原始 3DGS，AGG 的变化不是只替换损失函数，而是把表示、初始化、约束或网络预测路径重新组织到任务流程中：该工作被召回是因为论文文本包含 Gaussian/3D Gaussian 相关表述；当前分类先确定其任务域和与 3DGS 的相关性，后续证据流水线再补代码、企业和引用信息。
\textbf{实现细节：} 具体实现可以按输入、核心模块和渲染/优化路径来读。输入端首先取得任务所需的图像、相机、文本、时间序列、语义或几何线索，并把这些线索转换成可以约束高斯中心、尺度、旋转、不透明度、颜色或附加特征的条件。核心结构中，方法使用的主要机制是：从摘要层面看，AGG 围绕“压缩或加速 3DGS 表示与渲染”展开。若其正文或代码后续可读，应进一步核对它具体改动的是表示、初始化、渲染、优化、损失、动态建模、压缩或应用接口中的哪一环。 训练或推理时，这些模块通常先给出初始化或中间特征，再通过可微 Gaussian rasterization、重建损失、感知/语义约束、几何正则或任务损失回传到高斯参数；如果方法包含前馈网络，则网络输出的是高斯属性、变形场、压缩码、语义特征或控制变量，而不是只输出最终 2D 图像。渲染路径仍以显式高斯投影、深度排序和 alpha 合成为基础，只是在采样、属性预测、局部编辑、动态变形、压缩编码或下游任务接口处加入额外约束。这样写入本章时，它与前后方法的关系是：相对原始 3DGS 和本章已有锚点方法，它可以理解为 3DGS extends\_or\_uses 宽召回分类阶段的关系判断，后续需要按论文方法部分复核。
\textbf{证据与边界：} 本轮召回分类尚未逐项核对实验表格；后续证据补全阶段需要读取论文源码/PDF或项目页确认数据集、指标和对比方法。 该条来自宽召回批次；源码下载受 arXiv e-print 网络限制，部分论文暂时只能依据标题、摘要和类别分类，需要在入可信书前做人工/LLM精读复核。 本书只采用本地 paper pack、数据库分类、源码/引用登记和论文摘要中可核验的信息；没有本地证据支撑的速度、排行榜或开源完整性不在正文中强化。\textbf{方法信息：} NVIDIA、NVIDIA research page lists AGG under NVIDIA DAIR publication 参与；代码：\url{https://github.com/kakaobrain/coyo-dataset}；引用信号：64 citations (2.1/month, domain-relative cutoff 4.2/month, semantic\_scholar)。
\textbf{设计动机与边界：} 相对原始 3DGS 以独立高斯和图像重建误差为主的设定，AGG 把问题落在以下可实现的改动上：从摘要层面看，AGG 围绕“压缩或加速 3DGS 表示与渲染”展开。若其正文或代码后续可读，应进一步核对它具体改动的是表示、初始化、渲染、优化、损失、动态建模、压缩或应用接口中的哪一环。其设计目的不是单独增加一个模块，而是用上述表示、监督或渲染约束消除原流程在目标任务中的具体失配。该工作被召回是因为论文文本包含 Gaussian/3D Gaussian 相关表述；当前分类先确定其任务域和与 3DGS 的相关性，后续证据流水线再补代码、企业和引用信息。适用时还应保留论文给出的边界：该条来自宽召回批次；源码下载受 arXiv e-print 网络限制，部分论文暂时只能依据标题、摘要和类别分类，需要在入可信书前做人工/LLM精读复核。
% METHOD: 2402.10128
\paragraph{GES（压缩、量化与加速）（2024.02）}
GES~\cite{2402_10128} 关注的问题是：压缩或加速 3DGS 表示与渲染 GES 是宽召回补全中发现的论文，主要根据标题、摘要和 arXiv 类别判定其与 3DGS 的关系。
相对原始 3DGS，GES 的变化不是只替换损失函数，而是把表示、初始化、约束或网络预测路径重新组织到任务流程中：该工作被召回是因为论文文本包含 Gaussian/3D Gaussian 相关表述；当前分类先确定其任务域和与 3DGS 的相关性，后续证据流水线再补代码、企业和引用信息。
\textbf{实现细节：} 具体实现可以按输入、核心模块和渲染/优化路径来读。输入端首先取得任务所需的图像、相机、文本、时间序列、语义或几何线索，并把这些线索转换成可以约束高斯中心、尺度、旋转、不透明度、颜色或附加特征的条件。核心结构中，方法使用的主要机制是：从摘要层面看，GES 围绕“压缩或加速 3DGS 表示与渲染”展开。若其正文或代码后续可读，应进一步核对它具体改动的是表示、初始化、渲染、优化、损失、动态建模、压缩或应用接口中的哪一环。 训练或推理时，这些模块通常先给出初始化或中间特征，再通过可微 Gaussian rasterization、重建损失、感知/语义约束、几何正则或任务损失回传到高斯参数；如果方法包含前馈网络，则网络输出的是高斯属性、变形场、压缩码、语义特征或控制变量，而不是只输出最终 2D 图像。渲染路径仍以显式高斯投影、深度排序和 alpha 合成为基础，只是在采样、属性预测、局部编辑、动态变形、压缩编码或下游任务接口处加入额外约束。这样写入本章时，它与前后方法的关系是：相对原始 3DGS 和本章已有锚点方法，它可以理解为 3DGS extends\_or\_uses 宽召回分类阶段的关系判断，后续需要按论文方法部分复核。
\textbf{证据与边界：} 本轮召回分类尚未逐项核对实验表格；后续证据补全阶段需要读取论文源码/PDF或项目页确认数据集、指标和对比方法。 该条来自宽召回批次；源码下载受 arXiv e-print 网络限制，部分论文暂时只能依据标题、摘要和类别分类，需要在入可信书前做人工/LLM精读复核。 本书只采用本地 paper pack、数据库分类、源码/引用登记和论文摘要中可核验的信息；没有本地证据支撑的速度、排行榜或开源完整性不在正文中强化。\textbf{方法信息：} 代码：\url{https://github.com/rwightman/pytorch-image-models}；引用信号：121 citations (4.2/month, domain-relative cutoff 4.2/month, semantic\_scholar)；代码关注：ajhamdi/ges-splatting has 238 stars (8.3/month, domain-relative cutoff 13.9/month)。
\textbf{设计动机与边界：} 相对原始 3DGS 以独立高斯和图像重建误差为主的设定，GES 把问题落在以下可实现的改动上：从摘要层面看，GES 围绕“压缩或加速 3DGS 表示与渲染”展开。若其正文或代码后续可读，应进一步核对它具体改动的是表示、初始化、渲染、优化、损失、动态建模、压缩或应用接口中的哪一环。其设计目的不是单独增加一个模块，而是用上述表示、监督或渲染约束消除原流程在目标任务中的具体失配。该工作被召回是因为论文文本包含 Gaussian/3D Gaussian 相关表述；当前分类先确定其任务域和与 3DGS 的相关性，后续证据流水线再补代码、企业和引用信息。适用时还应保留论文给出的边界：该条来自宽召回批次；源码下载受 arXiv e-print 网络限制，部分论文暂时只能依据标题、摘要和类别分类，需要在入可信书前做人工/LLM精读复核。
% METHOD: 2403.08551
\paragraph{GaussianImage（压缩、量化与加速）（2024.03）}
GaussianImage~\cite{2403_08551} 关注的问题是：压缩或加速 3DGS 表示与渲染 GaussianImage 是宽召回补全中发现的论文，主要根据标题、摘要和 arXiv 类别判定其与 3DGS 的关系。
相对原始 3DGS，GaussianImage 的变化不是只替换损失函数，而是把表示、初始化、约束或网络预测路径重新组织到任务流程中：该工作被召回是因为论文文本包含 Gaussian/3D Gaussian 相关表述；当前分类先确定其任务域和与 3DGS 的相关性，后续证据流水线再补代码、企业和引用信息。
\textbf{实现细节：} 具体实现可以按输入、核心模块和渲染/优化路径来读。输入端首先取得任务所需的图像、相机、文本、时间序列、语义或几何线索，并把这些线索转换成可以约束高斯中心、尺度、旋转、不透明度、颜色或附加特征的条件。核心结构中，方法使用的主要机制是：从摘要层面看，GaussianImage 围绕“压缩或加速 3DGS 表示与渲染”展开。若其正文或代码后续可读，应进一步核对它具体改动的是表示、初始化、渲染、优化、损失、动态建模、压缩或应用接口中的哪一环。 训练或推理时，这些模块通常先给出初始化或中间特征，再通过可微 Gaussian rasterization、重建损失、感知/语义约束、几何正则或任务损失回传到高斯参数；如果方法包含前馈网络，则网络输出的是高斯属性、变形场、压缩码、语义特征或控制变量，而不是只输出最终 2D 图像。渲染路径仍以显式高斯投影、深度排序和 alpha 合成为基础，只是在采样、属性预测、局部编辑、动态变形、压缩编码或下游任务接口处加入额外约束。这样写入本章时，它与前后方法的关系是：相对原始 3DGS 和本章已有锚点方法，它可以理解为 3DGS extends\_or\_uses 宽召回分类阶段的关系判断，后续需要按论文方法部分复核。
\textbf{证据与边界：} 本轮召回分类尚未逐项核对实验表格；后续证据补全阶段需要读取论文源码/PDF或项目页确认数据集、指标和对比方法。 该条来自宽召回批次；源码下载受 arXiv e-print 网络限制，部分论文暂时只能依据标题、摘要和类别分类，需要在入可信书前做人工/LLM精读复核。 本书只采用本地 paper pack、数据库分类、源码/引用登记和论文摘要中可核验的信息；没有本地证据支撑的速度、排行榜或开源完整性不在正文中强化。\textbf{方法信息：} 代码：\url{https://github.com/xinjie-q/gaussianimage}；引用信号：106 citations (3.8/month, domain-relative cutoff 4.2/month, semantic\_scholar)；代码关注：xinjie-q/gaussianimage has 411 stars (14.7/month, domain-relative cutoff 13.9/month)。
\textbf{设计动机与边界：} 相对原始 3DGS 以独立高斯和图像重建误差为主的设定，GaussianImage 把问题落在以下可实现的改动上：从摘要层面看，GaussianImage 围绕“压缩或加速 3DGS 表示与渲染”展开。若其正文或代码后续可读，应进一步核对它具体改动的是表示、初始化、渲染、优化、损失、动态建模、压缩或应用接口中的哪一环。其设计目的不是单独增加一个模块，而是用上述表示、监督或渲染约束消除原流程在目标任务中的具体失配。该工作被召回是因为论文文本包含 Gaussian/3D Gaussian 相关表述；当前分类先确定其任务域和与 3DGS 的相关性，后续证据流水线再补代码、企业和引用信息。适用时还应保留论文给出的边界：该条来自宽召回批次；源码下载受 arXiv e-print 网络限制，部分论文暂时只能依据标题、摘要和类别分类，需要在入可信书前做人工/LLM精读复核。
% METHOD: 2403.14530
\paragraph{HAC（压缩、量化与加速）（2024.03）}
HAC~\cite{2403_14530} 关注的问题是：压缩或加速 3DGS 表示与渲染 HAC 是宽召回补全中发现的论文，主要根据标题、摘要和 arXiv 类别判定其与 3DGS 的关系。
相对原始 3DGS，HAC 的变化不是只替换损失函数，而是把表示、初始化、约束或网络预测路径重新组织到任务流程中：该工作被召回是因为论文文本包含 Gaussian/3D Gaussian 相关表述；当前分类先确定其任务域和与 3DGS 的相关性，后续证据流水线再补代码、企业和引用信息。
\textbf{实现细节：} 具体实现可以按输入、核心模块和渲染/优化路径来读。输入端首先取得任务所需的图像、相机、文本、时间序列、语义或几何线索，并把这些线索转换成可以约束高斯中心、尺度、旋转、不透明度、颜色或附加特征的条件。核心结构中，方法使用的主要机制是：从摘要层面看，HAC 围绕“压缩或加速 3DGS 表示与渲染”展开。若其正文或代码后续可读，应进一步核对它具体改动的是表示、初始化、渲染、优化、损失、动态建模、压缩或应用接口中的哪一环。 训练或推理时，这些模块通常先给出初始化或中间特征，再通过可微 Gaussian rasterization、重建损失、感知/语义约束、几何正则或任务损失回传到高斯参数；如果方法包含前馈网络，则网络输出的是高斯属性、变形场、压缩码、语义特征或控制变量，而不是只输出最终 2D 图像。渲染路径仍以显式高斯投影、深度排序和 alpha 合成为基础，只是在采样、属性预测、局部编辑、动态变形、压缩编码或下游任务接口处加入额外约束。这样写入本章时，它与前后方法的关系是：相对原始 3DGS 和本章已有锚点方法，它可以理解为 3DGS extends\_or\_uses 宽召回分类阶段的关系判断，后续需要按论文方法部分复核。
\textbf{证据与边界：} 本轮召回分类尚未逐项核对实验表格；后续证据补全阶段需要读取论文源码/PDF或项目页确认数据集、指标和对比方法。 该条来自宽召回批次；源码下载受 arXiv e-print 网络限制，部分论文暂时只能依据标题、摘要和类别分类，需要在入可信书前做人工/LLM精读复核。 本书只采用本地 paper pack、数据库分类、源码/引用登记和论文摘要中可核验的信息；没有本地证据支撑的速度、排行榜或开源完整性不在正文中强化。\textbf{方法信息：} 代码：\url{https://github.com/yihangchen-ee/hac}；引用信号：253 citations (9.2/month, domain-relative cutoff 4.2/month, semantic\_scholar)；代码关注：yihangchen-ee/hac has 331 stars (12.0/month, domain-relative cutoff 13.9/month)。
\textbf{原始方法补充：} HAC 的出发点是 Scaffold-GS 的锚点特征、尺度和偏移属性仍需以高码率逐项存储。它把锚点位置送入可学习的二值哈希网格并插值得到上下文特征，使用共享 MLP 同时预测自适应量化步长和各锚点属性的概率分布参数；概率再供算术编码降低熵。训练时以均匀噪声替代不可导的量化，测试时按预测步长舍入；偏移掩码会删除冗余高斯和锚点。为了在保持 Scaffold-GS 视图自适应表示的同时去除属性冗余，HAC 以锚点位置和属性为输入，输出可解码的压缩属性码流和重建高斯；训练在率失真目标下联合平衡渲染质量与编码长度，推理阶段依次查询哈希上下文、量化、熵编码/解码并恢复高斯。
% METHOD: 2505.23044
\paragraph{SpatialSplat（压缩、量化与加速）（2025.05）}
SpatialSplat~\cite{2505_23044} 关注的问题是：压缩或加速 3DGS 表示与渲染 SpatialSplat 是宽召回补全中发现的论文，主要根据标题、摘要和 arXiv 类别判定其与 3DGS 的关系。
相对原始 3DGS，SpatialSplat 的变化不是只替换损失函数，而是把表示、初始化、约束或网络预测路径重新组织到任务流程中：该工作被召回是因为论文文本包含 Gaussian/3D Gaussian 相关表述；当前分类先确定其任务域和与 3DGS 的相关性，后续证据流水线再补代码、企业和引用信息。
\textbf{实现细节：} 具体实现可以按输入、核心模块和渲染/优化路径来读。输入端首先取得任务所需的图像、相机、文本、时间序列、语义或几何线索，并把这些线索转换成可以约束高斯中心、尺度、旋转、不透明度、颜色或附加特征的条件。核心结构中，方法使用的主要机制是：从摘要层面看，SpatialSplat 围绕“压缩或加速 3DGS 表示与渲染”展开。若其正文或代码后续可读，应进一步核对它具体改动的是表示、初始化、渲染、优化、损失、动态建模、压缩或应用接口中的哪一环。 训练或推理时，这些模块通常先给出初始化或中间特征，再通过可微 Gaussian rasterization、重建损失、感知/语义约束、几何正则或任务损失回传到高斯参数；如果方法包含前馈网络，则网络输出的是高斯属性、变形场、压缩码、语义特征或控制变量，而不是只输出最终 2D 图像。渲染路径仍以显式高斯投影、深度排序和 alpha 合成为基础，只是在采样、属性预测、局部编辑、动态变形、压缩编码或下游任务接口处加入额外约束。这样写入本章时，它与前后方法的关系是：相对原始 3DGS 和本章已有锚点方法，它可以理解为 3DGS extends\_or\_uses 宽召回分类阶段的关系判断，后续需要按论文方法部分复核。
\textbf{证据与边界：} 本轮召回分类尚未逐项核对实验表格；后续证据补全阶段需要读取论文源码/PDF或项目页确认数据集、指标和对比方法。 该条来自宽召回批次；源码下载受 arXiv e-print 网络限制，部分论文暂时只能依据标题、摘要和类别分类，需要在入可信书前做人工/LLM精读复核。 本书只采用本地 paper pack、数据库分类、源码/引用登记和论文摘要中可核验的信息；没有本地证据支撑的速度、排行榜或开源完整性不在正文中强化。\textbf{方法信息：} 代码：\url{https://github.com/shengyuuu/spatialsplat}。
\textbf{设计动机与边界：} 相对原始 3DGS 以独立高斯和图像重建误差为主的设定，SpatialSplat 把问题落在以下可实现的改动上：从摘要层面看，SpatialSplat 围绕“压缩或加速 3DGS 表示与渲染”展开。若其正文或代码后续可读，应进一步核对它具体改动的是表示、初始化、渲染、优化、损失、动态建模、压缩或应用接口中的哪一环。其设计目的不是单独增加一个模块，而是用上述表示、监督或渲染约束消除原流程在目标任务中的具体失配。该工作被召回是因为论文文本包含 Gaussian/3D Gaussian 相关表述；当前分类先确定其任务域和与 3DGS 的相关性，后续证据流水线再补代码、企业和引用信息。适用时还应保留论文给出的边界：该条来自宽召回批次；源码下载受 arXiv e-print 网络限制，部分论文暂时只能依据标题、摘要和类别分类，需要在入可信书前做人工/LLM精读复核。
% METHOD: 2508.03017
\paragraph{SA-3DGS（压缩、量化与加速）（2025.08）}
SA-3DGS~\cite{2508_03017} 关注的问题是：压缩或加速 3DGS 表示与渲染 SA-3DGS 是宽召回补全中发现的论文，主要根据标题、摘要和 arXiv 类别判定其与 3DGS 的关系。
相对原始 3DGS，SA-3DGS 的变化不是只替换损失函数，而是把表示、初始化、约束或网络预测路径重新组织到任务流程中：该工作被召回是因为论文文本包含 Gaussian/3D Gaussian 相关表述；当前分类先确定其任务域和与 3DGS 的相关性，后续证据流水线再补代码、企业和引用信息。
\textbf{实现细节：} 具体实现可以按输入、核心模块和渲染/优化路径来读。输入端首先取得任务所需的图像、相机、文本、时间序列、语义或几何线索，并把这些线索转换成可以约束高斯中心、尺度、旋转、不透明度、颜色或附加特征的条件。核心结构中，方法使用的主要机制是：从摘要层面看，SA-3DGS 围绕“压缩或加速 3DGS 表示与渲染”展开。若其正文或代码后续可读，应进一步核对它具体改动的是表示、初始化、渲染、优化、损失、动态建模、压缩或应用接口中的哪一环。 训练或推理时，这些模块通常先给出初始化或中间特征，再通过可微 Gaussian rasterization、重建损失、感知/语义约束、几何正则或任务损失回传到高斯参数；如果方法包含前馈网络，则网络输出的是高斯属性、变形场、压缩码、语义特征或控制变量，而不是只输出最终 2D 图像。渲染路径仍以显式高斯投影、深度排序和 alpha 合成为基础，只是在采样、属性预测、局部编辑、动态变形、压缩编码或下游任务接口处加入额外约束。这样写入本章时，它与前后方法的关系是：相对原始 3DGS 和本章已有锚点方法，它可以理解为 3DGS extends\_or\_uses 宽召回分类阶段的关系判断，后续需要按论文方法部分复核。
\textbf{证据与边界：} 本轮召回分类尚未逐项核对实验表格；后续证据补全阶段需要读取论文源码/PDF或项目页确认数据集、指标和对比方法。 该条来自宽召回批次；源码下载受 arXiv e-print 网络限制，部分论文暂时只能依据标题、摘要和类别分类，需要在入可信书前做人工/LLM精读复核。 本书只采用本地 paper pack、数据库分类、源码/引用登记和论文摘要中可核验的信息；没有本地证据支撑的速度、排行榜或开源完整性不在正文中强化。\textbf{方法信息：} 代码：\url{https://github.com/sssdavid529/satellite-3dgs-eval-protocol}；本条的源码/章节抽取不完整，正文只使用摘要、分类和已登记证据。
\textbf{设计动机与边界：} 相对原始 3DGS 以独立高斯和图像重建误差为主的设定，SA-3DGS 把问题落在以下可实现的改动上：从摘要层面看，SA-3DGS 围绕“压缩或加速 3DGS 表示与渲染”展开。若其正文或代码后续可读，应进一步核对它具体改动的是表示、初始化、渲染、优化、损失、动态建模、压缩或应用接口中的哪一环。其设计目的不是单独增加一个模块，而是用上述表示、监督或渲染约束消除原流程在目标任务中的具体失配。该工作被召回是因为论文文本包含 Gaussian/3D Gaussian 相关表述；当前分类先确定其任务域和与 3DGS 的相关性，后续证据流水线再补代码、企业和引用信息。适用时还应保留论文给出的边界：该条来自宽召回批次；源码下载受 arXiv e-print 网络限制，部分论文暂时只能依据标题、摘要和类别分类，需要在入可信书前做人工/LLM精读复核。
% METHOD: 2606.30017
\paragraph{Monte Carlo Energy Aggregation（移动端高阶辐射能量聚合）（2026.06）}
Monte Carlo Energy Aggregation for Mobile 3D Gaussian Splatting~\cite{2606_30017} 也称 Flux-GS，面向移动端 3DGS。原始 3DGS 常用三阶球谐表达视角相关外观，每个高斯需要 48 个 SH 浮点参数；在百万级高斯场景中，这会带来存储、带宽和移动 GPU 读取压力。Flux-GS 的改变，是把高阶 SH 的 specular radiance residual 采样并聚合成更紧凑的能量表示，再回到一阶 SH 上渲染。

训练流程先用三阶 SH 优化若干初始迭代，使模型捕获完整视角相关外观；随后 Monte Carlo Specular Energy Aggregator 在不同观察方向上采样三阶 radiance residual，把高光能量统计进 compact latent。为了弥补降到一阶 SH 后损失的高频细节，Attribute-Conditioned SH Enhancement 使用高斯自身属性预测 Gaussian-aware offsets，并在推理前把这些 offset 烘焙回一阶 SH 参数；也就是说，推理时不再额外跑一个 MLP，而是直接读取压缩后的显式属性。这样它不是简单删除高阶 SH，而是先学习高阶响应，再把可见的 specular 能量折叠到低阶表示。

Flux-GS 还处理 primitive 数量问题。原始单视角梯度 densification 容易对某个训练视角过拟合，生成对多视角结构帮助不大的高斯；它提出 Multi-view Alpha-based Densification and Pruning，用分层采样相机和跨视角 alpha 加权误差累计来决定哪里新增、哪里删除。相对 QuantumGS 的函数族替换，Flux-GS 更像“辐射参数压缩 + 多视角密度控制”的移动端组合；相对一般文件 codec，它明确服务 WebGL/移动渲染，公开 README 也提供压缩文件解码、WebGL 移动 renderer 和训练脚本。
\textbf{锚点细化：} Monte Carlo Energy Aggregation for Mobile 3D Gaussian Splatting：在 3DGS 的“Mobile Gaussian rendering and radiance compression”基础上推进，主要改变是：Monte Carlo Specular Energy Aggregator 与紧凑 latent 代替高阶 SH 开销。它要解决的问题是：针对 high-order Spherical Harmonics 带来的显著 overhead，提出面向移动端的实时高保真 GS 渲染方法。核心思想是：Flux-GS 通过聚合高阶辐射残差中的镜面能量，降低移动平台上 3DGS 的存储和推理开销。
\textbf{方法信息：} Baidu 参与了这篇论文；代码：\url{https://github.com/xiaobiaodu/Flux-GS}；GitHub 约 33 stars。

\begin{takeaway}
本章可以用“先少长、再少存、最后少画”来概括。Taming 3DGS、DashGaussian、GSta 和 SteepGS 试图在训练阶段减少无效增长和无效更新；LightGaussian、PUP、MaskGaussian、Gaussian Herding、REFINE 和 VEDAL 判断哪些高斯可以删、可以合并或可以异步处理；RDO-Gaussian、CompGS、ContextGS、HAC++、MesonGS、HybridGS、GSCodec Studio 等把属性和结构写进码率模型；Scaffold-GS、SAGS、ProtoGS、V3DG、CLoD-GS 则给高斯建立锚点、原型或 LoD 层次。对读者来说，压缩方法不应只看最终 MB 数，而要看它改的是哪一层：densification、primitive 数量、属性表示、概率模型、bitstream 渐进性、渲染时可见集合，还是移动端硬件执行路径。
\end{takeaway}
